% Les bruits et les nuisances sonores — Allemand 1ère A — Cours signé OnBuch+
% Programme officiel MINESEC. Compiler avec : tectonic ALA13-les-bruits-et-les-nuisances-sonores.tex
\newcommand{\DOCMATIERE}{Allemand}
\newcommand{\DOCNIVEAU}{1ère A}
\newcommand{\DOCMODULE}{Chapitre 3 : Environnement, santé et bien-être}
\newcommand{\DOCLECON}{ALA13}
\newcommand{\DOCTITRE}{Les bruits et les nuisances sonores}
% OnBuch+ — préambule « cours signés » ENRICHI (compilé avec tectonic / XeLaTeX)
% Système visuel « Pop » d'OnBuch+ : fond crème (jamais blanc), bordures encre
% épaisses, ombres « dures » décalées, titres Archivo Black en capitales.
% Version MATHÉMATIQUES : graphes (pgfplots/tikz), boîtes pédagogiques
% (définition, propriété, méthode, attention, le savais-tu, exercice résolu,
% exercice, TP, bilan), page de garde et signature OnBuch+ ancrée en bas.
%
% Macros attendues dans main.tex AVANT \input{preamble} :
%   \DOCMATIERE  \DOCNIVEAU  \DOCMODULE  \DOCLECON (numéro)  \DOCTITRE
\documentclass[11pt]{article}

\usepackage{fontspec}
\usepackage[ngerman,french]{babel} % français = langue principale ; l'allemand via \all{..} / textebox
\usepackage[a4paper,margin=2cm,top=3.4cm,bottom=2.8cm,headheight=1.4cm,footskip=1.3cm]{geometry}
\usepackage{amsmath,amssymb,amsfonts}
\usepackage{mathtools} % \xmapsto, \coloneqq... souvent utilisés par les modèles
\usepackage{cancel} % simplifications barrées (bilans de forces)
% \abs{z} (module, valeur absolue) et \norm{u} (norme), utilisés par les modèles
\providecommand{\abs}{}\let\abs\relax\DeclarePairedDelimiter{\abs}{\lvert}{\rvert}
\providecommand{\norm}{}\let\norm\relax\DeclarePairedDelimiter{\norm}{\lVert}{\rVert}
\usepackage[table]{xcolor} % [table] : fournit aussi \rowcolors (alternance de lignes)
\usepackage{pagecolor}
\usepackage{tikz}
\usetikzlibrary{arrows.meta,positioning,calc,shapes.geometric,shapes.misc,shapes.arrows,decorations.pathmorphing,decorations.pathreplacing,decorations.markings,patterns,shadows,mindmap,trees,fit,backgrounds,matrix,intersections,angles,quotes}
\usepackage{pgfplots}
\usepackage[version=4]{mhchem} % équations nucléaires \ce{^{235}_{92}U}
\usepackage[european,siunitx]{circuitikz} % schémas électriques
\pgfplotsset{compat=1.18}
\usepgfplotslibrary{fillbetween,groupplots} % aires entre courbes (calcul intégral)
\usepackage{enumitem}
\usepackage{fancyhdr}
% [most] charge toutes les bibliothèques tcolorbox (listings, minted, raster,
% documentation...) et gonfle inutilement la mémoire TeX sur un document long
% (~300 boîtes) : on ne charge que ce qui est réellement utilisé.
\usepackage[skins,breakable]{tcolorbox}
\usepackage{graphicx}
\usepackage{colortbl}
\usepackage{booktabs}
\usepackage{tabularx}
\usepackage{diagbox} % case d'en-tête diagonale des tableaux à double entrée
% Colonnes extensibles centrée / gauche / droite (C, L, R), souvent utilisées par les modèles
\newcolumntype{C}{>{\centering\arraybackslash}X}
\newcolumntype{L}{>{\raggedright\arraybackslash}X}
\newcolumntype{R}{>{\raggedleft\arraybackslash}X}
\usepackage{array}
\usepackage{multicol}
\usepackage{multirow}
\usepackage{siunitx}
% parse-numbers=false : accepte aussi \SI{2,5 \times 10^{-3}}{...}
\sisetup{locale=FR,output-decimal-marker={,},group-minimum-digits=4,parse-numbers=false}
\DeclareSIUnit\ohm{\ensuremath{\symup{\Omega}}} % U+2126 absent de la police
\DeclareSIUnit\division{div} % réglages d'oscilloscope (ms/div, V/div)
\DeclareSIUnit\tr{tr} % tours (vitesse de rotation, ex. \SI{1500}{\tr\per\minute})
\DeclareSIUnit\molar{M} % concentration molaire (ex. \SI{3}{\molar}, \SI{1}{\micro\molar})
\DeclareSIUnit\an{an} % année (le modèle confond parfois avec \year, un compteur TeX) — ex. \SI{5}{\kilo\meter\per\an}
\DeclareSIUnit\FCFA{FCFA} % franc CFA (ex. \SI{500}{\FCFA\per\kilo\gram})
\DeclareSIUnit\degreeBrix{\textdegree Brix} % teneur en sucre d'un sirop/jus (réfractométrie)
\DeclareSIUnit\month{mois} % le modèle utilise parfois \month, un compteur TeX, comme unité de durée
\usepackage{qrcode}
% \degree : commande fréquemment utilisée par les modèles pour un angle ou une
% température en dehors de \SI{}{} (non fournie par les paquets chargés).
\providecommand{\degree}{^{\circ}}
% \qed (fin de démonstration), utilisé par les modèles sans amsthm
\providecommand{\qed}{\hfill$\square$}
% \barre : double barre des valeurs interdites dans un tableau de variations (tkz-tab)
\providecommand{\barre}{\ensuremath{\|}}
% environnement proof (amsthm non chargé) utilisé par les modèles
\ifdefined\proof\else\newenvironment{proof}[1][Démonstration]{\par\noindent\textit{#1.}\ }{\qed\par}\fi
\usepackage{lastpage}
\usepackage{hyperref}
\hypersetup{hidelinks}

% ============================================================
% Palette « Pop » OnBuch+ — mêmes hex que la classe P (plus_ui.dart)
% ============================================================
\definecolor{poporange}{HTML}{F59321}
\definecolor{popdark}{HTML}{E07A0C}
\definecolor{popcream}{HTML}{FFF6E8}
\definecolor{popink}{HTML}{1C1714}
\definecolor{poppurple}{HTML}{7A5AE0}
\definecolor{popgreen}{HTML}{2E9E6A}
\definecolor{poppink}{HTML}{E0457B}
\definecolor{popblue}{HTML}{2F7DE1}
\definecolor{popgold}{HTML}{C98A12}
\definecolor{popmuted}{HTML}{8A7F76}
\definecolor{popline}{HTML}{EADFD2}
\definecolor{popgray}{HTML}{8A7F76}
\colorlet{popgrey}{popgray}
% Alias tolérants (le modèle nomme parfois les couleurs différemment)
\colorlet{popwhite}{white}
\colorlet{popblack}{popink}
\colorlet{popred}{poppink}
\colorlet{popyellow}{popgold}
% Teintes claires (fonds de tableaux/encadrés) — le modèle nomme parfois
% les couleurs avec un suffixe L (ex. popblueL) sans les avoir définies.
\colorlet{poporangeL}{poporange!20}
\colorlet{popdarkL}{popdark!20}
\colorlet{poppurpleL}{poppurple!20}
\colorlet{popgreenL}{popgreen!20}
\colorlet{poppinkL}{poppink!20}
\colorlet{popblueL}{popblue!20}
\colorlet{popgoldL}{popgold!20}
\colorlet{popredL}{popred!20}
\colorlet{popgrayL}{popgray!20}
% Teintes claires pour fonds de tableaux / zones
\colorlet{popblueL}{popblue!10!white}
\colorlet{poporangeL}{poporange!14!white}
\colorlet{popgreenL}{popgreen!12!white}
\colorlet{poppurpleL}{poppurple!12!white}
\colorlet{poppinkL}{poppink!12!white}
\colorlet{popgoldL}{popgold!14!white}

\pagecolor{popcream}

% ============================================================
% Typographie
% ============================================================
\defaultfontfeatures{Ligatures=TeX}
\setmainfont{PlusJakartaSans-Medium.ttf}[Path=../../../fonts/,BoldFont=PlusJakartaSans-Bold.ttf]
% Maths en sans-serif (Fira Math) avec chiffres et lettres droites de Plus
% Jakarta, pour que les formules se fondent dans le texte.
\usepackage[math-style=ISO,bold-style=ISO]{unicode-math}
\setmathfont{FiraMath-Regular.otf}[Path=../../../fonts/]
\setmathfont{PlusJakartaSans-Medium.ttf}[Path=../../../fonts/,range={up/num,up/{Latin,latin},\mathup}]
\usepackage{icomma} % virgule décimale sans espace en mode math
% Caractères mathématiques Unicode absents des polices de texte (Plus Jakarta,
% Archivo Black) : rendus par leur équivalent mathématique, en texte comme en
% formule — sinon ils s'affichent en carré vide (ex. « Arithmétique dans ℤ »).
\usepackage{newunicodechar}
\newunicodechar{ℝ}{\ensuremath{\mathbb{R}}}
\newunicodechar{ℤ}{\ensuremath{\mathbb{Z}}}
\newunicodechar{ℂ}{\ensuremath{\mathbb{C}}}
\newunicodechar{ℕ}{\ensuremath{\mathbb{N}}}
\newunicodechar{ℚ}{\ensuremath{\mathbb{Q}}}
\newunicodechar{𝔻}{\ensuremath{\mathbb{D}}}
\newunicodechar{α}{\ensuremath{\alpha}}
\newunicodechar{β}{\ensuremath{\beta}}
\newunicodechar{γ}{\ensuremath{\gamma}}
\newunicodechar{δ}{\ensuremath{\delta}}
\newunicodechar{ε}{\ensuremath{\varepsilon}}
\newunicodechar{θ}{\ensuremath{\theta}}
\newunicodechar{λ}{\ensuremath{\lambda}}
\newunicodechar{μ}{\ensuremath{\mu}}
\newunicodechar{σ}{\ensuremath{\sigma}}
\newunicodechar{τ}{\ensuremath{\tau}}
\newunicodechar{φ}{\ensuremath{\varphi}}
\newunicodechar{ψ}{\ensuremath{\psi}}
\newunicodechar{ω}{\ensuremath{\omega}}
\newunicodechar{Σ}{\ensuremath{\Sigma}}
% majuscules produites par \MakeUppercase dans les titres (α-aminés -> Α)
\newunicodechar{Α}{\ensuremath{\alpha}}
\newunicodechar{Β}{\ensuremath{\beta}}
\newunicodechar{Γ}{\ensuremath{\Gamma}}
\newunicodechar{Δ}{\ensuremath{\Delta}}
\newunicodechar{Ε}{\ensuremath{\varepsilon}}
\newunicodechar{Θ}{\ensuremath{\Theta}}
\newunicodechar{Λ}{\ensuremath{\Lambda}}
\newunicodechar{Μ}{\ensuremath{\mu}}
\newunicodechar{Τ}{\ensuremath{\tau}}
\newunicodechar{Φ}{\ensuremath{\Phi}}
\newunicodechar{Ψ}{\ensuremath{\Psi}}
\newunicodechar{Ω}{\ensuremath{\Omega}}
\newunicodechar{↦}{\ensuremath{\mapsto}}
\newunicodechar{∈}{\ensuremath{\in}}
\newunicodechar{∉}{\ensuremath{\notin}}
\newunicodechar{∧}{\ensuremath{\wedge}}
\newunicodechar{⇔}{\ensuremath{\Leftrightarrow}}
\newunicodechar{⟺}{\ensuremath{\Longleftrightarrow}}
\newunicodechar{⇒}{\ensuremath{\Rightarrow}}
\newunicodechar{⟹}{\ensuremath{\Longrightarrow}}
\newunicodechar{∘}{\ensuremath{\circ}}
\newunicodechar{∪}{\ensuremath{\cup}}
\newunicodechar{∩}{\ensuremath{\cap}}
\newunicodechar{≡}{\ensuremath{\equiv}}
\newunicodechar{⊂}{\ensuremath{\subset}}
\newunicodechar{⊥}{\ensuremath{\perp}}
\newunicodechar{∃}{\ensuremath{\exists}}
\newunicodechar{∀}{\ensuremath{\forall}}
\newunicodechar{∛}{\ensuremath{\sqrt[3]{\,}}}
\newunicodechar{∜}{\ensuremath{\sqrt[4]{\,}}}
\newunicodechar{⃗}{\ensuremath{{}^{\to}}}
\newunicodechar{ⁿ}{\ifmmode^{n}\else\textsuperscript{\ensuremath{n}}\fi}
\newunicodechar{ˣ}{\ifmmode^{x}\else\textsuperscript{\ensuremath{x}}\fi}
\newunicodechar{ᵇ}{\ifmmode^{b}\else\textsuperscript{\ensuremath{b}}\fi}
\newunicodechar{ʳ}{\ifmmode^{r}\else\textsuperscript{\ensuremath{r}}\fi}
\newunicodechar{ᵘ}{\ifmmode^{u}\else\textsuperscript{\ensuremath{u}}\fi}
\newunicodechar{ʸ}{\ifmmode^{y}\else\textsuperscript{\ensuremath{y}}\fi}
\newunicodechar{ᵃ}{\ifmmode^{a}\else\textsuperscript{\ensuremath{a}}\fi}
\newunicodechar{ᵉ}{\ifmmode^{e}\else\textsuperscript{\ensuremath{e}}\fi}
\newunicodechar{ᵏ}{\ifmmode^{k}\else\textsuperscript{\ensuremath{k}}\fi}
\newunicodechar{ᵖ}{\ifmmode^{p}\else\textsuperscript{\ensuremath{p}}\fi}
\newunicodechar{ᴺ}{\ifmmode^{N}\else\textsuperscript{\ensuremath{N}}\fi}
\newunicodechar{ᶜ}{\ifmmode^{c}\else\textsuperscript{\ensuremath{c}}\fi}
\newunicodechar{ʰ}{\ifmmode^{h}\else\textsuperscript{\ensuremath{h}}\fi}
\newunicodechar{ᵠ}{\ifmmode^{q}\else\textsuperscript{\ensuremath{q}}\fi}
\newunicodechar{ₙ}{\ifmmode_{n}\else\textsubscript{\ensuremath{n}}\fi}
\newunicodechar{ₐ}{\ifmmode_{a}\else\textsubscript{\ensuremath{a}}\fi}
\newunicodechar{ᵢ}{\ifmmode_{i}\else\textsubscript{\ensuremath{i}}\fi}
\newunicodechar{ᵣ}{\ifmmode_{r}\else\textsubscript{\ensuremath{r}}\fi}
\newunicodechar{ₖ}{\ifmmode_{k}\else\textsubscript{\ensuremath{k}}\fi}
\newunicodechar{ᵦ}{\ifmmode_{\beta}\else\textsubscript{\ensuremath{\beta}}\fi}
\newunicodechar{ₓ}{\ifmmode_{x}\else\textsubscript{\ensuremath{x}}\fi}
\newfontfamily\popdisplay{ArchivoBlack-Regular.ttf}[Path=../../../fonts/]
\newfontfamily\poplogo{SpaceGrotesk-Bold.ttf}[Path=../../../fonts/]
\newfontfamily\popsemi{PlusJakartaSans-SemiBold.ttf}[Path=../../../fonts/]
\newfontfamily\popxbold{PlusJakartaSans-ExtraBold.ttf}[Path=../../../fonts/]
\newfontfamily\popxboldls{PlusJakartaSans-ExtraBold.ttf}[Path=../../../fonts/,LetterSpace=3.0]

\setlist[enumerate]{leftmargin=1.7em,itemsep=5pt,topsep=4pt}
\setlist[itemize]{leftmargin=1.7em,itemsep=4pt,topsep=4pt,label={\color{poporange}\textbullet}}
\setlength{\parindent}{0pt}
\setlength{\parskip}{5pt}
\renewcommand{\arraystretch}{1.25}

% pgfplots : style Pop par défaut
\pgfplotsset{
  every axis/.append style={axis line style={popink,line width=1pt},tick style={popink},
    grid style={popline,line width=0.5pt},label style={font=\small},tick label style={font=\footnotesize}},
  popcurve/.style={poporange,line width=1.8pt,no markers,smooth},
}
% Même style disponible dans un \draw TikZ simple (hors pgfplots)
\tikzset{popcurve/.style={poporange,line width=1.8pt}}
\tikzset{popgrid/.style={popline,very thin}}
\colorlet{popcurve}{poporange} % le nom du style est parfois utilisé comme couleur

% ============================================================
% Pill / squiggle
% ============================================================
\newcommand{\poppill}[3]{%
  \tikz[baseline=-0.55ex]\node[draw=popink,line width=1.2pt,rounded corners=2.6pt,%
    fill=#2,inner xsep=6.5pt,inner ysep=3pt]{%
    \popxboldls\fontsize{7}{7}\selectfont\color{#3}\MakeUppercase{#1}};%
}
\newcommand{\popsquiggle}[1]{%
  \tikz[baseline=-0.4ex]{
    \draw[#1,line width=2.6pt,line cap=round]
      plot [smooth] coordinates {(0,0.09) (0.34,0.01) (0.68,0.21) (1.02,0.01) (1.36,0.10)};
  }%
}
% Mot-clé surligné (marqueur orange sous le texte)
\newcommand{\cle}[1]{{\popxbold\color{popdark}#1}}

% ============================================================
% En-tête / pied
% ============================================================
\pagestyle{fancy}
\fancyhf{}
\renewcommand{\headrulewidth}{0pt}
\renewcommand{\footrulewidth}{0pt}
\lhead{\raisebox{-0.15ex}{\poplogo\fontsize{13}{13}\selectfont\color{popink}OnBuch\color{poporange}+}}
\rhead{\poppill{\DOCMATIERE\ \textbullet\ \DOCNIVEAU\ \textbullet\ Leçon \DOCLECON}{white}{popink}}
\lfoot{\popxbold\fontsize{7.6}{7.6}\selectfont\color{popmuted}\MakeUppercase{Cours signé OnBuch+}\ \textbullet\ {\popsemi\DOCTITRE}}
\rfoot{\tikz[baseline=-0.6ex]\node[rounded corners=8.5pt,fill=popink,minimum height=17pt,minimum width=17pt,inner xsep=5pt,inner sep=0]{\popxbold\fontsize{8}{8}\selectfont\color{popcream}\thepage\,/\,\pageref*{LastPage}};}

% ============================================================
% Titres
% ============================================================
\newcommand{\chaptitre}[2]{%
  \begin{tcolorbox}[enhanced,colback=poporange,colframe=popink,boxrule=2.6pt,arc=11pt,
      drop shadow=popink,left=15pt,right=15pt,top=12pt,bottom=11pt,before skip=3pt,after skip=18pt]
    \poppill{#2}{popink}{popcream}\par\vspace{9pt}
    {\raggedright\hyphenpenalty=10000\exhyphenpenalty=10000\popdisplay\fontsize{22}{24}\selectfont\color{popcream}\MakeUppercase{#1}\par}
  \end{tcolorbox}
}
% Section numérotée : pastille orange avec le numéro + titre Archivo
\newcounter{popsec}
\newcommand{\coursec}[1]{%
  \stepcounter{popsec}\setcounter{popsub}{0}%
  \par\vspace{16pt}\noindent
  \begin{minipage}{\linewidth}
  \tikz[baseline=-0.7ex]\node[draw=popink,line width=1.6pt,fill=poporange,rounded corners=4pt,minimum size=22pt,inner sep=0]{\popdisplay\fontsize{12}{12}\selectfont\color{popcream}\thepopsec};%
  \hspace{8pt}{\popdisplay\fontsize{15}{17}\selectfont\color{popink}\MakeUppercase{#1}}\par
  \vspace{4pt}\noindent{\color{poporange}\rule{36pt}{3.6pt}}
  \end{minipage}\par\nopagebreak\vspace{8pt}
}
\newcounter{popsub}
\newcommand{\courssub}[1]{%
  \stepcounter{popsub}%
  \par\vspace{9pt}\noindent{\popxbold\fontsize{11.5}{13}\selectfont\color{popink}{\color{poporange}\thepopsec.\thepopsub}\hspace{6pt}#1}\par\nopagebreak\vspace{4pt}
}

% ============================================================
% Boîtes pédagogiques — même recette : fond blanc, bordure épaisse colorée,
% ombre dure encre, badge plein. Titre optionnel [#1].
% ============================================================
\tcbset{popbase/.style={breakable,enhanced,colback=white,boxrule=2.3pt,arc=9pt,
  shadow={3pt}{-3pt}{0pt}{popink},left=14pt,right=14pt,top=11pt,bottom=11pt,before skip=11pt,after skip=15pt,
  fontupper=\popsemi\fontsize{10.6}{14.5}\selectfont\color{popink},
  fontlower=\popsemi\fontsize{10.6}{14.5}\selectfont\color{popink}}}
% Coche dessinée (le glyphe ✓ n'existe pas dans les polices du document)
\newcommand{\popcheck}[1]{\tikz[baseline=-0.5ex]\draw[#1,line width=1.6pt,line cap=round,line join=round](-0.13,0.0)--(-0.04,-0.09)--(0.13,0.1);}
\providecommand{\coche}{\popcheck{popgreen}}
\providecommand{\resultat}[1]{\par\smallskip\noindent\fcolorbox{popgreen}{popgreenL}{\parbox{\dimexpr\linewidth-2\fboxsep-2\fboxrule\relax}{\textbf{Résultat.} #1}}\par\smallskip}
\newcommand{\popboxhead}[3]{\poppill{#1}{#2}{white}\ifx\relax#3\relax\else\hspace{6pt}{\popxbold\fontsize{10.6}{12}\selectfont\color{popink}#3}\fi\par\vspace{7pt}}

\newtcolorbox{aretenir}[1][]{popbase,colframe=popgreen,before upper={\popboxhead{À retenir}{popgreen}{#1}}}
% Texte en allemand (document de lecture, dialogue, extrait) : cadre
% d'étude. Un titre optionnel (source, auteur) s'affiche dans l'en-tête.
\newtcolorbox{textebox}[1][]{popbase,colframe=popblue,before upper={\popboxhead{Text}{popblue}{#1}\begin{otherlanguage}{ngerman}},after upper={\end{otherlanguage}}}
% Marques ✓ / ✗ (forme correcte / fautive) souvent utilisées par les modèles
\providecommand{\cross}{\textcolor{poppink}{\ensuremath{\times}}}
\providecommand{\xmark}{\cross}\providecommand{\crossmark}{\cross}\providecommand{\cmark}{\ensuremath{\checkmark}}
% Ligne de réponse / justification à compléter (inventée par les modèles)
\providecommand{\justification}[1]{\par\noindent\textit{Justification~:} \dotfill\par}
% \textipa (package tipa non chargé) : la police ne couvre pas l'API -> simple passage
\providecommand{\textipa}[1]{#1}
% Soulignement (ulem non chargé) : \uline, \ul
\providecommand{\uline}[1]{\underline{#1}}\providecommand{\ul}[1]{\underline{#1}}
% \alert (beamer) inventée par les modèles : texte mis en évidence
\providecommand{\alert}[1]{\textbf{\textcolor{poppink}{#1}}}
% variantes ASCII d'ordinaux écrites en macro
\providecommand{\iere}{\textsuperscript{re}}\providecommand{\ieme}{\textsuperscript{e}}\providecommand{\ere}{\textsuperscript{re}}\providecommand{\eme}{\textsuperscript{e}}\providecommand{\er}{\textsuperscript{er}}\providecommand{\ie}{\textsuperscript{e}}
% Ordinaux français écrits en macro par les modèles : 1\ière, 3\ième
\newcommand{\ière}{\textsuperscript{re}}\newcommand{\ième}{\textsuperscript{e}}
% \exemple (inventée par les modèles) : étiquette « Exemple : »
\providecommand{\exemple}{\textbf{Exemple~:}\ }
% \square utilisable aussi hors mode math (cases à cocher)
\AtBeginDocument{\let\squareMath\square\renewcommand{\square}{\ensuremath{\squareMath}}}
% Mot ou phrase en allemand à l'intérieur d'un paragraphe français
\newcommand{\all}[1]{\ifmmode\text{\textit{#1}}\else\begingroup\selectlanguage{ngerman}\itshape #1\endgroup\fi}
\let\esp\all % alias toléré
\newtcolorbox{exemplebox}[1][]{popbase,colframe=poporange,before upper={\popboxhead{Exemple}{poporange}{#1}}}
\newtcolorbox{definition}[1][]{popbase,colframe=popblue,before upper={\popboxhead{Définition}{popblue}{#1}}}
\newtcolorbox{propriete}[1][]{popbase,colframe=poppurple,before upper={\popboxhead{Propriété}{poppurple}{#1}}}
\newtcolorbox{methode}[1][]{popbase,colframe=popgold,before upper={\popboxhead{Méthode}{popgold}{#1}}}
\newtcolorbox{attention}[1][]{popbase,colframe=poppink,before upper={\popboxhead{Attention}{poppink}{#1}}}
\newtcolorbox{experience}[1][]{popbase,colframe=popblue,colback=popblueL,before upper={\popboxhead{Expérience}{popblue}{#1}}}
% « Le savais-tu ? » : carte violette pleine (contexte, culture, Cameroun)
\newtcolorbox{savaistu}[1][]{popbase,colback=poppurple,colframe=popink,
  fontupper=\popsemi\fontsize{10.4}{14.2}\selectfont\color{white},
  before upper={\poppill{Le savais-tu\,?}{white}{poppurple}\ifx\relax#1\relax\else\hspace{6pt}{\popxbold\color{white}#1}\fi\par\vspace{7pt}}}
% Exercice résolu : énoncé en haut, \tcblower puis la solution sur fond crème
\newcounter{popexo}\newcounter{popexr}
\newtcolorbox{exoresolu}[1][]{popbase,colframe=poporange,colbacklower=popcream,
  segmentation style={popink,line width=1.2pt,dashed},
  before upper={\stepcounter{popexr}\popboxhead{Exercice résolu \thepopexr}{poporange}{#1}},
  before lower={\poppill{Solution}{popink}{popcream}\par\vspace{6pt}}}
% Exercice (non résolu en place) — la correction est dans la partie Corrigés
\newtcolorbox{exercice}[1][]{popbase,colframe=popink,
  before upper={\stepcounter{popexo}\popboxhead{Exercice \thepopexo}{popink}{#1}}}
% Corrigé
\newtcolorbox{corrige}[1][]{popbase,colframe=popgreen,colback=popgreenL,
  before upper={\popboxhead{Corrigé}{popgreen}{#1}}}
% Objectifs / prérequis / bilan
\newtcolorbox{objectifs}{popbase,colframe=popink,colback=poporangeL,
  before upper={\popboxhead{Objectifs de la leçon}{popink}{}}}
\newtcolorbox{prerequis}{popbase,colframe=popink,colback=popblueL,
  before upper={\popboxhead{Prérequis}{popblue}{}}}
\newtcolorbox{bilan}{popbase,colframe=popink,colback=white,boxrule=2.8pt,
  before upper={\popboxhead{Fiche bilan}{poporange}{L'essentiel à connaître pour l'examen}}}
% Figure encadrée avec légende
\newtcolorbox{popfigure}[1][]{enhanced,breakable=false,colback=white,colframe=popline,boxrule=1.4pt,arc=8pt,
  left=10pt,right=10pt,top=10pt,bottom=8pt,before skip=10pt,after skip=12pt,halign=center,
  lower separated=false,halign lower=center,
  fontlower=\popsemi\fontsize{9}{11.5}\selectfont\color{popmuted},
  #1}
\newcommand{\legende}[1]{\par\vspace{6pt}{\popsemi\fontsize{9}{11.5}\selectfont\color{popmuted}#1}}

% ============================================================
% Signature OnBuch+
% ============================================================
\newcommand{\poplogoinline}[1]{{\poplogo\fontsize{#1}{#1}\selectfont\color{popink}OnBuch\color{poporange}+}}
\newcommand{\signatureonbuch}{%
  \begin{tcolorbox}[enhanced,colback=popink,colframe=popink,boxrule=2.2pt,arc=10pt,
      drop shadow=poporange,left=14pt,right=14pt,top=10pt,bottom=10pt,before skip=0pt,after skip=0pt]
    \tikz[baseline=-0.7ex]\node[circle,fill=popgreen,draw=popcream,line width=1.3pt,minimum size=22pt,inner sep=0]{\popcheck{white}};%
    \hspace{9pt}\begin{minipage}[c]{0.62\linewidth}
      {\popxbold\fontsize{12}{13}\selectfont\color{popcream}Signé\ \textbullet\ {\poplogo OnBuch\color{poporange}+}}\par\vspace{2pt}
      {\raggedright\popsemi\fontsize{8}{10.5}\selectfont\color{popline}\MakeUppercase{Équipe pédagogique OnBuch+}\\\MakeUppercase{Conforme au programme officiel MINESEC}\par}
    \end{minipage}\hfill
    {\poplogo\fontsize{22}{22}\selectfont\color{popcream}OnBuch\color{poporange}+}
  \end{tcolorbox}%
}

% ============================================================
% Page de garde
% #1 titre · #2 matière · #3 niveau · #4 module · #5 n° leçon · #6 durée ·
% #7 description / objectifs (texte ou itemize)
% ============================================================
\newcommand{\pagegarde}[7]{%
  \thispagestyle{empty}
  \newgeometry{a4paper,margin=1.7cm,top=1.6cm,bottom=1.4cm}%
  \begin{tikzpicture}[remember picture,overlay]
    \fill[poporange!18] ([xshift=-3.2cm,yshift=-3.6cm]current page.north east) circle (4.6cm);
    \fill[poppurple!14] ([xshift=2.4cm,yshift=7.6cm]current page.south west) circle (3.8cm);
  \end{tikzpicture}%
  {\poplogo\fontsize{30}{30}\selectfont\color{popink}OnBuch\color{poporange}+}\hfill
  \raisebox{0.6ex}{\poppill{Cours signé \textbullet\ Premium}{popink}{popcream}}\par
  \vspace{1pt}\popsquiggle{poppurple}\par
  \vspace{8pt}
  \begin{tcolorbox}[enhanced,colback=poporange,colframe=popink,boxrule=2.8pt,arc=13pt,
      drop shadow=popink,left=18pt,right=18pt,top=12pt,bottom=13pt,after skip=10pt]
    \poppill{#2 \textbullet\ #3}{popink}{popcream}\hspace{5pt}\poppill{#4}{white}{popink}\par\vspace{8pt}
    {\popdisplay\fontsize{14}{14}\selectfont\color{popink}LEÇON #5}\par\vspace{4pt}
    {\raggedright\hyphenpenalty=10000\exhyphenpenalty=10000\popdisplay\fontsize{25}{27}\selectfont\color{popcream}\MakeUppercase{#1}\par}
    \vspace{8pt}
    {\popxbold\fontsize{9.5}{9.5}\selectfont\color{popink}\MakeUppercase{Durée conseillée : #6}}
  \end{tcolorbox}
  \begin{tcolorbox}[enhanced,colback=white,colframe=popink,boxrule=2.2pt,arc=10pt,
      drop shadow=popink,left=16pt,right=16pt,top=10pt,bottom=10pt,after skip=8pt]
    \poppill{Dans cette leçon}{poppurple}{white}\par\vspace{5pt}
    \popsemi\fontsize{9.3}{12.6}\selectfont\color{popink}#7
  \end{tcolorbox}
  \noindent\begin{minipage}[t]{0.487\linewidth}
    \begin{tcolorbox}[enhanced,colback=popgreen,colframe=popink,boxrule=2.2pt,arc=10pt,
        drop shadow=popink,left=11pt,right=11pt,top=10pt,bottom=10pt,height=3.1cm,valign=center]
      \tcbox[colback=white,colframe=popink,boxrule=1.3pt,arc=5pt,boxsep=0pt,nobeforeafter,
        left=3pt,right=3pt,top=3pt,bottom=3pt,valign=center]{\qrcode[height=1.9cm]{https://play.google.com/store/apps/details?id=com.luvvix.onbuch&hl=fr}}%
      \hspace{8pt}\begin{minipage}[c]{0.5\linewidth}
        {\popdisplay\fontsize{11}{12.5}\selectfont\color{white}\MakeUppercase{Télécharge OnBuch}}\par\vspace{4pt}
        {\raggedright\popsemi\fontsize{8.4}{11}\selectfont\color{white}Gratuit \textbullet\ Google Play\\Tuteur IA, annales, concours, résultats\par}
      \end{minipage}
    \end{tcolorbox}
  \end{minipage}\hfill
  \begin{minipage}[t]{0.487\linewidth}
    \begin{tcolorbox}[enhanced,colback=poppurple,colframe=popink,boxrule=2.2pt,arc=10pt,
        drop shadow=popink,left=11pt,right=11pt,top=10pt,bottom=10pt,height=3.1cm,valign=center]
      \tikz[baseline=-0.7ex]\node[circle,fill=white,draw=popink,line width=1.3pt,minimum size=1.6cm,inner sep=0]{\popdisplay\fontsize{24}{24}\selectfont\color{poppurple}+};%
      \hspace{8pt}\begin{minipage}[c]{0.6\linewidth}
        {\popdisplay\fontsize{11}{12.5}\selectfont\color{white}\MakeUppercase{Passe à OnBuch+}}\par\vspace{4pt}
        {\raggedright\popsemi\fontsize{8.4}{11}\selectfont\color{white}Tous les cours signés, Tuteur IA illimité, fiches PDF — onglet OnBuch+ de l'app\par}
      \end{minipage}
    \end{tcolorbox}
  \end{minipage}
  \vspace{8pt}
  \popcontenu
  \vfill
  \signatureonbuch
  \restoregeometry
  \newpage
}

% Rangée « Ce cours contient » (page de garde)
\newcommand{\poptile}[3]{% #1 couleur #2 gros texte #3 légende
  \begin{tcolorbox}[enhanced,colback=#1,colframe=popink,boxrule=2pt,arc=9pt,shadow={2.5pt}{-2.5pt}{0pt}{popink},
      left=6pt,right=6pt,top=9pt,bottom=9pt,halign=center,valign=center,height=1.75cm,width=\linewidth,nobeforeafter]
    {\popdisplay\fontsize{12.5}{13}\selectfont\color{white}\MakeUppercase{#2}}\par\vspace{3pt}
    {\centering\popsemi\fontsize{7.8}{9.5}\selectfont\color{white}#3\par}
  \end{tcolorbox}}
\newcommand{\popcontenu}{%
  \par\noindent{\popxboldls\fontsize{7.5}{7.5}\selectfont\color{popmuted}CE COURS SIGNÉ CONTIENT}\par\vspace{7pt}
  \noindent\begin{minipage}[t]{0.235\linewidth}\poptile{popblue}{Cours}{complet, illustré, pas à pas}\end{minipage}\hfill
  \begin{minipage}[t]{0.235\linewidth}\poptile{poporange}{Méthodes}{et exercices résolus}\end{minipage}\hfill
  \begin{minipage}[t]{0.235\linewidth}\poptile{poppink}{Exercices}{type examen + corrigés}\end{minipage}\hfill
  \begin{minipage}[t]{0.235\linewidth}\poptile{popgreen}{Fiche bilan}{l'essentiel à retenir}\end{minipage}\par}

% Sommaire
\newtcolorbox{sommaire}{enhanced,colback=white,colframe=popink,boxrule=2.2pt,arc=10pt,
  drop shadow=popink,left=18pt,right=18pt,top=13pt,bottom=13pt,before skip=6pt,after skip=20pt,
  before upper={\poppill{Sommaire}{poppurple}{white}\par\vspace{9pt}\popsemi\fontsize{10.6}{14.5}\selectfont\color{popink}}}

% Signature de fin de cours — poussée tout en bas de la dernière page
\newcommand{\findecours}{%
  \par\vfill\nopagebreak
  \begin{center}\popsquiggle{poporange}\end{center}\vspace{4pt}
  \signatureonbuch
}
\AtBeginDocument{\def\llbracket{\mathopen{[\mkern-3.5mu[}}\def\rrbracket{\mathclose{]\mkern-3.5mu]}}} % intervalles d'entiers (glyphe ⟦ absent de la police)
% Glyphes absents de Fira Math : redéfinis à partir de glyphes disponibles
\AtBeginDocument{%
  \def\setminus{\mathbin{\backslash}}\let\smallsetminus\setminus
  \def\vdots{\mathord{\vbox{\baselineskip3.2pt\lineskiplimit0pt\kern4pt\hbox{.}\hbox{.}\hbox{.}}}}%
  \def\ddots{\mathinner{\mkern1mu\raise6.5pt\hbox{.}\mkern2mu\raise3.6pt\hbox{.}\mkern2mu\raise0.7pt\hbox{.}\mkern1mu}}%
  \def\triangle{\mathord{\scalebox{0.9}{$\symup{\Delta}$}}}%
  \def\nabla{\mathord{\raisebox{\depth}{\scalebox{1}[-1]{$\symup{\Delta}$}}}}%
  \def\checkmark{\popcheck{popdark}}%
  \def\star{\mathbin{\ast}}\def\bigstar{\mathord{\ast}}%
  \def\diamond{\mathbin{\scalebox{0.6}{$\Diamond$}}}%
  \def\bigcup{\mathop{\mathchoice{\raisebox{-0.3em}{\scalebox{1.7}{$\cup$}}}{\scalebox{1.25}{$\cup$}}{\cup}{\cup}}\displaylimits}%
  \def\bigcap{\mathop{\mathchoice{\raisebox{-0.3em}{\scalebox{1.7}{$\cap$}}}{\scalebox{1.25}{$\cap$}}{\cap}{\cap}}\displaylimits}%
  \def\lesseqqgtr{\mathrel{\lessgtr}}\def\bigoplus{\mathop{\oplus}\displaylimits}%
}
\providecommand{\ssi}{si et seulement si} % abréviation fréquente
\AtBeginDocument{\providecommand{\wideparen}{\overparen}} % arc de cercle

%% FIGFIX-BEGIN v5 — correctifs globaux des figures (docs/onbuch-generation/tools/figfix)
\usepackage{adjustbox}\usepackage{etoolbox}\tcbuselibrary{raster}
% toute figure TikZ/pgfplots plus large que la ligne est réduite à la largeur disponible
\BeforeBeginEnvironment{tikzpicture}{\begin{adjustbox}{max width=\linewidth}}
\AfterEndEnvironment{tikzpicture}{\end{adjustbox}}
% légendes pgfplots lisibles par-dessus les courbes
\pgfplotsset{every axis/.append style={legend style={fill=white,fill opacity=0.92,text opacity=1,draw=black!15,inner sep=3pt}}}
% étiquettes d'arêtes (syntaxe edge["..."]) avec fond blanc
\tikzset{every edge quotes/.append style={fill=white,inner sep=1.2pt}}
%% FIGFIX-END
\begin{document}

\pagegarde{Les bruits et les nuisances sonores}{Allemand}{1ère A}{Chapitre 3 : Environnement, santé et bien-être}{ALA13}{2 h}{Cette leçon de type « texte » propose la lecture et le commentaire d'un texte allemand sur les bruits et les nuisances sonores, leurs causes et leurs effets sur la santé. Elle permet d'acquérir le vocabulaire thématique (der Lärm, laut, leise, stören, die Ruhe...), de maîtriser les subordonnées temporelles avec wenn/als, les propositions relatives simples et le prétérit de sein, haben et des verbes modaux.
\vspace{4pt}
\begin{multicols}{2}\small
\begin{itemize}
\item À la fin de cette leçon, je sais comprendre globalement puis en détail un texte allemand sur les bruits et les nuisances sonores.
\item Je sais employer correctement le vocabulaire du bruit et du silence (der Lärm, laut, leise, die Nachbarn, stören, die Ruhe) avec articles et pluriels.
\item Je sais construire des subordonnées temporelles avec wenn et als et les distinguer sans erreur.
\item Je sais former et placer correctement une proposition relative simple avec der, die, das.
\item Je sais conjuguer sein, haben et les verbes modaux au prétérit et les utiliser dans un récit au passé.
\item Je sais répondre par écrit et à l'oral à des questions de compréhension sur un texte.
\end{itemize}\end{multicols}}

\begin{sommaire}
\begin{multicols}{2}
\begin{enumerate}
\item Texte support : Lärm in der Stadt
\item Compréhension du texte
\item Lexique thématique : der Lärm und die Ruhe
\item Grammaire 1 : les subordonnées temporelles (wenn, als)
\item Grammaire 2 : les propositions relatives simples
\item Conjugaison : prétérit de sein, haben et des modaux
\item Production guidée et méthode du commentaire
\item Activité d'intégration
\item Méthodes et exercices résolus
\item Exercices
\item Corrigés des exercices
\item Fiche bilan
\end{enumerate}
\end{multicols}
\end{sommaire}

\begin{objectifs}
\begin{itemize}
\item À la fin de cette leçon, je sais comprendre globalement puis en détail un texte allemand sur les bruits et les nuisances sonores.
\item Je sais employer correctement le vocabulaire du bruit et du silence (der Lärm, laut, leise, die Nachbarn, stören, die Ruhe) avec articles et pluriels.
\item Je sais construire des subordonnées temporelles avec wenn et als et les distinguer sans erreur.
\item Je sais former et placer correctement une proposition relative simple avec der, die, das.
\item Je sais conjuguer sein, haben et les verbes modaux au prétérit et les utiliser dans un récit au passé.
\item Je sais répondre par écrit et à l'oral à des questions de compréhension sur un texte.
\item Je sais résumer un texte en quelques phrases avec mes propres mots.
\item Je sais rédiger un court paragraphe guidé (plainte, description, conseil) sur le thème du bruit.
\end{itemize}
\end{objectifs}

\begin{prerequis}
\begin{itemize}
\item La phrase déclarative et la place du verbe (V2) en allemand (Seconde)
\item Les articles définis et indéfinis au nominatif et à l'accusatif
\item Le présent des verbes forts et faibles, y compris sein et haben
\item Le vocabulaire de la ville, de la maison et de la vie quotidienne (Seconde)
\item Notion de subordonnée introduite par weil ou dass (début d'année de Première)
\end{itemize}
\end{prerequis}

\newpage

\coursec{Texte support : Lärm in der Stadt}

\courssub{Situation d'accroche et problématique}

Il est 23 heures à Yaoundé. Demain, tu as cours à 7 h 30, mais ton voisin organise une fête : la musique est très forte, les invités rient, les motos passent et repassent. Tu te retournes dans ton lit, tu as mal à la tête, tu ne peux pas dormir. Le lendemain, tu es fatigué en classe et tu ne peux pas suivre la leçon d'allemand...

En Allemagne, ce problème est pris \textbf{très au sérieux} : il existe des \all{Ruhezeiten} (heures de silence) imposées par la loi, et les voisins peuvent porter plainte pour \all{Lärmbelästigung} (nuisance sonore). Le bruit n'est pas seulement un désagrément : c'est un vrai problème de \textbf{santé publique}.

\textbf{Problématique :} Comment parler du bruit, décrire ses sources et ses effets sur la santé, et exprimer un souhait de calme en allemand ? C'est ce que nous allons apprendre à travers un texte authentique : le témoignage d'une habitante de Munich qui souffre du bruit de sa rue.

\begin{savaistu}[Le saviez-vous ? Les « Ruhezeiten » en Allemagne]
En Allemagne, la loi impose des \all{Ruhezeiten} (heures de repos) : le silence est obligatoire la nuit, en général de 22 h à 6 h, ainsi que le dimanche et les jours fériés. Tondre sa pelouse (\all{den Rasen mähen}) un dimanche peut coûter une amende ! Beaucoup d'immeubles ont aussi une \all{Hausordnung} (règlement intérieur) qui interdit la musique forte entre 13 h et 15 h. Le mot \all{Lärmbelästigung} signifie « nuisance sonore » : c'est un motif légal de plainte contre un voisin. Remarque au passage la construction typique de l'allemand : \all{Lärm} (bruit) + \all{Belästigung} (gêne) = un mot composé très clair.
\end{savaistu}

\courssub{Comment lire un texte allemand : la méthode en trois étapes}

Avant de découvrir le texte, rappelons la bonne démarche de lecture. Ne commence jamais par traduire mot à mot !

\begin{methode}[Lire un texte en trois étapes]
\begin{enumerate}
\item \textbf{Lecture globale} (sans dictionnaire) : lis le texte une première fois rapidement. Objectif : saisir l'idée générale. Qui parle ? De quoi parle-t-on ? Où et quand ? Ne t'arrête pas sur les mots inconnus : devine leur sens grâce au contexte et aux mots transparents (\all{Musik}, \all{Stress}, \all{Party}...).
\item \textbf{Lecture détaillée} (avec le glossaire) : relis phrase par phrase en utilisant le glossaire. Repère les informations précises : les causes, les conséquences, les moments, les personnes.
\item \textbf{Relecture grammaticale} : observe les structures de la langue : les temps des verbes, la place du verbe, les subordonnées, les propositions relatives. C'est là que le texte devient un modèle pour ta propre expression.
\end{enumerate}
\end{methode}

\courssub{Première lecture globale : le texte}

Lis le texte suivant une première fois, sans t'aider du glossaire. Cherche simplement à répondre à cette question : \emph{de quoi se plaint la narratrice ?}

\begin{textebox}[Zu viel Lärm in unserer Straße]
Als ich letztes Jahr nach München zog, war unsere Straße sehr ruhig. Am Abend konnte man die Vögel hören, und ich habe immer gut geschlafen. Aber jetzt ist alles anders : unsere Straße ist viel zu laut!

Zuerst gibt es eine große Baustelle neben unserem Haus. Die Maschinen fangen jeden Morgen um sechs Uhr an, auch am Samstag. Wenn die Maschinen arbeiten, kann ich nicht schlafen. Dann kommt der Verkehr : viele Autos, Busse und Motorräder fahren den ganzen Tag durch unsere Straße. Wenn ich am Fenster sitze und lerne, höre ich nur Hupen und Motoren.

Meine Nachbarn, die sehr laute Musik hören, stören mich auch. Der junge Mann, der über mir wohnt, feiert fast jedes Wochenende eine Party. Seine Gäste sprechen laut, und die Musik dauert oft bis zwei Uhr nachts.

Der Lärm hat Folgen für meine Gesundheit : ich hatte gestern starke Kopfschmerzen und ich war sehr müde. Ich konnte nicht lernen und ich war in der Schule nervös. Meine Mutter sagt, dass ich mehr schlafen muss, aber wie? Ich wünsche mir mehr Ruhe! Vielleicht schreibe ich einen Brief an die Hausverwaltung. Oder ich ziehe um — aufs Land, wo es still ist!
\end{textebox}

\textbf{Idée générale :} une jeune habitante de Munich décrit les bruits de sa rue (travaux, circulation, voisins) et leurs effets sur sa santé (fatigue, maux de tête, nervosité). Elle rêve de calme et envisage d'écrire à la régie ou de déménager à la campagne.

\courssub{Glossaire bilingue du texte}

Voici les mots nouveaux du texte, avec l'article et le pluriel des noms. Apprends-les toujours \textbf{avec leur article} : c'est la clé pour maîtriser les cas allemands.

\begin{center}
\begin{tabularx}{\linewidth}{|l|l|X|}
\hline
\rowcolor{popblueL} \textbf{Deutsch} & \textbf{Pluriel} & \textbf{Français} \\
\hline
der Lärm & (nur Singular) & le bruit (nuisible) \\
\hline
das Geräusch & die Geräusche & le bruit (petit son, neutre) \\
\hline
die Ruhe & (nur Singular) & le calme, le silence \\
\hline
die Baustelle & die Baustellen & le chantier \\
\hline
der Verkehr & (nur Singular) & la circulation \\
\hline
der Nachbar & die Nachbarn & le voisin \\
\hline
die Maschine & die Maschinen & la machine \\
\hline
die Hupe & die Hupen & le klaxon \\
\hline
das Motorrad & die Motorräder & la moto \\
\hline
der Motor & die Motoren & le moteur \\
\hline
der Vogel & die Vögel & l'oiseau \\
\hline
der Gast & die Gäste & l'invité \\
\hline
die Kopfschmerzen & (nur Plural) & les maux de tête \\
\hline
die Gesundheit & (nur Singular) & la santé \\
\hline
die Folge & die Folgen & la conséquence \\
\hline
die Hausverwaltung & die Hausverwaltungen & la régie (gestion de l'immeuble) \\
\hline
das Land & die Länder & la campagne (ici) ; le pays \\
\hline
\end{tabularx}
\end{center}

\begin{center}
\begin{tabularx}{\linewidth}{|l|X|}
\hline
\rowcolor{popblueL} \textbf{Verben und Adjektive} & \textbf{Français} \\
\hline
stören & déranger, gêner \\
\hline
laut / leise & fort (bruyant) / doux (silencieux) \\
\hline
ruhig / still & calme / silencieux \\
\hline
anfangen (fängt an, fing an, hat angefangen) & commencer \\
\hline
umziehen (zieht um, zog um, ist umgezogen) & déménager \\
\hline
sich wünschen & souhaiter \\
\hline
feiern & fêter, faire la fête \\
\hline
dauern & durer \\
\hline
müde / nervös & fatigué / nerveux \\
\hline
\end{tabularx}
\end{center}

\begin{definition}[Lärm, Geräusch ou Ruhe ? Trois mots à ne pas confondre]
\begin{itemize}
\item \all{der Lärm} (uniquement au singulier) = le bruit \textbf{nuisible}, désagréable : \all{Der Lärm stört mich.} « Le bruit me dérange. »
\item \all{das Geräusch}, \all{-e} = un bruit neutre, un petit son : \all{Ich höre ein Geräusch in der Küche.} « J'entends un bruit dans la cuisine. »
\item \all{die Ruhe} (uniquement au singulier) = le calme, le silence : \all{Ich brauche Ruhe.} « J'ai besoin de calme. »
\end{itemize}
Attention : pour dire « le bruit » en général (celui qui gêne), on utilise presque toujours \all{der Lärm}, jamais \all{der Ton} ou \all{der Klang} dans ce contexte.
\end{definition}

\begin{attention}[Faux amis et pièges de vocabulaire]
\begin{itemize}
\item \all{laut} ne veut pas dire « lourd » mais « fort, bruyant » : \all{Die Musik ist zu laut.} « La musique est trop forte. »
\item \all{leise} = « doux, faible » (pour le son) : \all{Sprich bitte leise!} « Parle doucement, s'il te plaît ! »
\item \all{der Verkehr} = la circulation, pas « le commerce ».
\item \all{das Land} : au sens de « campagne », on dit \all{auf dem Land} (datif, lieu) ou \all{aufs Land} (accusatif, mouvement) : \all{Ich ziehe aufs Land.} « Je déménage à la campagne. »
\item \all{die Ruhe} est féminin : on dit \all{die Ruhe}, jamais \all{der Ruhe} au nominatif.
\item N'oublie pas la majuscule à tous les noms : \all{der Lärm}, \all{die Baustelle}, \all{der Schlaf} (le sommeil).
\end{itemize}
\end{attention}

\courssub{Carte mentale du vocabulaire : der Lärm}

Pour mémoriser le vocabulaire, rien de tel qu'une carte mentale : au centre le mot-clé, autour les sources du bruit (\all{Quellen}) et ses conséquences (\all{Folgen}).

\begin{popfigure}
\begin{center}
\begin{tikzpicture}[
  mindmap,
  grow cyclic,
  every node/.style={concept, align=center, font=\small},
  root concept/.append style={concept color=poporange, font=\large\bfseries},
  level 1/.append style={level distance=3.2cm, sibling angle=90},
  level 2/.append style={level distance=2.4cm, sibling angle=45, font=\footnotesize},
]
\node[root concept]{der Lärm}
  child[concept color=popblue] { node[align=center]{Quellen\\ (sources)}
    child { node[concept color=popblueL, text=popdark] {die Baustelle} }
    child { node[concept color=popblueL, text=popdark] {der Verkehr} }
    child { node[concept color=popblueL, text=popdark] {die Nachbarn} }
    child { node[concept color=popblueL, text=popdark] {die Musik} }
  }
  child[concept color=popgreen] { node[align=center]{Folgen\\ (effets)}
    child { node[concept color=popgreenL, text=popdark] {der Stress} }
    child { node[concept color=popgreenL, text=popdark, align=center]{die Kopf-\\schmerzen} }
    child { node[concept color=popgreenL, text=popdark] {kein Schlaf} }
  };
\end{tikzpicture}
\end{center}
\legende{Carte mentale : les sources du bruit et ses effets sur la santé.}
\end{popfigure}

\courssub{Deuxième lecture détaillée : causes et conséquences}

Reprenons le texte phrase par phrase. Consigne de travail en classe : \textbf{souligne} les causes du bruit et \textbf{entoure} les conséquences sur la santé.

\textbf{Les causes du bruit (à souligner) :}
\begin{itemize}
\item \all{eine große Baustelle neben unserem Haus} — un grand chantier à côté de la maison ;
\item \all{der Verkehr : viele Autos, Busse und Motorräder} — la circulation : voitures, bus, motos ;
\item \all{Meine Nachbarn, die sehr laute Musik hören} — les voisins qui écoutent de la musique trop forte ;
\item \all{Der junge Mann, der über mir wohnt, feiert fast jedes Wochenende eine Party} — le jeune voisin du dessus qui fait la fête presque chaque week-end.
\end{itemize}

\textbf{Les conséquences sur la santé (à entourer) :}
\begin{itemize}
\item \all{ich kann nicht schlafen} — je ne peux pas dormir ;
\item \all{ich hatte gestern starke Kopfschmerzen} — j'ai eu de forts maux de tête hier ;
\item \all{ich war sehr müde} — j'étais très fatiguée ;
\item \all{ich konnte nicht lernen} — je n'ai pas pu travailler/réviser ;
\item \all{ich war in der Schule nervös} — j'étais nerveuse à l'école.
\end{itemize}

Observe la progression logique du texte : \all{zuerst} (d'abord) → \all{dann} (ensuite) → \all{auch} (aussi) → les \all{Folgen} (conséquences) → le souhait final. Ces connecteurs simples te serviront pour tes propres productions écrites.

\begin{exoresolu}[Questions de compréhension]
Réponds en allemand, par des phrases complètes :
\begin{enumerate}
\item \all{Wann fangen die Maschinen an?}
\item \all{Warum kann die Erzählerin nicht lernen?}
\item \all{Was wünscht sie sich?}
\end{enumerate}
\tcblower
\textbf{Solutions :}
\begin{enumerate}
\item \all{Die Maschinen fangen jeden Morgen um sechs Uhr an.} « Les machines commencent chaque matin à six heures. » (Attention : \all{anfangen} est un verbe à particule séparable : \all{fängt ... an}.)
\item \all{Sie kann nicht lernen, weil sie müde ist und Kopfschmerzen hat.} « Elle ne peut pas travailler parce qu'elle est fatiguée et qu'elle a mal à la tête. » (Après \all{weil}, le verbe conjugué va à la fin.)
\item \all{Sie wünscht sich mehr Ruhe.} « Elle souhaite plus de calme. » (Verbe réfléchi : \all{sich wünschen} + accusatif.)
\end{enumerate}
\end{exoresolu}

\courssub{Relecture grammaticale : observer les structures}

Troisième étape de la méthode : le texte est un modèle de grammaire. Observons trois phénomènes que nous étudierons en détail dans les leçons de grammaire qui suivent.

\textbf{1. Deux subordonnées temporelles avec \all{wenn} :}
\begin{itemize}
\item \all{Wenn die Maschinen arbeiten, kann ich nicht schlafen.} « Quand les machines fonctionnent, je ne peux pas dormir. »
\item \all{Wenn ich am Fenster sitze und lerne, höre ich nur Hupen und Motoren.} « Quand je suis assise à la fenêtre et que je travaille, je n'entends que des klaxons et des moteurs. »
\end{itemize}
Observe : dans la subordonnée introduite par \all{wenn}, le verbe conjugué va \textbf{à la fin} ; quand la subordonnée est en tête de phrase, le verbe de la principale vient juste après la virgule, en première position (\all{..., kann ich ...} / \all{..., höre ich ...}).

\textbf{2. Une subordonnée avec \all{als} :}
\begin{itemize}
\item \all{Als ich letztes Jahr nach München zog, war unsere Straße sehr ruhig.} « Quand j'ai emménagé à Munich l'année dernière, notre rue était très calme. »
\end{itemize}
\all{als} s'emploie pour un événement \textbf{unique au passé} ; \all{wenn} pour une répétition, une habitude ou le présent. Comparaison :

\begin{center}
\begin{tabularx}{\linewidth}{|X|X|X|}
\hline
\rowcolor{popblueL} \textbf{Français} & \textbf{Deutsch} & \textbf{Remarque} \\
\hline
Quand j'étais enfant (répétition, passé) & \all{Wenn ich ein Kind war, ...} & habitude au passé : \all{wenn} \\
\hline
Quand je suis arrivé (une fois, passé) & \all{Als ich ankam, ...} & événement unique : \all{als} \\
\hline
Quand il pleut (présent, général) & \all{Wenn es regnet, ...} & présent ou futur : \all{wenn} \\
\hline
\end{tabularx}
\end{center}

\textbf{3. Deux propositions relatives :}
\begin{itemize}
\item \all{Meine Nachbarn, \textbf{die} sehr laute Musik hören, stören mich auch.} « Mes voisins, qui écoutent une musique très forte, me dérangent aussi. » (\all{die} = qui, pour \all{die Nachbarn}, pluriel.)
\item \all{Der junge Mann, \textbf{der} über mir wohnt, feiert fast jedes Wochenende eine Party.} « Le jeune homme qui habite au-dessus de chez moi fait la fête presque chaque week-end. » (\all{der} = qui, pour \all{der Mann}, masculin singulier.)
\end{itemize}
Le pronom relatif allemand reprend le \textbf{genre et le nombre} du nom auquel il se réfère, et le verbe de la relative se place à la fin.

\textbf{4. Trois verbes au prétérit :}
\begin{itemize}
\item \all{war} (prétérit de \all{sein}, être) : \all{ich war sehr müde} ;
\item \all{hatte} (prétérit de \all{haben}, avoir) : \all{ich hatte starke Kopfschmerzen} ;
\item \all{konnte} (prétérit du modal \all{können}, pouvoir) : \all{ich konnte nicht lernen}.
\end{itemize}
Ces trois prétérits sont les plus fréquents de la langue parlée et écrite : retiens-les par cœur dès maintenant. Remarque aussi le parfait \all{ich habe immer gut geschlafen} (verbe fort \all{schlafen} : participe \all{geschlafen}).

\begin{aretenir}[L'essentiel de cette section]
\begin{itemize}
\item \all{der Lärm} = le bruit nuisible ; \all{die Ruhe} = le calme ; \all{stören} = déranger.
\item Les sources du bruit : \all{die Baustelle}, \all{der Verkehr}, \all{die Nachbarn}, \all{laute Musik}.
\item Les effets : \all{Kopfschmerzen haben}, \all{müde sein}, \all{nicht schlafen können}, \all{nervös sein}.
\item \all{wenn} + verbe à la fin (présent ou habitude) ; \all{als} + verbe à la fin (événement unique au passé).
\item Prétérits à connaître : \all{war}, \all{hatte}, \all{konnte}.
\end{itemize}
\end{aretenir}

\courssub{Lecture à voix haute : travaillons la prononciation}

Lis le texte à voix haute après le professeur. Quelques points de prononciation à soigner (notation en lettres françaises) :

\begin{itemize}
\item \all{Lärm} se prononce « lèrm » : le \textbf{ä} long se dit « è » ouvert.
\item \all{Geräusch} : le \textbf{ch} après une voyelle antérieure (ä, i, e) se prononce comme un « ch » doux, proche de « hy » : « gue-roïch ». Après a, o, u, au, le \textbf{ch} est guttural : \all{nach} « rakh », \all{auch} « aoukh ».
\item La lettre \textbf{w} se prononce « v » : \all{Woche} « vokh-e », \all{Wochenende} « vokh-n-end-e », \all{wohnen} « vo-nen ».
\item La lettre \textbf{v} se prononce souvent « f » : \all{Verkehr} « fèr-kèr », \all{Vögel} « feu-guel ».
\item \all{sechs} se prononce « zeks » ; \all{zwei} « tsvaï » ; \all{zuerst} « tsou-èrst ».
\item \all{Straße} : le \textbf{ß} se prononce « s » sec, et le \textbf{a} est long : « chtra-sse ».
\end{itemize}

\begin{experience}[Activité orale : la chaîne des bruits]
Par groupes de quatre. Chaque élève répète la phrase de son camarade et ajoute une source de bruit, sur le modèle : \all{In meiner Straße ist es laut : es gibt eine Baustelle.} — \all{In meiner Straße ist es laut : es gibt eine Baustelle und viele Autos.} — et ainsi de suite avec \all{Motorräder}, \all{laute Musik}, \all{Nachbarn}, \all{Hupen}... Le groupe qui fait la chaîne la plus longue sans erreur d'article gagne !
\end{experience}

\begin{exercice}[À toi de jouer]
\begin{enumerate}
\item Traduis en français : \all{Wenn die Maschinen arbeiten, kann ich nicht schlafen.}
\item Complète avec \all{wenn} ou \all{als} : \all{... ich letztes Jahr nach Yaoundé kam, war ich müde.} / \all{... mein Nachbar Musik hört, kann ich nicht lernen.}
\item Trouve dans le texte les deux propositions relatives et souligne le pronom relatif.
\item Réponds par une phrase complète : \all{Was macht der junge Mann fast jedes Wochenende?}
\end{enumerate}
\end{exercice}

\begin{corrige}[Exercice « À toi de jouer »]
\begin{enumerate}
\item « Quand les machines fonctionnent, je ne peux pas dormir. »
\item \all{\textbf{Als} ich letztes Jahr nach Yaoundé kam, war ich müde.} (événement unique au passé) ; \all{\textbf{Wenn} mein Nachbar Musik hört, kann ich nicht lernen.} (répétition, présent).
\item \all{Meine Nachbarn, \textbf{die} sehr laute Musik hören, stören mich auch.} et \all{Der junge Mann, \textbf{der} über mir wohnt, feiert fast jedes Wochenende eine Party.}
\item \all{Er feiert fast jedes Wochenende eine Party.} « Il fait la fête presque chaque week-end. »
\end{enumerate}
\end{corrige}

\coursec{Compréhension du texte}

Après la lecture du texte support „Lärm in der Stadt“ (section précédente), nous passons maintenant à l'exploitation systématique de ce document. La compréhension suit une progression classique : du global vers le détaillé, puis vers la reformulation et l'expression personnelle. Chaque étape est accompagnée d'une méthode rigoureuse et d'exemples corrigés pour vous permettre d'acquérir les automatismes attendus au baccalauréat.

\courssub{Compréhension globale — questions de repérage}

Avant d'entrer dans le détail, répondez aux questions suivantes en français. Elles visent à vérifier que vous avez saisi la situation de communication dans son ensemble.

\begin{exercice}[Compréhension globale]
Répondez en français par des phrases complètes.
\begin{enumerate}
\item De quoi parle le texte ? Donnez le thème général en une phrase.
\item Qui est le narrateur ? Quels indices du texte permettent de l'identifier (âge, situation, lieu) ?
\item Où et quand se déroule l'action ? Citez les éléments spatio-temporels explicites.
\item Quel est le problème principal évoqué par le narrateur ?
\item Quelle est l'intention communicative de ce texte (témoignage, plainte, description, argumentation) ?
\end{enumerate}
\end{exercice}

\begin{corrige}[Exercice 1 — Compréhension globale]
\begin{enumerate}
\item Le texte traite des nuisances sonores urbaines (bruit de chantier, circulation, voisins) et de leurs conséquences sur la santé et le travail scolaire d'une lycéenne.
\item Le narrateur est une élève de Première (lycéenne) qui prépare son baccalauréat. Indices : „\all{ich muss für das Abitur lernen}“, „\all{mein Zimmer}“, „\all{Hausaufgaben}“.
\item L'action se déroule à Yaoundé (Cameroun), dans le quartier du narrateur, actuellement et de manière récurrente : „\all{jeden Morgen}“, „\all{seit drei Wochen}“, „\all{um sechs Uhr}“.
\item Le problème principal : le bruit constant d'un chantier de construction qui commence trop tôt le matin l'empêche de dormir, de se concentrer et lui donne des maux de tête.
\item C'est un témoignage personnel à visée descriptive et plaintive : la narratrice décrit sa situation vécue pour faire prendre conscience de l'impact du bruit sur la santé et la scolarité.
\end{enumerate}
\end{corrige}

\courssub{Compréhension détaillée — questions en allemand}

Les questions détaillées sont formulées en allemand. Vous devez y répondre \emph{en allemand}, par des phrases complètes, en reprenant le vocabulaire du texte mais avec votre propre structure syntaxique.

\begin{exercice}[Compréhension détaillée]
Répondez en allemand, phrases complètes obligatoires.
\begin{enumerate}
\item Wann fängt die Baustelle jeden Morgen an?
\item Was für Geräusche hört die Erzählerin außer dem Bagger?
\item Warum konnte die Erzählerin gestern Abend nicht lernen?
\item Wie fühlt sich die Erzählerin wegen des Lärms? Nennen Sie zwei Symptome.
\item Was macht die Erzählerin gegen den Lärm? Nennen Sie zwei Maßnahmen.
\item Was fordern die Nachbarn von der Stadtverwaltung?
\end{enumerate}
\end{exercice}

\begin{methode}[Répondre à une question de compréhension en allemand]
\textbf{Étape 1 — Repérer l'information dans le texte.} Soulignez la phrase ou le segment qui contient la réponse. Notez le numéro de ligne.
\textbf{Étape 2 — Identifier la structure de la question.} Repérez le mot interrogatif (Wann, Was, Warum, Wie, Wer, Wo) et le verbe conjugué.
\textbf{Étape 3 — Construire la réponse.}
\begin{itemize}
\item Reprenez le sujet de la question (ou le sujet logique).
\item Utilisez le verbe au même temps que dans la question (souvent le même que dans le texte).
\item Remplacez l'élément interrogé par l'information trouvée.
\item Adaptez la structure : en allemand, le verbe conjugué reste en \textbf{deuxième position} dans la phrase principale.
\end{itemize}
\textbf{Étape 4 — Vérifier.} Lisez votre phrase à voix haute : est-elle correcte (ordre des mots, cas, accord) ? Répond-elle exactement à la question posée ?
\end{methode}

\begin{exemplebox}[Modèles de réponses détaillées]
\begin{itemize}
\item \textbf{Frage:} Wann fängt die Baustelle an? \\
\textbf{Antwort:} \all{Die Baustelle fängt jeden Morgen um sechs Uhr an.} — « Le chantier commence chaque matin à six heures. »
\item \textbf{Frage:} Was stört die Erzählerin am meisten? \\
\textbf{Antwort:} \all{Der Lärm der Baustelle und der laute Verkehr stören sie am meisten.} — « Le bruit du chantier et la circulation bruyante la dérangent le plus. »
\item \textbf{Frage:} Warum konnte sie nicht lernen? \\
\textbf{Antwort:} \all{Sie konnte nicht lernen, weil der Bagger zu laut war und sie starke Kopfschmerzen hatte.} — « Elle n'a pas pu apprendre parce que la pelleteuse était trop bruyante et qu'elle avait de forts maux de tête. »
\item \textbf{Frage:} Was macht sie gegen den Lärm? \\
\textbf{Antwort:} \all{Sie schließt die Fenster und trägt Ohrstöpsel.} — « Elle ferme les fenêtres et porte des bouchons d'oreilles. »
\end{itemize}
\end{exemplebox}

\begin{attention}[Piège fréquent : traduction mot à mot]
Ne traduisez \textbf{pas} mot à mot depuis le français. En allemand, l'ordre des mots change radicalement :
\begin{itemize}
\item \textbf{Français :} « Pourquoi n'a-t-elle pas pu apprendre ? » → \textbf{Allemand :} \all{Warum konnte sie nicht lernen?} (verbe en position 2, sujet après le verbe).
\item \textbf{Français :} « Elle ferme les fenêtres. » → \textbf{Allemand :} \all{Sie schließt die Fenster.} (verbe en position 2).
\item \textbf{Français :} « Parce que c'est trop bruyant. » → \textbf{Allemand :} \all{Weil es zu laut ist.} (verbe à la fin en subordonnée !).
\end{itemize}
Réutilisez le vocabulaire du texte (\all{Baustelle}, \all{Bagger}, \all{Kopfschmerzen}, \all{Ohrstöpsel}, \all{Ruhe}) mais construisez \emph{votre propre phrase} allemande.
\end{attention}

\begin{corrige}[Exercice 2 — Compréhension détaillée (propositions de réponses)]
\begin{enumerate}
\item \all{Die Baustelle fängt jeden Morgen um sechs Uhr an.}
\item \all{Sie hört den lauten Verkehr, die Hupen der Autos und die lauten Nachbarn.}
\item \all{Sie konnte nicht lernen, weil der Bagger sehr laut war und sie starke Kopfschmerzen hatte.}
\item \all{Sie fühlt sich müde und gestresst, und sie hat Kopfschmerzen.}
\item \all{Sie schließt die Fenster und trägt Ohrstöpsel.}
\item \all{Die Nachbarn fordern, dass die Stadtverwaltung die Baustelle später beginnen lässt und Ruhezeiten einhält.}
\end{enumerate}
\end{corrige}

\courssub{Exercice Vrai / Faux avec justification par citation}

Lisez les affirmations suivantes. Cochez \textbf{Richtig} (vrai) ou \textbf{Falsch} (faux) et justifiez \textbf{obligatoirement} par une citation exacte du texte (copiez le segment entre guillemets allemands „…“).

\begin{exercice}[Vrai / Faux — Justification textuelle]
\begin{enumerate}
\item Die Baustelle beginnt um sieben Uhr.
\item Die Erzählerin wohnt in einer ruhigen Seitenstraße.
\item Der Lärm verursacht bei der Erzählerin Schlafprobleme und Kopfschmerzen.
\item Die Nachbarn sind mit der Situation zufrieden.
\item Die Stadtverwaltung hat schon reagiert und die Baustelle gestoppt.
\end{enumerate}
\end{exercice}

\begin{corrige}[Exercice 4 — Vrai / Faux avec citations]
\begin{enumerate}
\item \textbf{Falsch.} — Citation : „\all{Die Baustelle fängt jeden Morgen um sechs Uhr an.}“
\item \textbf{Falsch.} — Citation : „\all{Ich wohne an einer Hauptstraße, wo der Verkehr schon früh morgens sehr laut ist.}“
\item \textbf{Richtig.} — Citations : „\all{Ich schlafe schlecht, weil der Lärm mich jede Nacht weckt.}“ et „\all{Gestern Abend hatte ich so starke Kopfschmerzen, dass ich nicht für das Abitur lernen konnte.}“
\item \textbf{Falsch.} — Citation : „\all{Die Nachbarn beschweren sich auch und fordern Ruhezeiten.}“
\item \textbf{Falsch.} — Citation : „\all{Bisher hat die Stadtverwaltung noch nicht reagiert.}“
\end{enumerate}
\end{corrige}

\courssub{Recherche d'équivalents lexicaux}

Retrouvez dans le texte les mots ou expressions allemands qui correspondent aux sens français suivants. Donnez le mot \textbf{avec son article} et, pour les noms, son \textbf{pluriel}.

\begin{exercice}[Équivalents lexicaux]
\begin{enumerate}
\item les voisins (nom masculin pluriel)
\item le mal de tête (nom féminin singulier / Pluraltantum)
\item le calme / la tranquillité (nom féminin singulier)
\item la pelleteuse / l'excavatrice (nom masculin singulier)
\item les bouchons d'oreilles (nom masculin pluriel / Pluraltantum)
\item se plaindre (verbe)
\item exiger / réclamer (verbe)
\item respecter (règles, délais) (verbe)
\end{enumerate}
\end{exercice}

\begin{corrige}[Exercice 3 — Équivalents lexicaux]
\begin{enumerate}
\item \all{die Nachbarn} (Pl. \all{die Nachbarn})
\item \all{die Kopfschmerzen} (Pluraltantum, pas de singulier usuel)
\item \all{die Ruhe} (pas de pluriel usuel)
\item \all{der Bagger} (Pl. \all{die Bagger})
\item \all{die Ohrstöpsel} (Pluraltantum)
\item \all{sich beschweren} (réfléchi + über + Akk.)
\item \all{fordern} (transitif + Akk.)
\item \all{einhalten} (séparable : \all{hält … ein})
\end{enumerate}
\end{corrige}

\courssub{Réformulation : résumé en trois phrases en allemand}

\begin{methode}[Rédiger un résumé en 3 phrases (niveau A2/B1)]
\textbf{Étape 1 — Sélectionner l'essentiel.} Gardez uniquement : \textbf{QUI} (la narratrice) + \textbf{QUOI} (bruit de chantier + circulation + voisins) + \textbf{OÙ} (Yaoundé, rue principale) + \textbf{CONSEQUENCE} (sommeil, maux de tête, apprentissage impossible) + \textbf{ACTION/DEMANDE} (fenêtres fermées, bouchons, plainte des voisins).
\textbf{Étape 2 — Choisir le temps.} Présent (description actuelle) ou Prétérit (récit). Au niveau A2, le \textbf{Présent} est plus sûr.
\textbf{Étape 3 — Rédiger 3 phrases maximum.} Utilisez des connecteurs logiques (\all{weil}, \all{deshalb}, \all{und}, \all{aber}). Évitez les subordonnées trop complexes.
\textbf{Étape 4 — Vérifier.} Chaque phrase commence par une majuscule, finit par un point, verbe en position 2 (principale) ou à la fin (subordonnée). Vocabulaire du texte réutilisé.
\end{methode}

\begin{corrige}[Résumé en 3 phrases — proposition corrigée]
\all{Die Schülerin wohnt in Yaoundé an einer lauten Hauptstraße, wo eine Baustelle jeden Morgen um sechs Uhr beginnt. Wegen des Lärms vom Bagger und des Verkehrs schläft sie schlecht und bekommt Kopfschmerzen, deshalb kann sie nicht für das Abitur lernen. Die Nachbarn fordern Ruhezeiten, aber die Stadtverwaltung hat noch nicht reagiert.}

\textbf{Analyse :}
\begin{itemize}
\item Phrase 1 : QUI + OÙ + QUOI (situation initiale).
\item Phrase 2 : CAUSE (Lärm) → CONSÉQUENCE (schläft schlecht, Kopfschmerzen) → CONSÉQUENCE FINALE (kann nicht lernen). Connecteur \all{deshalb}. \textbf{Attention :} \all{Wegen} + Génitif (\all{des Lärms}, \all{des Verkehrs}).
\item Phrase 3 : RÉACTION (Nachbarn fordern) + CONSTAT (Stadtverwaltung nicht reagiert). Contraste \all{aber}.
\end{itemize}
\end{corrige}

\courssub{Tableau des causes et conséquences (Ursachen / Folgen)}

Complétez le tableau ci-dessous en extrayant les informations du texte. Ce travail de structuration visuelle prépare l'expression écrite et orale.

\begin{popfigure}
\begin{tikzpicture}[
    node distance=0.8cm and 1.2cm,
    box/.style={rectangle, draw=popblue, thick, rounded corners=4pt, minimum height=1.1cm, text width=7.2cm, align=center, font=\small, fill=popblue!5},
    header/.style={rectangle, draw=popdark, thick, rounded corners=4pt, minimum height=1cm, text width=7.2cm, align=center, font=\small\bfseries, fill=popblue!30},
    arrow/.style={-Stealth, thick, poporange}
]
% En-têtes
\node[header, align=center] (u) at (0,0){\textbf{Ursachen (Causes)}\\ \textit{Was löst den Lärm aus?}};
\node[header, align=center] (f) at (16,0){\textbf{Folgen (Conséquences)}\\ \textit{Was bewirkt der Lärm?}};

% Causes
\node[box, align=center] (u1) [below=of u]{\all{die Baustelle (Bagger, LKWs)}\\\all{ab 6:00 Uhr morgens}};
\node[box, align=center] (u2) [below=of u1]{\all{der laute Verkehr (Autos, Hupen)}\\\all{auf der Hauptstraße}};
\node[box, align=center] (u3) [below=of u2]{\all{laute Nachbarn (Musik, Streit)}\\\all{in den Abendstunden}};
\node[box, align=center] (u4) [below=of u3]{\all{fehlende Ruhezeiten / keine Reaktion}\\\all{der Stadtverwaltung}};

% Conséquences
\node[box, align=center] (f1) [below=of f]{\all{Schlafstörungen}\\\all{(wird jede Nacht geweckt)}};
\node[box, align=center] (f2) [below=of f1]{\all{Kopfschmerzen, Müdigkeit,}\\\all{Stress, Reizbarkeit}};
\node[box, align=center] (f3) [below=of f2]{\all{Keine Konzentration möglich:}\\\all{Lernen für Abitur unmöglich}};
\node[box, align=center] (f4) [below=of f3]{\all{Maßnahmen der Betroffenen:}\\\all{Fenster zu, Ohrstöpsel,}\\\all{Beschwerde bei der Stadt}};

% Flèches de liaison (optionnel, pour montrer la causalité)
\draw[arrow] (u1.east) -- ++(0.8,0) |- (f1.west);
\draw[arrow] (u2.east) -- ++(0.8,0) |- (f2.west);
\draw[arrow] (u3.east) -- ++(0.8,0) |- (f2.west);
\draw[arrow] (u4.east) -- ++(0.8,0) |- (f4.west);
\end{tikzpicture}
\legende{Tableau Ursachen / Folgen à compléter d'après le texte „Lärm in der Stadt“.}
\end{popfigure}

\courssub{Expression orale guidée : comparaison Cameroun / Allemagne}

Pour ancrer le vocabulaire et la structure dans un échange authentique, menez l'activité suivante en binôme ou en groupe-classe.

\begin{experience}[Débat guidé : Le bruit chez nous et en Allemagne]
\textbf{Consigne :} En binôme, préparez un dialogue de 2 minutes. L'élève A vit à Yaoundé (ou Douala), l'élève B vit à Berlin (ou Munich). Vous vous parlez par visioconférence.

\textbf{Points à aborder (utilisez le vocabulaire de la leçon) :}
\begin{enumerate}
\item \textbf{Description du bruit :} Quelles sont les sources principales ? (chantiers, circulation, bendskins, bars, églises, marchés, \all{Baustelle}, \all{Verkehr}, \all{Nachbarn}, \all{Kirchenglocken}, \all{Markt})
\item \textbf{Horaires :} À quelle heure commence/ finit le bruit ? (\all{um sechs Uhr}, \all{abends}, \all{nachts}, \all{Ruhezeiten})
\item \textbf{Réglementation :} Existe-t-il des heures de calme légales ? (\all{Ruhezeiten}, \all{Nachtruhe 22--6 Uhr}, \all{Mittagsruhe}, \all{Sonntagsruhe})
\item \textbf{Conséquences personnelles :} Sommeil, concentration, santé. (\all{schlecht schlafen}, \all{Kopfschmerzen}, \all{sich nicht konzentrieren können})
\item \textbf{Solutions individuelles / collectives :} (\all{Fenster schließen}, \all{Ohrstöpsel}, \all{sich beschweren}, \all{Stadtverwaltung}, \all{Ordnungsamt})
\end{enumerate}

\textbf{Structures utiles (à afficher au tableau) :}
\begin{itemize}
\item \all{Bei mir in Yaoundé ist es …, weil …}
\item \all{In Deutschland gibt es strenge Ruhezeiten: …}
\item \all{Das stört mich am meisten: …}
\item \all{Ich fordere, dass …} (+ subordonnée \all{dass} verbe à la fin)
\item \all{Was machst du dagegen?}
\end{itemize}

\textbf{Évaluation par les pairs :} Après chaque dialogue, la classe note sur une grille : vocabulaire thématique (5), exactitude grammaticale (5), fluidité/prononciation (5), interaction (5).
\end{experience}

\begin{savaistu}[Le bruit en Allemagne : cadre légal et culturel]
En Allemagne, la protection contre le bruit (\all{der Lärmschutz}) est strictement encadrée par la \all{Bundesimmissionsschutzgesetz} (loi fédérale sur la protection contre les immissions) et les \all{Landesimmissionsschutzgesetze} des Länder. Les \textbf{Ruhezeiten} (heures de calme) sont sacralisées :
\begin{itemize}
\item \textbf{Nachtruhe} : 22 h -- 6 h (interdiction de tout bruit anormal).
\item \textbf{Mittagsruhe} : 12 h -- 15 h (selon les Länder et communes, souvent facultative).
\item \textbf{Sonn- und Feiertagsruhe} : toute la journée le dimanche et jours fériés.
\end{itemize}
Tondre la pelouse, faire du bricolage bruyant, organiser une fête bruyante en dehors de ces créneaux expose à une amende (\all{Bußgeld}) et à l'intervention de l'\all{Ordnungsamt} (police municipale). Les chantiers (\all{Baustellen}) ont des horaires stricts (généralement 7 h -- 20 h en semaine, interdits le dimanche). En Suisse, la \all{Ruhezeit} est encore plus stricte (souvent 22 h -- 7 h, pas de lessive le dimanche dans certains immeubles !). Au Cameroun, la loi-cadre sur l'environnement (1996) et les arrêtés municipaux prévoient des normes, mais l'application reste inégale : bruits des bendskins, sono des bars/églises, vendeurs ambulants, chantiers nocturnes. La sensibilisation (\all{Bewusstseinsbildung}) progresse via les associations de quartier et les médias.
\end{savaistu}

\begin{textebox}[Texte support de référence — Lärm in der Stadt (rappel)]
\all{Ich wohne in Yaoundé an einer Hauptstraße, wo der Verkehr schon früh morgens sehr laut ist. Seit drei Wochen gibt es direkt vor unserem Haus eine große Baustelle. Die Baustelle fängt jeden Morgen um sechs Uhr an. Dann hört man den Bagger, die Lastwagen und das Hämmern. Der Lärm ist so laut, dass er durch die geschlossenen Fenster dringt.}

\all{Ich schlafe schlecht, weil der Lärm mich jede Nacht weckt. Gestern Abend hatte ich so starke Kopfschmerzen, dass ich nicht für das Abitur lernen konnte. Ich habe die Fenster geschlossen und Ohrstöpsel getragen, aber das half nicht viel. Die Nachbarn beschweren sich auch und fordern Ruhezeiten. Sie haben einen Brief an die Stadtverwaltung geschrieben. Bisher hat die Stadtverwaltung noch nicht reagiert. Ich hoffe, dass die Baustelle bald fertig ist und wir wieder Ruhe haben.}

\textbf{Glossaire des mots difficiles :}
\begin{itemize}
\item \all{die Hauptstraße, -n} — avenue, grande rue
\item \all{die Baustelle, -n} — chantier
\item \all{der Bagger, -} — pelleteuse, excavatrice
\item \all{der Lastwagen, -} (Pl. \all{die Lastwagen}) — camion poids lourd
\item \all{hämmern} — marteler, faire du bruit de marteau
\item \all{dringen} (v. fort : drang, gedrungen) — pénétrer, s'infiltrer
\item \all{der Ohrstöpsel, -} (souvent Pl. \all{die Ohrstöpsel}) — bouchon d'oreille
\item \all{sich beschweren über + Akk.} — se plaindre de
\item \all{fordern} — exiger, réclamer
\item \all{die Ruhezeit, -en} — heure de calme, période de silence réglementaire
\item \all{die Stadtverwaltung} — mairie, administration municipale
\item \all{reagieren auf + Akk.} — réagir à
\end{itemize}
\end{textebox}

\coursec{Lexique thématique : der Lärm und die Ruhe}

Après avoir travaillé la compréhension globale et détaillée du texte \emph{Lärm in der Stadt}, nous allons maintenant fixer et étendre le vocabulaire essentiel pour parler des bruits, des nuisances sonores et du calme. En allemand, la distinction entre le \emph{bruit gênant} et le \emph{simple son} est lexicale (deux mots différents), ce qui n'est pas le cas en français. Maîtriser ces nuances, les familles de mots et les verbes associés vous permettra de comprendre des articles de presse, de rédiger une lettre de plainte ou de participer à un débat sur la qualité de vie en ville — à Yaoundé comme à Berlin.

\courssub{Le champ lexical du bruit : der Lärm und seine Quellen}

Commençons par les noms désignant les bruits et leurs sources. En allemand, tout nom s'apprend \textbf{impérativement avec son article (der/die/das) et son pluriel}. Observez le tableau ci-dessous : les mots sont classés par catégorie sémantique.

\begin{center}
\begin{tabularx}{\linewidth}{|X|X|X|}
\hline
\rowcolor{popblueL} \textbf{Allemand (Art. + Pluriel)} & \textbf{Français} & \textbf{Remarque / Contexte} \\
\hline
der Lärm (kein Pl.) & le bruit (gênant, nuisance) & \textbf{Mot-clé du thème.} Sens négatif : bruit qui dérange, pollution sonore. Pas de pluriel. \\
das Geräusch, die Geräusche & le bruit, le son (neutre) & Son perçu objectivement, sans jugement de valeur (ex. \all{das Geräusch des Regens} — le bruit de la pluie). \\
der Krach (kein Pl.) & le vacarme, le raffut (fam.) & Familier. Souvent pour des disputes ou bruits violents : \all{Was für ein Krach!} (Quel vacarme !). \\
der Verkehr (kein Pl.) & la circulation, le trafic & Source majeure de \all{Lärm} en ville. \all{der Straßenverkehr} (circulation routière). \\
die Baustelle, die Baustellen & le chantier & Source de bruit temporaire mais intense. \all{Die Baustelle vor dem Haus}. \\
die Musik (kein Pl.) & la musique & Peut être \all{Lärm} pour les voisins : \all{laute Musik}. \\
der Wecker, die Wecker & le réveil & Source de bruit matinal. \all{Der Wecker klingelt}. \\
die Maschine, die Maschinen & la machine & Usine, outils, appareils ménagers. \all{der Maschinenlärm}. \\
\hline
\end{tabularx}
\end{center}

\begin{propriete}[Familles de mots : laut / stören]
L'allemand construit beaucoup de mots par dérivation (préfixes, suffixes). Apprendre la racine permet de deviner le sens de mots inconnus.

\textbf{1. La racine \cle{laut} (bruyant, fort) :}
\begin{itemize}
    \item \all{laut} (adj.) : bruyant, fort $\rightarrow$ \all{laute Musik, ein lautes Geräusch}.
    \item \all{die Lautstärke} (f., die Lautstärken) : le volume sonore. \all{Die Lautstärke reduzieren} (baisser le volume).
    \item \all{läuten} (verbe faible) : sonner (cloche, téléphone, réveil). \all{Das Telefon läutet}.
    \item \all{lärmen} (verbe faible) : faire du bruit, faire du vacarme (souvent péjoratif). \all{Die Kinder lärmten im Garten}.
\end{itemize}

\textbf{2. La racine \cle{stören} (déranger) :}
\begin{itemize}
    \item \all{stören} (verbe faible, transitif) : déranger, gêner. \all{Der Lärm stört mich.} (Le bruit me dérange.)
    \item \all{die Störung, die Störungen} : le trouble, la perturbation, la panne. \all{eine Störung im Netz} (une panne réseau).
    \item \all{störend} (adj.) : gênant, importun. \all{ein störendes Geräusch}.
    \item \all{ungestört} (adj.) : sans être dérangé, tranquille. \all{ungestört schlafen}.
\end{itemize}
\end{propriete}

\begin{exemplebox}[Exemples contextualisés — Bruit]
\all{Der Verkehrslärm in Yaoundé ist sehr stark.} — La pollution sonore due à la circulation à Yaoundé est très forte. \\
\all{Auf der Baustelle nebenan arbeiten laute Maschinen.} — Sur le chantier d'à côté, des machines bruyantes travaillent. \\
\all{Mein Nachbar hört oft laute Musik am Wochenende.} — Mon voisin écoute souvent de la musique forte le week-end. \\
\all{Der Wecker läutet um 5:30 Uhr. Das stört meinen Schlaf.} — Le réveil sonne à 5h30. Cela perturbe mon sommeil. \\
\all{Bitte lärmt nicht so! Das Baby schläft.} — Ne faites pas tant de vacarme, s'il vous plaît ! Le bébé dort.
\end{exemplebox}

\courssub{Le champ lexical du calme : die Ruhe und ihre Qualitäten}

Le calme n'est pas seulement l'absence de bruit ; l'allemand dispose d'un vocabulaire riche pour le décrire (sérénité, silence, détente). Voici les termes clés et leurs contraires directs.

\begin{center}
\begin{tabularx}{\linewidth}{|X|X|X|}
\hline
\rowcolor{popgreenL} \textbf{Allemand (Art. + Pluriel / Forme)} & \textbf{Français} & \textbf{Contraire / Antonyme} \\
\hline
die Ruhe (kein Pl.) & le calme, la tranquillité, le repos & der Lärm, die Unruhe \\
die Stille (kein Pl.) & le silence (absence totale de son) & der Krach, das Geräusch \\
\hline
\all{leise} (adj.) & bas (son), silencieux, tout bas & \all{laut} \\
\all{ruhig} (adj.) & calme, tranquille & \all{unruhig, laut} \\
\all{still} (adj.) & silencieux, muet & \all{laut, gesprächig} \\
\all{entspannend} (adj.) & relaxant, apaisant & \all{stressig, anstrengend} \\
\hline
\end{tabularx}
\end{center}

\begin{propriete}[Famille de mots : ruhen / Ruhe]
La racine \cle{ruhen} (reposer, être en repos) est très productive :
\begin{itemize}
    \item \all{ruhen} (verbe faible, intransitif) : reposer, être au repos. \all{Die Maschinen ruhen am Sonntag.} (Les machines sont au repos le dimanche.)
    \item \all{sich ausruhen} (verbe réfléchi, faible, à particule séparable) : se reposer, se détendre. \all{Ich ruhe mich eine Stunde aus.} (Je me repose une heure.)
    \item \all{die Ruhe} (f., pas de pluriel) : le calme, la tranquillité. \all{Ruhe bewahren!} (Garder son calme !)
    \item \all{ruhig} (adj.) : calme, tranquille. \all{ein ruhiges Zimmer} (une chambre calme).
    \item \all{unruhig} (adj.) : agité, inquiet, sans repos. \all{ein unruhiger Schlaf} (un sommeil agité).
    \item \all{beruhigen} (verbe faible) : calmer, apaiser. \all{Beruhige dich!} (Calme-toi !)
\end{itemize}
\end{propriete}

\begin{exemplebox}[Exemples contextualisés — Calme]
\all{Nachts ist es in meinem Viertel endlich ruhig.} — La nuit, dans mon quartier, c'est enfin calme. \\
\all{Ich brauche absolute Stille zum Lernen.} — J'ai besoin de silence absolu pour étudier. \\
\all{Der Park ist ein Ort der Ruhe und Entspannung.} — Le parc est un lieu de calme et de détente. \\
\all{Mach bitte das Fenster zu, damit es leiser wird.} — Ferme la fenêtre s'il te plaît, pour que ce soit moins bruyant. \\
\all{Nach dem stressigen Tag brauche ich Ruhe.} — Après cette journée stressante, j'ai besoin de calme.
\end{exemplebox}

\courssub{Verbes clés et expressions utiles pour réagir}

Pour interagir face au bruit (se plaindre, demander le silence, décrire son état), quelques verbes et tournures figées sont indispensables. Notez la conjugaison du verbe irrégulier \all{schlafen} au présent et au prétérit (revu en section 6).

\begin{center}
\begin{tabularx}{\linewidth}{|X|X|X|}
\hline
\rowcolor{poporangeL} \textbf{Verbe / Expression} & \textbf{Signification / Structure} & \textbf{Exemple allemand + traduction} \\
\hline
\all{stören} (faible) & déranger qqn (Akk.) & \all{Der Lärm stört mich beim Lesen.} (Le bruit me dérange quand je lis.) \\
\all{schlafen} (irr. : schläft, schlief, hat geschlafen) & dormir & \all{Ich kann vor Lärm nicht schlafen.} (Je ne peux pas dormir à cause du bruit.) \\
\all{sich ausruhen} (réfl., faible, séparable) & se reposer & \all{Er ruht sich auf dem Balkon aus.} (Il se repose sur le balcon.) \\
\all{sich beschweren über + Akk.} & se plaindre de & \all{Die Nachbarn beschweren sich über den Krach.} (Les voisins se plaignent du vacarme.) \\
\all{lärmen} (faible) & faire du bruit / vacarme & \all{Die Jugendlichen lärmten bis Mitternacht.} (Les jeunes ont fait du bruit jusqu'à minuit.) \\
\hline
\end{tabularx}
\end{center}

\begin{definition}[Verbes à retenir absolument]
\all{stören} = déranger, gêner (verbe faible, transitif direct). Construction : \textbf{Etwas/jemand stört jemanden (Akkusativ)}. \\
Ex. \all{Die Musik stört mich.} (La musique me dérange.) — \all{Stört dich der Lärm?} (Le bruit te dérange ?)

\all{sich beschweren über + Akkusativ} = se plaindre de. Verbe réfléchi + préposition \textbf{über} gouvernant l'accusatif. \\
Ex. \all{Er beschwert sich über den Lärm.} (Il se plaint du bruit.) — \all{Wir beschweren uns über die laute Musik.} (Nous nous plaignons de la musique forte.)
\end{definition}

\begin{attention}[Piège majeur : \cle{laut} — Adjectif ou Préposition ?]
C'est l'erreur n°1 des francophones ! Le mot \all{laut} a deux fonctions grammaticales totalement différentes :
\begin{enumerate}
    \item \textbf{Adjectif} (très fréquent) : « bruyant, fort ». S'accorde en genre/nombre/cas. \\
    \all{laute Musik} (musique forte), \all{ein lautes Auto} (une voiture bruyante), \all{Sei leise, nicht laut!} (Sois silencieux, pas bruyant !).
    \item \textbf{Préposition} (formel, écrit) : « selon, d'après, suivant ». Gouverne normalement le \textbf{GÉNITIF} (le datif est admis à l'oral et dans la langue courante). Ne s'accorde PAS. \\
    \all{Laut Wetterbericht regnet es morgen.} (Selon la météo, il pleut demain.) \\
    \all{Laut Vertrag muss er zahlen.} (Selon le contrat, il doit payer.)
\end{enumerate}
\textbf{Astuce :} Si \all{laut} est devant un nom sans article et signifie « selon » $\rightarrow$ Préposition + Génitif. Si \all{laut} est devant un nom avec article et qualifie ce nom $\rightarrow$ Adjectif.
\end{attention}

\begin{savaistu}[Lärm vs. Geräusch : une distinction culturelle]
En français, le mot « bruit » couvre les deux réalités. En allemand, la distinction est nette et culturellement ancrée :
\begin{itemize}
    \item \cle{der Lärm} = bruit \emph{subjectif}, \emph{gênant}, \emph{indésirable}, souvent nocif pour la santé (\all{die Lärmbelästigung} = nuisance sonore, \all{der Lärmschutz} = protection contre le bruit). C'est un problème de société et de droit (\all{die Ruhezeiten} = heures de calme légales).
    \item \cle{das Geräusch} = son, bruit \emph{objectif}, neutre, physique. \all{Ein seltsames Geräusch} (un bruit étrange), \all{Geräusche im Motor} (bruits dans le moteur).
\end{itemize}
\textbf{Conseil de traduction :} Si vous pouvez remplacer « bruit » par « vacarme », « nuisance », « tapage » $\rightarrow$ \all{Lärm}. Si c'est un son neutre, technique, naturel $\rightarrow$ \all{Geräusch}. \emph{Ex. : « Le bruit de la mer » = \all{das Geräusch des Meeres} (agréable). « Le bruit des voisins » = \all{der Lärm der Nachbarn} (gênant).}
\end{savaistu}

\courssub{Entraînement lexical : Carte mentale « Lärm \& Ruhe »}

La figure ci-dessous résume visuellement les trois piliers du vocabulaire de cette leçon. Utilisez-la pour réviser ou préparer une prise de parole.

\begin{popfigure}
\begin{tikzpicture}[
    mindmap,
    grow cyclic,
    every node/.style={concept, execute at begin node=\hskip0pt},
    root concept/.append style={concept color=popblue, minimum size=2.5cm, text=white, font=\bfseries\large},
    level 1 concept/.append style={level distance=4.5cm, sibling angle=120, minimum size=1.8cm, font=\bfseries\normalsize},
    level 2 concept/.append style={level distance=3.2cm, sibling angle=45, minimum size=1.4cm, font=\small},
    concept color=popblue!80,
    text=white
]
\node[root concept] {Lärm \& Ruhe}
    child[concept color=poporange] {
        node[align=center]{Geräusche \\ (Sources)}
        child { node {Verkehr} }
        child { node {Baustelle} }
        child { node {Musik} }
        child { node {Wecker} }
        child { node {Maschine} }
        child { node {Nachbarn} }
    }
    child[concept color=popgreen] {
        node[align=center]{Adjektive \\ (Qualités)}
        child { node {laut / leise} }
        child { node {ruhig / unruhig} }
        child { node {still} }
        child { node {stressig} }
        child { node {entspannend} }
        child { node {störend} }
    }
    child[concept color=poppurple] {
        node[align=center]{Verben \\ (Actions)}
        child { node {stören} }
        child { node {schlafen} }
        child { node {sich ausruhen} }
        child { node {sich beschweren} }
        child { node {lärmen} }
        child { node {beruhigen} }
    };
\end{tikzpicture}
\legende{Carte mentale du vocabulaire : sources de bruit (orange), adjectifs de description (vert), verbes d'action (violet).}
\end{popfigure}

\courssub{Texte d'entraînement : « Ruhestörung in Yaoundé »}

Le texte suivant réemploie le vocabulaire vu ci-dessus dans un contexte camerounais authentique. Lisez-le à haute voix, soulignez les mots du lexique cible, puis répondez aux questions.

\begin{textebox}[Ruhestörung in Yaoundé — Ein Leserbrief]
Liebe Redaktion,

ich schreibe Ihnen, weil die Lärmbelästigung in unserem Viertel in Yaoundé unerträglich geworden ist. Tagsüber donnern die schweren Lastwagen und die alten Taxis auf der Hauptstraße vorbei. Der Verkehrslärm ist so laut, dass man sein eigenes Wort nicht mehr versteht. Dazu kommen die Baustellen: Seit drei Monaten baut man ein neues Geschäftshaus direkt neben unserem Wohnhaus. Die Maschinen — Bagger, Betonmischer und Presslufthämmer — lärmen von 7 Uhr morgens bis 19 Uhr abends. Selbst sonntags ist keine Ruhe.

Nachts ist die Situation nicht besser. Viele junge Leute hören extrem laute Musik aus ihren Autos oder aus den Bars an der Ecke. Der Bass dröhnt bis in unser Schlafzimmer. Meine kleine Tochter wacht oft auf und weint, weil der Krach sie erschreckt. Mein Mann und ich können nicht schlafen und sind am nächsten Tag total unruhig und müde. Wir haben uns schon bei der Polizei beschwert, aber nichts hat sich geändert.

Ich appelliere an die Stadtverwaltung: Bitte sorgen Sie für die Einhaltung der Ruhezeiten! Wir brauchen Orte der Stille und Erholung. Lärm macht krank — das ist bewiesen. Eine ruhige Nacht ist ein Menschenrecht, kein Luxus.

Mit freundlichen Grüßen \\
Marie N., Yaoundé-Bastos
\end{textebox}

\begin{center}
\begin{tabularx}{\linewidth}{|X|X|}
\hline
\rowcolor{popblueL} \textbf{Allemand (Glossaire)} & \textbf{Français} \\
\hline
die Lärmbelästigung (kein Pl.) & la nuisance sonore, la pollution sonore \\
unerträglich & insupportable \\
donnern (faible) & gronder, faire un bruit sourd (moteurs lourds) \\
das Geschäftshaus, die Geschäftshäuser & l'immeuble de bureaux, le bâtiment commercial \\
der Bagger, die Bagger & la pelleteuse, l'excavatrice \\
der Betonmischer, die Betonmischer & la bétonnière, le malaxeur \\
der Presslufthammer, die Presslufthämmer & le marteau-piqueur \\
dröhnen (faible) & résonner, bourdonner (basse fréquence) \\
das Schlafzimmer, die Schlafzimmer & la chambre à coucher \\
aufwachen (wacht auf, wachte auf, ist aufgewacht) & se réveiller \\
erschrecken (erschrickt, erschrak, ist erschrocken) & être effrayé, avoir peur \\
sich beschweren bei + Dat. & se plaindre auprès de \\
die Stadtverwaltung, die Stadtverwaltungen & l'administration municipale, la mairie \\
die Einhaltung (kein Pl.) & le respect, l'observance (d'une règle) \\
die Ruhezeit, die Ruhezeiten & l'heure de calme (légale), la plage de silence \\
die Erholung (kein Pl.) & la récupération, le repos, la détente \\
bewiesen & prouvé \\
das Menschenrecht, die Menschenrechte & le droit de l'homme \\
\hline
\end{tabularx}
\end{center}

\begin{experience}[Compréhension et réemploi lexical]
Répondez en allemand (phrases complètes) :
\begin{enumerate}
    \item Nennen Sie drei Quellen für Lärm, die Marie im Text nennt. (Citez trois sources de bruit mentionnées par Marie.)
    \item Welche Adjektive benutzt sie, um die Musik zu beschreiben? (Quels adjectifs utilise-t-elle pour décrire la musique ?)
    \item Was bedeutet „Lärm macht krank“? Nennen Sie zwei Folgen für Marie und ihre Familie. (Que signifie « Le bruit rend malade » ? Citez deux conséquences pour Marie et sa famille.)
    \item Finden Sie im Text ein Verb aus der Familie „stören“ und eines aus der Familie „Ruhe“. (Trouvez dans le texte un verbe de la famille « stören » et un de la famille « Ruhe ».)
\end{enumerate}
\end{experience}

\begin{corrige}[Compréhension et réemploi lexical]
\begin{enumerate}
    \item \all{Marie nennt den Verkehrslärm (Lastwagen und Taxis), die Baustelle (Bagger, Betonmischer, Presslufthämmer) und die laute Musik aus Autos und Bars.} — Marie cite le bruit de la circulation (camions et taxis), le chantier (pelleteuses, bétonnières, marteaux-piqueurs) et la musique forte des voitures et des bars.
    \item \all{Sie beschreibt die Musik als „extrem laut“.} — Elle décrit la musique comme « extrêmement forte ».
    \item \all{„Lärm macht krank“ bedeutet: Lärm ist schlecht für die Gesundheit. Folgen: Marie und ihr Mann können nicht schlafen und sind müde und unruhig; die kleine Tochter wacht auf und weint.} — « Le bruit rend malade » signifie : le bruit est mauvais pour la santé. Conséquences : Marie et son mari ne peuvent pas dormir et sont fatigués et agités ; la petite fille se réveille et pleure.
    \item \all{Familie „stören“: die Ruhestörung (im Titel) / beschweren (sich beschweren über etwas, das stört). Familie „Ruhe“: die Ruhezeiten, die Erholung, ruhig.} — Famille « stören » : \all{die Ruhestörung} (dans le titre). Famille « Ruhe » : \all{die Ruhezeiten}, \all{die Erholung}, \all{ruhig}.
\end{enumerate}
\end{corrige}

\courssub{Exercice résolu : Mots manquants et traduction}

\begin{exoresolu}[Compléter et traduire]
Complétez les phrases allemandes avec le mot approprié du lexique (nom, adjectif ou verbe conjugué), puis traduisez la phrase en français.

\textbf{1.} Die Nachbarn hören \_\_\_\_\_\_\_\_ Musik. (bruyante / forte) \\
\textbf{2.} Nachts brauche ich absolute \_\_\_\_\_\_\_\_ zum Schlafen. (silence) \\
\textbf{3.} Der \_\_\_\_\_\_\_\_ von der Baustelle ist sehr laut. (bruit de chantier) \\
\textbf{4.} Bitte \_\_\_\_\_\_\_\_ Sie nicht so! Das Baby schläft. (faire du vacarme — verbe) \\
\textbf{5.} Ich \_\_\_\_\_\_\_\_ mich bei der Polizei \_\_\_\_\_\_\_\_ den Krach. (me plains / du vacarme — verbe + préposition) \\
\textbf{6.} Die \_\_\_\_\_\_\_\_ im Park ist sehr entspannend. (calme) \\
\textbf{7.} \_\_\_\_\_\_\_\_ Arzt ist Lärm schlecht für das Herz. (Selon / D'après — préposition)

\tcblower
\textbf{Solution détaillée :}
\begin{enumerate}
    \item \textbf{laute} (adjectif, accusatif féminin devant \all{Musik}) — \textit{Les voisins écoutent de la musique \textbf{forte}.}
    \item \textbf{Stille} (ou \textbf{Ruhe}) — \textit{La nuit, j'ai besoin de \textbf{silence} (ou \textbf{calme}) absolu pour dormir.}
    \item \textbf{Lärm} (\all{der Lärm von der Baustelle}) ou mot composé \textbf{Baustellenlärm} — \textit{Le \textbf{bruit du chantier} est très fort.}
    \item \textbf{lärmen} (impératif poli, 2\up{e} pers. pl. : \all{lärmen Sie}) — \textit{S'il vous plaît, \textbf{ne faites pas tant de vacarme} ! Le bébé dort.}
    \item \textbf{beschwere ... über} (\all{ich beschwere mich über + Akk.}) — \textit{Je \textbf{me plains auprès de la police du vacarme}.}
    \item \textbf{Ruhe} (ou \textbf{Stille}) — \textit{Le \textbf{calme} (ou \textbf{silence}) dans le parc est très relaxant.}
    \item \textbf{Laut} (préposition + génitif : \all{Laut Arzt...}) — \textit{\textbf{Selon le médecin}, le bruit est mauvais pour le cœur.}
\end{enumerate}
\end{exoresolu}

\begin{methode}[Comment enrichir son lexique « Environnement \& Bruit » durablement]
\begin{enumerate}
    \item \textbf{Apprenez par « couples » (Wortpaare) :} \all{laut – leise}, \all{ruhig – unruhig}, \all{Lärm – Ruhe}, \all{stören – nicht stören}. Le cerveau retient mieux les contrastes.
    \item \textbf{Construisez votre « Dictionnaire personnel » :} une page par racine (\all{laut}, \all{stör}, \all{ruh}). Notez \textbf{tous} les dérivés rencontrés (noms, verbes, adjectifs, prépositions) avec un exemple de phrase personnelle (ex. \all{Mein Zimmer ist ruhig.}).
    \item \textbf{Écoutez activement :} regardez un reportage (Deutsche Welle, ZDF) sur \all{Lärmschutz} ou \all{Verkehrslärm}. Notez 5 mots nouveaux par vidéo.
    \item \textbf{Produisez :} écrivez une courte plainte formelle (\all{Beschwerdebrief}) à votre mairie (Yaoundé, Douala, Buea...) en réemployant \all{sich beschweren über}, \all{die Lärmbelästigung}, \all{die Ruhezeiten}, \all{unerträglich}.
\end{enumerate}
\end{methode}

\coursec{Grammaire 1 : les subordonnées temporelles (wenn, als)}

Dans la section précédente, nous avons découvert le lexique du bruit et du calme : \all{der Lärm}, \all{laut}, \all{leise}, \all{die Nachbarn}, \all{stören}, \all{die Ruhe}. Nous allons maintenant observer comment ces mots s'organisent dans des phrases plus complexes. En effet, pour raconter une nuisance sonore, il faut souvent dire \emph{quand} elle se produit : « quand les machines travaillent », « quand j'ai emménagé ». L'allemand possède deux mots pour dire « quand » : \all{als} et \all{wenn}. C'est l'une des distinctions les plus importantes — et les plus rentables à l'examen — de la grammaire allemande.

\courssub{Observation : deux exemples tirés du texte}

Reprenons deux phrases de notre texte support \emph{Lärm in der Stadt} :

\begin{exemplebox}[Deux phrases à observer]
\all{Als ich nach München zog, war die Straße ruhig.} « Quand j'ai emménagé à Munich, la rue était calme. »\par
\all{Wenn die Maschinen arbeiten, kann ich nicht schlafen.} « Quand les machines fonctionnent, je ne peux pas dormir. »
\end{exemplebox}

Observons attentivement ces deux phrases et posons-nous trois questions :

\begin{enumerate}
\item \textbf{Où est le verbe conjugué dans la subordonnée ?} Dans \all{Als ich nach München \textbf{zog}} et \all{Wenn die Maschinen \textbf{arbeiten}}, le verbe conjugué se trouve \textbf{à la fin} de la subordonnée.
\item \textbf{Que se passe-t-il après la virgule ?} Le verbe de la principale vient \textbf{immédiatement} après la virgule : \all{..., \textbf{war} die Straße ruhig.} et \all{..., \textbf{kann} ich nicht schlafen.} Le sujet vient après le verbe : c'est l'\textbf{inversion}.
\item \textbf{Pourquoi \all{als} dans la première phrase et \all{wenn} dans la seconde ?} La première phrase décrit un événement \textbf{unique et passé} (un seul déménagement). La seconde décrit une situation \textbf{répétée, générale} (chaque fois que les machines tournent). C'est exactement la clé de la distinction.
\end{enumerate}

\courssub{La règle de construction : verbe à la fin et inversion}

\begin{propriete}[Construction de la subordonnée temporelle]
\begin{enumerate}
\item \all{als} et \all{wenn} sont des \textbf{conjonctions de subordination} : le verbe conjugué de la subordonnée est rejeté \textbf{à la fin} de celle-ci.
\item La subordonnée et la principale sont \textbf{toujours séparées par une virgule}.
\item Si la subordonnée est placée \textbf{en tête de phrase}, elle occupe la position 1 : le verbe conjugué de la principale vient alors \textbf{immédiatement après la virgule} (position 2), et le sujet suit le verbe (\textbf{inversion}).
\item Si la subordonnée est placée \textbf{après} la principale, la principale garde son ordre normal : \all{Ich kann nicht schlafen, wenn die Maschinen arbeiten.}
\end{enumerate}
\end{propriete}

\begin{popfigure}
\begin{tikzpicture}[font=\small, >=Stealth]
\node[draw=popblue, fill=popblueL, rounded corners, align=center, minimum width=6.2cm, minimum height=1.3cm] (sub) at (0,0) {\textbf{Subordonnée}\\ \all{Als ich nach München \textbf{zog}}};
\node[draw=poporange, fill=poporangeL, rounded corners, align=center, minimum width=6.2cm, minimum height=1.3cm] (main) at (8.4,0) {\textbf{Principale (inversion)}\\ \all{\textbf{war} die Straße ruhig}};
\draw[->, very thick, popdark] (sub) -- (main) node[midway, above, align=center]{virgule\\ obligatoire};
\draw[->, very thick, poppurple, bend left=25] (3.2,-0.75) to node[below, align=center]{verbe conjugué\\ en position 2} (5.4,-0.75);
\node[align=center, popgreen!70!black] at (4.2,1.6) {La subordonnée = position 1 ; le verbe de la principale suit tout de suite.};
\end{tikzpicture}
\legende{Structure de la phrase : subordonnée en tête, puis inversion dans la principale.}
\end{popfigure}

\begin{exemplebox}[Même phrase, deux constructions]
\all{Wenn es laut ist, kann ich nicht arbeiten.} « Quand il y a du bruit, je ne peux pas travailler. »\par
\all{Ich kann nicht arbeiten, wenn es laut ist.} « Je ne peux pas travailler quand il y a du bruit. »\par
\medskip
Dans le premier cas, la subordonnée ouvre la phrase : inversion (\all{kann ich}). Dans le second, la principale garde l'ordre sujet + verbe (\all{Ich kann}).
\end{exemplebox}

\courssub{Emploi de als : l'événement unique au passé}

\begin{propriete}[Emploi de als]
\all{als} s'emploie pour un événement \textbf{passé, accompli et unique} (une seule fois, ou une période passée considérée comme un tout). Il se combine toujours avec un \textbf{temps passé} : Präteritum ou Perfekt. Il ne s'emploie \textbf{jamais} au présent ni au futur.\par
\all{Als ich klein war, wohnten wir in Douala.} « Quand j'étais petit, nous habitions à Douala. » (une période unique de ma vie)\par
\all{Als der Zug ankam, wurde es still.} « Quand le train est arrivé, le silence est revenu. » (une seule fois)
\end{propriete}

Attention : une période de vie (\all{als ich klein war}, \all{als ich in Yaoundé wohnte}) compte comme \textbf{un seul bloc} : on utilise donc \all{als}. En revanche, ce qui se \textbf{répétait} à l'intérieur de cette période exige \all{wenn} (voir plus bas).

\courssub{Emploi de wenn : présent, futur, habitude, condition}

\begin{propriete}[Emploi de wenn]
\all{wenn} s'emploie dans tous les autres cas :
\begin{itemize}
\item \textbf{Présent} (généralité, fait actuel) : \all{Wenn die Nachbarn Musik hören, ist es laut.} « Quand les voisins écoutent de la musique, il y a du bruit. »
\item \textbf{Futur} : \all{Wenn ich nach Deutschland reise, besuche ich München.} « Quand je voyagerai en Allemagne, je visiterai Munich. »
\item \textbf{Habitude au présent} : \all{Wenn ich müde bin, gehe ich früh ins Bett.} « Quand je suis fatigué, je me couche tôt. »
\item \textbf{Habitude au passé} (répétition) : \all{Wenn ich früher nach Hause kam, kochte meine Mutter.} « Quand je rentrais à la maison, ma mère cuisinait. »
\item \textbf{Condition} (si) : \all{Wenn es ruhig ist, kann ich gut lernen.} « Si c'est calme, je peux bien travailler. »
\end{itemize}
\end{propriete}

\begin{attention}[Le cas particulier qui piège tous les francophones]
« Quand j'étais enfant, je jouais dehors » décrit une \textbf{habitude passée} (une action répétée). En allemand, on dit donc :\par
\all{Wenn ich ein Kind war, spielte ich draußen.} — avec \all{wenn}, et non \all{als} !\par
Retenez : \all{als} = \textbf{une seule fois} dans le passé. Dès qu'il y a \textbf{répétition} (chaque fois, toujours, souvent, d'habitude), même au passé, c'est \all{wenn}. En allemand soutenu, on peut renforcer l'idée de répétition avec \all{immer wenn} : \all{Immer wenn er kam, war es laut.} « Chaque fois qu'il venait, c'était bruyant. »
\end{attention}

\begin{attention}[Ne pas confondre wenn et wann]
\all{wann} (« quand ? ») sert uniquement à poser une \textbf{question}, directe ou indirecte : \all{Wann kommst du?} « Quand viens-tu ? » ; \all{Ich weiß nicht, wann er kommt.} « Je ne sais pas quand il vient. »\par
Dans une subordonnée temporelle qui raconte un fait (et non une question), on utilise toujours \all{als} ou \all{wenn}, jamais \all{wann}. Prononciation : \all{wann} se dit « vanne », \all{wenn} se dit « vène ».
\end{attention}

\courssub{Tableau récapitulatif et comparaison avec le français}

\begin{center}
\begin{tabularx}{\linewidth}{|X|X|X|}
\hline
\rowcolor{popblueL} Situation & Conjonction & Exemple \\
\hline
Passé, une seule fois (événement unique) & \all{als} & \all{Als ich ankam, war es still.} \\
\hline
Passé, habitude (répétition) & \all{wenn} & \all{Wenn ich kam, war es immer laut.} \\
\hline
Présent (généralité, habitude) & \all{wenn} & \all{Wenn es laut ist, kann ich nicht schlafen.} \\
\hline
Futur & \all{wenn} & \all{Wenn ich Zeit habe, rufe ich dich an.} \\
\hline
Condition (« si ») & \all{wenn} & \all{Wenn du leise bist, störst du niemanden.} \\
\hline
\end{tabularx}
\end{center}

\begin{center}
\begin{tabularx}{\linewidth}{|X|X|X|}
\hline
\rowcolor{popblueL} Français & Deutsch & Remarque \\
\hline
Quand j'ai emménagé, la rue était calme. & \all{Als ich einzog, war die Straße ruhig.} & une seule fois, passé : \all{als} \\
\hline
Quand je rentrais, ma mère cuisinait. & \all{Wenn ich nach Hause kam, kochte meine Mutter.} & habitude passée : \all{wenn} \\
\hline
Quand il y a du bruit, je ne travaille pas. & \all{Wenn es laut ist, arbeite ich nicht.} & généralité : \all{wenn} \\
\hline
Quand tu viendras, nous irons au marché. & \all{Wenn du kommst, gehen wir auf den Markt.} & futur : \all{wenn} + présent \\
\hline
\end{tabularx}
\end{center}

En français, un seul mot — « quand » — couvre toutes les situations. En allemand, c'est le \textbf{contexte} qui décide : le français « quand » se traduit tantôt par \all{als}, tantôt par \all{wenn}. Il faut donc toujours se poser deux questions : \emph{passé ou non ?} puis \emph{une seule fois ou répétition ?}

\begin{popfigure}
\begin{tikzpicture}[font=\small, >=Stealth, node distance=1.1cm]
\node[draw=popdark, fill=popblueL, rounded corners, align=center] (q1) {L'événement est-il\\ \textbf{au passé} ?};
\node[draw=popdark, fill=poporangeL, rounded corners, align=center, below left=1.1cm and 0.2cm of q1] (q2) {Une \textbf{seule} fois ?};
\node[draw=popgreen!70!black, fill=popgreenL, rounded corners, align=center, below right=1.1cm and 0.2cm of q1] (w1) {\all{wenn}\\ présent / futur / condition};
\node[draw=popgreen!70!black, fill=popgreenL, rounded corners, align=center, below left=1.1cm and -0.4cm of q2] (a1) {\all{als}\\ événement unique};
\node[draw=popgreen!70!black, fill=popgreenL, rounded corners, align=center, below right=1.1cm and -0.4cm of q2] (w2) {\all{wenn}\\ habitude au passé};
\draw[->, thick, popdark] (q1) -- node[above left]{oui} (q2);
\draw[->, thick, popdark] (q1) -- node[above right]{non} (w1);
\draw[->, thick, popdark] (q2) -- node[above left]{oui} (a1);
\draw[->, thick, popdark] (q2) -- node[above right]{non} (w2);
\end{tikzpicture}
\legende{Arbre de décision : \all{als} ou \all{wenn} ?}
\end{popfigure}

\courssub{Texte d'application}

\begin{textebox}[Ein lautes Wochenende in Yaoundé]
\all{Als ich letztes Jahr nach Yaoundé zog, freute ich mich auf mein neues Leben. Die Wohnung war schön, und die Nachbarn waren freundlich. Aber als das erste Wochenende kam, bekam ich einen Schock. Am Samstagmorgen begann der Nachbar um sechs Uhr mit seiner lauten Musik. Wenn er Musik hört, hört man sein Radio in der ganzen Straße! Wenn ich früher in meinem Dorf aufwachte, hörte ich nur die Vögel. Hier höre ich Motorräder, Hupen und Generatoren.\par
Als ich ein Kind war, spielte ich oft draußen, und es war immer ruhig. Wenn es heute Abend ruhig ist, werde ich endlich gut schlafen. Mein Freund sagt: „Wenn man in der Stadt lebt, muss man den Lärm akzeptieren.“ Aber ich glaube: Das Leben wäre schöner, wenn alle Nachbarn leiser wären. Als ich das dem Nachbarn sagte, lachte er nur. Vielleicht kaufe ich mir Ohrenstöpsel — wenn ich Geld habe!}
\end{textebox}

\textbf{Glossaire :} \all{sich freuen auf} + Akkusativ (se réjouir de) ; \all{der Schock, -s} (le choc) ; \all{die Hupe, -n} (le klaxon) ; \all{der Generator, -en} (le groupe électrogène) ; \all{der Vogel, die Vögel} (l'oiseau) ; \all{akzeptieren} (accepter) ; \all{der Ohrenstöpsel, -} (la boule Quiès) ; \all{aufwachen} (se réveiller).

\textbf{Questions de grammaire :} 1) Relevez toutes les subordonnées avec \all{als} et justifiez chaque emploi. 2) Relevez toutes les subordonnées avec \all{wenn} et précisez : habitude présente, habitude passée, futur ou condition ? 3) Pourquoi dit-on \all{Als ich ein Kind war} mais \all{Wenn ich früher aufwachte} ? 4) Relevez la phrase au Konjunktiv II : quelle nuance exprime-t-elle ?

\courssub{Exercice résolu}

\begin{exoresolu}[als oder wenn ?]
Complétez avec \all{als} ou \all{wenn}, puis traduisez.\par
1. \all{... ich in Douala ankam, regnete es.}\par
2. \all{... die Maschinen arbeiten, kann ich nicht schlafen.}\par
3. \all{... ich früher müde war, trank ich einen Kaffee.}\par
4. \all{... du morgen kommst, zeige ich dir die Stadt.}\par
5. \all{... mein Vater jung war, wohnte er in Bafoussam.}
\tcblower
\textbf{Solution détaillée :}\par
1. \all{\textbf{Als} ich in Douala ankam, regnete es.} — arrivée unique au passé. « Quand je suis arrivé à Douala, il pleuvait. »\par
2. \all{\textbf{Wenn} die Maschinen arbeiten, kann ich nicht schlafen.} — généralité au présent. « Quand les machines fonctionnent, je ne peux pas dormir. »\par
3. \all{\textbf{Wenn} ich früher müde war, trank ich einen Kaffee.} — habitude passée (\all{früher} = répétition). « Quand j'étais fatigué, je buvais un café. »\par
4. \all{\textbf{Wenn} du morgen kommst, zeige ich dir die Stadt.} — futur. « Quand tu viendras demain, je te montrerai la ville. »\par
5. \all{\textbf{Als} mein Vater jung war, wohnte er in Bafoussam.} — période unique passée (sa jeunesse, un seul bloc). « Quand mon père était jeune, il habitait à Bafoussam. »
\end{exoresolu}

\begin{exercice}[À vous : als oder wenn ?]
Complétez avec \all{als} ou \all{wenn}, puis traduisez les phrases 1 et 4.\par
1. \all{... der Generator lief, konnten wir nicht schlafen.} (chaque nuit, au passé)\par
2. \all{... ich gestern nach Hause kam, war niemand da.}\par
3. \all{... es in der Nacht ruhig ist, schlafe ich sehr gut.}\par
4. \all{... wir früher Fußball spielten, machten wir viel Lärm.}\par
5. \all{... du nach Deutschland reist, besuche ich dich.}\par
6. Remettez la subordonnée en tête : \all{Ich konnte nicht arbeiten, als der Lärm begann.}
\end{exercice}

\begin{corrige}[Exercice : als oder wenn ?]
1. \all{\textbf{Wenn} der Generator lief, konnten wir nicht schlafen.} — habitude passée (chaque nuit). « Quand le groupe électrogène tournait, nous ne pouvions pas dormir. »\par
2. \all{\textbf{Als} ich gestern nach Hause kam, war niemand da.} — événement unique au passé.\par
3. \all{\textbf{Wenn} es in der Nacht ruhig ist, schlafe ich sehr gut.} — généralité au présent.\par
4. \all{\textbf{Wenn} wir früher Fußball spielten, machten wir viel Lärm.} — habitude passée. « Quand nous jouions au football, nous faisions beaucoup de bruit. »\par
5. \all{\textbf{Wenn} du nach Deutschland reist, besuche ich dich.} — futur.\par
6. \all{\textbf{Als der Lärm begann, konnte ich nicht arbeiten.}} — subordonnée en tête, verbe conjugué de la principale (\all{konnte}) immédiatement après la virgule, puis inversion du sujet.
\end{corrige}

\begin{attention}[Les trois erreurs à ne jamais commettre]
\begin{enumerate}
\item \textbf{Oublier la virgule} entre la subordonnée et la principale : elle est \textbf{obligatoire} en allemand.
\item \textbf{Oublier l'inversion} quand la subordonnée ouvre la phrase : \all{Als er kam, \textbf{war ich} müde.} (et non : \emph{ich war müde}). La subordonnée entière compte comme la position 1.
\item \textbf{Employer \all{als} pour une habitude passée} : dès que l'action se répète (souvent, chaque fois, toujours), c'est \all{wenn}, même au Präteritum. Formule magique : \all{als} = une seule fois ; sinon = \all{wenn}.
\end{enumerate}
\end{attention}

\begin{aretenir}[L'essentiel]
\all{als} et \all{wenn} introduisent une subordonnée temporelle : verbe conjugué à la fin, virgule obligatoire, inversion si la subordonnée est en tête. \all{als} = événement passé \textbf{unique} ; \all{wenn} = présent, futur, condition et toute \textbf{répétition}, y compris les habitudes au passé. \all{wann} est réservé aux questions. En cas de doute, appliquez l'arbre de décision : passé ? une seule fois ?
\end{aretenir}

\coursec{Grammaire 2 : les propositions relatives simples}

Dans la section précédente, nous avons appris à relier deux événements dans le temps grâce aux subordonnées temporelles avec \all{wenn} et \all{als}. Nous allons maintenant découvrir un deuxième type de subordonnée, tout aussi fréquent : la \cle{proposition relative}. Elle ne sert pas à situer dans le temps, mais à \textbf{préciser un nom}, à donner une information supplémentaire sur une personne ou une chose. C'est l'outil idéal pour décrire le voisin qui fait du bruit, le chantier qui commence chaque matin, le réveil qui sonne trop tôt...

\courssub{Observation : repérons les relatives dans le texte}

Relisons deux phrases du texte support \all{Lärm in der Stadt} :

\begin{exemplebox}[Deux phrases du texte]
\all{Vor unserem Haus gibt es eine Baustelle, die jeden Morgen um sechs Uhr anfängt.}\\
« Devant notre maison, il y a un chantier qui commence chaque matin à six heures. »

\all{Meine Nachbarn, die sehr laute Musik hören, verstehen das nicht.}\\
« Mes voisins, qui écoutent de la musique très forte, ne comprennent pas cela. »
\end{exemplebox}

Observons attentivement :

\begin{itemize}
\item Dans les deux phrases, il y a un \textbf{nom} (\all{eine Baustelle}, \all{meine Nachbarn}) suivi d'une \textbf{virgule}, puis d'un petit mot (\all{die}), puis d'un groupe de mots dont le \textbf{verbe conjugué se trouve tout à la fin} (\all{anfängt}, \all{hören}).
\item Le petit mot \all{die} renvoie au nom qui le précède : on l'appelle le \cle{pronom relatif}. Le nom auquel il renvoie s'appelle l'\cle{antécédent}.
\item En français, ce petit mot correspond à « qui » : \emph{un chantier qui commence}, \emph{des voisins qui écoutent}.
\end{itemize}

\begin{definition}[La proposition relative]
Une \cle{proposition relative} est une subordonnée qui précise ou décrit un nom (l'\cle{antécédent}). Elle est introduite par un \cle{pronom relatif} (\all{der, die, das}...) placé juste après une virgule, et son verbe conjugué est rejeté à la fin de la subordonnée. Elle répond à la question : « lequel ? laquelle ? ».
\end{definition}

\courssub{La règle : construction de la proposition relative}

\begin{propriete}[Construction de la relative]
La proposition relative allemande obéit à trois règles strictes :
\begin{enumerate}
\item Elle est toujours précédée d'une \textbf{virgule}.
\item Elle commence par un \textbf{pronom relatif} qui prend le \textbf{genre} (masculin, féminin, neutre) et le \textbf{nombre} (singulier, pluriel) de son antécédent.
\item Le \textbf{verbe conjugué} de la relative va \textbf{à la toute dernière place} de la subordonnée (comme dans toutes les subordonnées allemandes).
\end{enumerate}
\end{propriete}

\begin{propriete}[Le pronom relatif sujet (nominatif)]
Quand le pronom relatif est \textbf{sujet} du verbe de la relative (c'est le cas le plus fréquent, il correspond au français « qui »), il est au \textbf{nominatif} :

\begin{center}
\begin{tabular}{|l|l|l|}
\hline
\rowcolor{popblueL} Antécédent & Pronom relatif & Exemple \\
\hline
masculin singulier & \all{der} & \all{der Mann, der ...} \\
\hline
féminin singulier & \all{die} & \all{die Frau, die ...} \\
\hline
neutre singulier & \all{das} & \all{das Kind, das ...} \\
\hline
pluriel (tous genres) & \all{die} & \all{die Leute, die ...} \\
\hline
\end{tabular}
\end{center}

Le verbe conjugué de la relative va à la fin : \all{Der Mann, der nebenan wohnt, ist nett.} « L'homme qui habite à côté est gentil. »
\end{propriete}

\begin{popfigure}
\begin{center}
\begin{tikzpicture}[node distance=0.4cm, every node/.style={font=\small}]
\node[draw=popblue, fill=popblueL, rounded corners, thick, align=center, minimum width=3.2cm, minimum height=1cm] (ant) {\textbf{Antécédent}\\ \all{eine Baustelle}};
\node[right=0.5cm of ant, font=\Large] (virg) {\textbf{,}};
\node[draw=poporange, fill=poporangeL, rounded corners, thick, align=center, minimum width=2.2cm, minimum height=1cm, right=0.5cm of virg] (pron) {\textbf{Pronom relatif}\\ \all{die}};
\node[draw=popgreen, fill=popgreenL, rounded corners, thick, align=center, minimum width=4.6cm, minimum height=1cm, right=0.5cm of pron] (verb) {\textbf{Reste + verbe à la fin}\\ \all{jeden Morgen anfängt}};
\draw[-{Stealth}, very thick, poppurple, bend right=35] (ant.south) to node[below, align=center, text=poppurple]{genre + nombre\\ (féminin singulier)} (pron.south);
\end{tikzpicture}
\end{center}
\legende{Structure de la relative : le pronom relatif hérite du genre et du nombre de l'antécédent ; le verbe conjugué part en fin de subordonnée.}
\end{popfigure}

\courssub{Le choix du pronom relatif : genre et nombre de l'antécédent}

C'est le point capital de la leçon. En français, « qui » ne change jamais. En allemand, il faut d'abord se poser deux questions :

\begin{methode}[Choisir le bon pronom relatif en trois étapes]
\begin{enumerate}
\item \textbf{Identifier l'antécédent} : quel nom la relative décrit-elle ? \all{Das Geräusch, ... ist laut.} → l'antécédent est \all{das Geräusch}.
\item \textbf{Déterminer son genre et son nombre} : \all{das Geräusch} est neutre singulier.
\item \textbf{Choisir le pronom en conséquence} : neutre singulier → \all{das}. Résultat : \all{Das Geräusch, das laut ist, stört mich.}
\end{enumerate}
\end{methode}

\begin{exemplebox}[Exemples traduits (pronom relatif SUJET = nominatif)]
\all{Der Lärm, der aus der Baustelle kommt, ist schrecklich.}\\
« Le bruit qui vient du chantier est terrible. » (\all{der Lärm} : masculin → \all{der})

\all{Das ist die Nachbarin, die jeden Tag lärmt.}\\
« C'est la voisine qui fait du bruit chaque jour. » (\all{die Nachbarin} : féminin → \all{die})

\all{Das Kind, das im Garten spielt, ist mein Bruder.}\\
« L'enfant qui joue dans le jardin est mon frère. » (\all{das Kind} : neutre → \all{das})

\all{Die Kinder, die im Hof spielen, sind sehr laut.}\\
« Les enfants qui jouent dans la cour sont très bruyants. » (pluriel → \all{die})

\all{Der bendskin, der vor dem Haus hupt, weckt mich auf.}\\
« La moto-taxi qui klaxonne devant la maison me réveille. » (\all{der} : masculin)
\end{exemplebox}

Remarquez bien la place du verbe conjugué : \all{kommt}, \all{lärmt}, \all{spielen}, \all{hupt} se trouvent \textbf{tous à la fin} de la relative. Et si la relative se termine par un verbe à particule séparable, la particule se \textbf{recolle} au verbe en fin de subordonnée : \all{anfängt} (et non \all{fängt ... an}), \all{weckt mich auf}.

\begin{attention}[Cas particulier : le pronom relatif COD (accusatif) — aperçu]
Dans les exemples ci-dessus, le pronom relatif est toujours \textbf{sujet} (« qui »). Il arrive que le pronom relatif soit \textbf{complément d'objet»} (« que »). Il passe alors à l'\textbf{accusatif}.
\begin{itemize}
\item Seule la forme \textbf{masculin singulier} change : \all{der} $\rightarrow$ \all{den}.
\item Féminin, neutre, pluriel : formes identiques au nominatif (\all{die, das, die}).
\end{itemize}
\all{Der Lärm, den ich hasse, kommt von der Baustelle.}\\
« Le bruit \textbf{que} je déteste vient du chantier. » (\all{ich} est le sujet de la relative, \all{den} est COD).
\emph{Vous approfondirez ce point en Terminale ; pour l'instant, retenez que si la relative a déjà un sujet (ex. \all{ich, mein Bruder}), le pronom relatif n'est pas sujet.}
\end{attention}

\courssub{Comparaison français / allemand}

\begin{center}
\begin{tabularx}{\linewidth}{|X|X|X|}
\hline
\rowcolor{popblueL} Français & Deutsch & Remarque \\
\hline
l'homme qui habite ici & \all{der Mann, der hier wohnt} & « qui » = \all{der} (masculin) \\
\hline
la femme qui parle & \all{die Frau, die spricht} & « qui » = \all{die} (féminin) \\
\hline
l'enfant qui joue & \all{das Kind, das spielt} & « qui » = \all{das} (neutre) \\
\hline
les voisins qui font du bruit & \all{die Nachbarn, die Lärm machen} & pluriel : toujours \all{die} \\
\hline
la jeune fille qui dort & \all{das Mädchen, das schläft} & \all{Mädchen} est neutre ! \\
\hline
\end{tabularx}
\end{center}

La grande différence : le français accorde rien du tout (« qui » est invariable), tandis que l'allemand accorde le pronom relatif avec le \textbf{genre grammatical} du nom, et non avec le sens ou le sexe réel de la personne.

\begin{attention}[das Mädchen, das ... et non die !]
En français, « la jeune fille » est féminin, donc on est tenté d'écrire \all{die Mädchen, die...}. Erreur ! En allemand, \all{das Mädchen} est \textbf{neutre} (à cause du suffixe diminutif \all{-chen}, toujours neutre). Le pronom relatif suit le genre grammatical :

\all{Das Mädchen, das im Zimmer schläft, ist meine Schwester.}\\
« La jeune fille qui dort dans la chambre est ma sœur. »

De même : \all{die Straße, die sehr laut ist} (féminin), \all{das Radio, das zu laut spielt} (neutre). Vérifiez toujours l'article du nom avant de choisir le pronom relatif !
\end{attention}

\begin{attention}[La virgule est obligatoire en allemand]
En français, on écrit : « L'homme qui habite à côté est gentil », sans virgule. En allemand, la virgule avant la relative est \textbf{toujours obligatoire}, et si la relative est insérée au milieu de la phrase, il faut une virgule \textbf{des deux côtés} :

\all{Der Mann, der nebenan wohnt, ist nett.}\\
« L'homme qui habite à côté est gentil. »

Oublier la virgule est une faute d'orthographe en allemand, sanctionnée à l'écrit.
\end{attention}

\courssub{Texte d'application : le bruit dans mon quartier}

\begin{textebox}[Lärm in meinem Viertel]
\all{Ich wohne in einem Viertel von Yaoundé, das sehr lebendig ist. Direkt neben unserem Haus gibt es eine Baustelle, die jeden Morgen um sechs Uhr anfängt. Die Arbeiter, die dort bauen, machen viel Lärm mit ihren Maschinen. Das Geräusch, das ich am meisten hasse, ist der Presslufthammer, der den ganzen Tag hämmert.}

\all{Mein Nachbar, der ein kleines Geschäft hat, spielt oft laute Musik. Seine Lautsprecher, die vor dem Laden stehen, stören die ganze Straße. Meine Schwester, die für das Bac lernt, kann sich nicht konzentrieren. Sie sagt: „Ich brauche Ruhe!“}

\all{Aber es gibt auch schöne Geräusche: die Vögel, die am Morgen singen, und die Kinder, die im Hof lachen. Am Abend, wenn die Baustelle schließt und die Musik endet, wird es endlich ruhig. Dann höre ich nur noch den Wind, der leise durch die Bäume geht. Das ist die beste Zeit des Tages.}
\end{textebox}

\textbf{Glossaire :}

\begin{itemize}
\item \all{das Viertel, -} : le quartier (pl. \all{die Viertel})
\item \all{lebendig} : vivant, animé
\item \all{der Presslufthammer, die Presslufthämmer} : le marteau-piqueur
\item \all{hämmern} : marteler
\item \all{das Geschäft, -e} : le magasin, la boutique
\item \all{der Lautsprecher, -} : le haut-parleur (pl. \all{die Lautsprecher})
\item \all{sich konzentrieren} : se concentrer
\item \all{schließen} : fermer
\item \all{der Wind, die Winde} : le vent
\end{itemize}

\textbf{Questions de compréhension :}

\begin{enumerate}
\item \all{Wann fängt die Baustelle an?} (Quand le chantier commence-t-il ?)
\item \all{Welches Geräusch hasst der Erzähler am meisten?} (Quel bruit le narrateur déteste-t-il le plus ?)
\item \all{Warum kann die Schwester nicht lernen?} (Pourquoi la sœur ne peut-elle pas travailler ?)
\item \all{Welche schönen Geräusche gibt es im Viertel?} (Quels beaux bruits y a-t-il dans le quartier ?)
\end{enumerate}

\textbf{Repérage grammatical :} soulignez dans le texte toutes les propositions relatives et entourez chaque antécédent. Il y en a \textbf{onze}. Vérifiez à chaque fois que le pronom relatif a bien le genre et le nombre de l'antécédent, et que le verbe conjugué est en fin de subordonnée. Identifiez aussi les deux cas où le pronom relatif n'est pas sujet (COD).

\courssub{Exercice résolu}

\begin{exoresolu}[Complétez avec le pronom relatif correct (nominatif / accusatif)]
Énoncé. Complétez avec \all{der}, \all{die}, \all{das} ou \all{den} :

1. \all{Der Wecker, ... jeden Morgen klingelt, ist zu laut.}\\
2. \all{Die Musik, ... mein Bruder hört, stört mich.}\\
3. \all{Das Radio, ... in der Küche steht, ist kaputt.}\\
4. \all{Die Nachbarn, ... nebenan wohnen, sind sehr nett.}\\
5. \all{Das Mädchen, ... dort spielt, heißt Amina.}\\
6. \all{Der Lärm, ... ich nicht mag, kommt von der Straße.}

\tcblower

Solution détaillée.

1. \all{Der Wecker, \textbf{der} jeden Morgen klingelt, ist zu laut.} — \all{der Wecker} est masculin singulier, le pronom est \textbf{sujet} de \all{klingelt} → nominatif masculin : \all{der}. « Le réveil qui sonne chaque matin est trop bruyant. »

2. \all{Die Musik, \textbf{die} mein Bruder hört, stört mich.} — \all{die Musik} est féminin singulier. Le sujet de la relative est \all{mein Bruder} → le pronom est \textbf{COD} (accusatif). Féminin accusatif = \all{die}. « La musique \textbf{que} mon frère écoute me dérange. »

3. \all{Das Radio, \textbf{das} in der Küche steht, ist kaputt.} — \all{das Radio} est neutre, pronom \textbf{sujet} de \all{steht} → nominatif neutre : \all{das}. « La radio qui se trouve dans la cuisine est en panne. »

4. \all{Die Nachbarn, \textbf{die} nebenan wohnen, sind sehr nett.} — pluriel → \all{die} (sujet de \all{wohnen}). « Les voisins qui habitent à côté sont très gentils. »

5. \all{Das Mädchen, \textbf{das} dort spielt, heißt Amina.} — Piège classique ! \all{Mädchen} est neutre malgré le sens féminin → \all{das} (sujet de \all{spielt}). « La jeune fille qui joue là-bas s'appelle Amina. »

6. \all{Der Lärm, \textbf{den} ich nicht mag, kommt von der Straße.} — \all{der Lärm} (masculin). Sujet de la relative = \all{ich} → pronom = \textbf{COD} (accusatif) → \all{den}. « Le bruit \textbf{que} je n'aime pas vient de la rue. »
\end{exoresolu}

\begin{exercice}[À vous de jouer]
Traduisez en allemand en utilisant une proposition relative :

1. Le chantier qui commence à six heures est très bruyant. (\all{die Baustelle})
2. La voisine qui écoute de la musique forte s'appelle Madame Ngo. (\all{die Nachbarin})
3. Le bruit que j'entends est un marteau-piqueur. (\all{das Geräusch})
4. Les enfants qui jouent dans la cour sont mes cousins. (\all{die Kinder})
\end{exercice}

\begin{corrige}[Exercice « À vous de jouer »]
1. \all{Die Baustelle, die um sechs Uhr anfängt, ist sehr laut.}\\
2. \all{Die Nachbarin, die laute Musik hört, heißt Frau Ngo.}\\
3. \all{Das Geräusch, das ich höre, ist ein Presslufthammer.}\\
4. \all{Die Kinder, die im Hof spielen, sind meine Cousins.}

Points de vigilance : virgule avant la relative, verbe conjugué à la fin, particule séparable recollée (\all{anfängt}), pronom accordé au genre grammatical. À la phrase 3, \all{das} est COD (accusatif neutre = \all{das}).
\end{corrige}

\courssub{Pour aller plus loin : le relatif à l'accusatif (niveau Bac)}

\begin{savaistu}[Le bruit « que » je déteste]
Jusqu'ici, le pronom relatif était \textbf{sujet»} (« qui »). Mais il peut aussi être \textbf{complément d'objet direct»} (« que »). Il passe alors à l'\textbf{accusatif}. Bonne nouvelle : seules les formes masculines changent !

\begin{center}
\begin{tabular}{|l|l|l|}
\hline
\rowcolor{popgreenL} Genre & Nominatif (sujet = qui) & Accusatif (COD = que) \\
\hline
masculin & \all{der} & \all{den} \\
\hline
féminin & \all{die} & \all{die} \\
\hline
neutre & \all{das} & \all{das} \\
\hline
pluriel & \all{die} & \all{die} \\
\hline
\end{tabular}
\end{center}

\all{Der Lärm, den ich hasse, kommt von der Baustelle.}\\
« Le bruit \textbf{que} je déteste vient du chantier. »

Comment le reconnaître ? Si la relative contient déjà un sujet (\all{ich, du, er, mein Vater}...), le pronom relatif ne peut pas être sujet : il est COD $\rightarrow$ \all{den} au masculin. Vous approfondirez ce point (et le datif/genitif) en classe de Terminale ; retenez pour l'instant la phrase modèle ci-dessus et le tableau.
\end{savaistu}

\begin{aretenir}[L'essentiel de la section]
\begin{itemize}
\item La proposition relative précise un nom ; elle est introduite par un pronom relatif après une \textbf{virgule obligatoire}, et le verbe conjugué va \textbf{à la fin}.
\item Au nominatif (sujet = « qui ») : \all{der} (masculin), \all{die} (féminin), \all{das} (neutre), \all{die} (pluriel).
\item Le pronom relatif suit le \textbf{genre grammatical} de l'antécédent, pas le sens : \all{das Mädchen, das...}.
\item Les verbes à particule séparable se recollent en fin de relative : \all{die Baustelle, die um sechs Uhr anfängt}.
\item Aperçu Bac : à l'accusatif (COD = « que »), seul le masculin change : \all{den Lärm, den ich hasse}. Féminin, neutre, pluriel inchangés.
\end{itemize}
\end{aretenir}

\coursec{Conjugaison : prétérit de sein, haben et des modaux}

Dans la section précédente, nous avons appris à relier deux idées grâce aux propositions relatives : \all{Die Nachbarn, die viel Lärm machen, wohnen nebenan.} Revenons maintenant au texte \all{Lärm in der Stadt} et regardons de près ses verbes. Trois formes reviennent sans cesse : \all{war}, \all{hatte}, \all{konnte}. À quel temps sont-elles conjuguées ? C'est ce que nous allons découvrir : le \cle{Präteritum} (le prétérit), le temps de référence pour raconter le passé avec \all{sein}, \all{haben} et les verbes modaux.

\courssub{Observation : war, hatte, konnte dans le texte}

\begin{textebox}[Extraits du texte support]
„Gestern Abend \textbf{war} es in unserer Straße sehr laut. Die Nachbarn \textbf{hatten} eine Party, und die Musik \textbf{war} bis zwei Uhr nachts zu hören. Ich \textbf{konnte} nicht schlafen, weil der Lärm so stark \textbf{war}. Mein Bruder \textbf{musste} heute früh aufstehen, aber er \textbf{hatte} keine Energie. Wir \textbf{wollten} uns bei den Nachbarn beschweren, aber wir \textbf{durften} nicht, denn unsere Eltern \textbf{waren} nicht zu Hause.“
\end{textebox}

Observons et complétons mentalement le tableau d'identification :

\begin{center}
\begin{tabularx}{\linewidth}{|X|X|X|X|}
\hline
\rowcolor{popblueL} Forme dans le texte & Infinitif & Personne & Sens \\
\hline
war & sein & 3\textsuperscript{e} pers. sg. (es) & était \\
\hline
hatten & haben & 3\textsuperscript{e} pers. pl. (die Nachbarn) & avaient \\
\hline
konnte & können & 1\textsuperscript{re} pers. sg. (ich) & pouvais \\
\hline
musste & müssen & 3\textsuperscript{e} pers. sg. (mein Bruder) & devait \\
\hline
wollten & wollen & 1\textsuperscript{re} pers. pl. (wir) & voulions \\
\hline
durften & dürfen & 1\textsuperscript{re} pers. pl. (wir) & avions le droit \\
\hline
waren & sein & 3\textsuperscript{e} pers. pl. (unsere Eltern) & étaient \\
\hline
\end{tabularx}
\end{center}

\textbf{Conclusion de l'observation :} ces formes expriment toutes le passé, mais ce ne sont \emph{pas} des parfaits (\all{ich bin gewesen}, \all{ich habe gehabt}). Ce sont des formes simples, en un seul mot, sans auxiliaire : c'est le \cle{Präteritum}.

\courssub{Le prétérit de sein : war}

\begin{propriete}[Prétérit de \all{sein} (être)]
Le verbe \all{sein} est totalement irrégulier au prétérit. Il faut apprendre ses formes par cœur :
\end{propriete}

\begin{center}
\begin{tabular}{|l|l|l|}
\hline
\rowcolor{popblueL} Personne & Forme & Exemple \\
\hline
ich & war & Ich war müde. (J'étais fatigué.) \\
\hline
du & warst & Du warst krank. (Tu étais malade.) \\
\hline
er/sie/es & war & Es war sehr laut. (C'était très bruyant.) \\
\hline
wir & waren & Wir waren zu Hause. (Nous étions à la maison.) \\
\hline
ihr & wart & Ihr wart in der Stadt. (Vous étiez en ville.) \\
\hline
sie/Sie & waren & Sie waren in Yaoundé. (Ils étaient à Yaoundé.) \\
\hline
\end{tabular}
\end{center}

Remarques importantes :
\begin{itemize}
\item La 1\textsuperscript{re} et la 3\textsuperscript{e} personne du singulier sont \textbf{identiques} et \textbf{sans terminaison} : \all{ich war}, \all{er war}. C'est une particularité générale du prétérit allemand.
\item Attention à la 2\textsuperscript{e} personne du pluriel : \all{ihr wart} (sans \emph{e} : on ne dit pas \all{ihr waret}).
\end{itemize}

\courssub{Le prétérit de haben : hatte}

\begin{propriete}[Prétérit de \all{haben} (avoir)]
Le prétérit de \all{haben} se construit sur le radical \all{hatt-} + les terminaisons du prétérit :
\end{propriete}

\begin{center}
\begin{tabular}{|l|l|l|}
\hline
\rowcolor{popblueL} Personne & Forme & Exemple \\
\hline
ich & hatte & Ich hatte Kopfschmerzen. (J'avais mal à la tête.) \\
\hline
du & hattest & Du hattest keine Zeit. (Tu n'avais pas le temps.) \\
\hline
er/sie/es & hatte & Sie hatte Hunger. (Elle avait faim.) \\
\hline
wir & hatten & Wir hatten keine Ruhe. (Nous n'avions pas de calme.) \\
\hline
ihr & hattet & Ihr hattet Glück. (Vous avez eu de la chance.) \\
\hline
sie/Sie & hatten & Sie hatten eine Party. (Ils avaient une fête.) \\
\hline
\end{tabular}
\end{center}

\courssub{Le prétérit des verbes modaux}

\begin{propriete}[Prétérit des modaux]
Les verbes modaux forment leur prétérit sur un radical \textbf{sans umlaut}, auquel s'ajoutent les mêmes terminaisons que \all{haben} (1\textsuperscript{re} et 3\textsuperscript{e} pers. sg. identiques, sans terminaison) :
\end{propriete}

\begin{center}
\begin{tabularx}{\linewidth}{|X|X|X|}
\hline
\rowcolor{popblueL} Infinitif (Präsens) & Präteritum & Sens au passé \\
\hline
können (kann) & konnte & pouvait, a pu \\
\hline
müssen (muss) & musste & devait, a dû \\
\hline
dürfen (darf) & durfte & avait le droit, a eu la permission \\
\hline
wollen (will) & wollte & voulait, a voulu \\
\hline
sollen (soll) & sollte & devait (obligation, conseil) \\
\hline
mögen (mag) & mochte & aimait \\
\hline
\end{tabularx}
\end{center}

Conjugaison complète de \all{können} au prétérit (le modèle des terminaisons est identique pour tous les modaux) :

\begin{center}
\begin{tabular}{|l|l|l|}
\hline
\rowcolor{popblueL} Personne & Forme & Exemple \\
\hline
ich & konnte & Ich konnte nicht schlafen. (Je ne pouvais pas dormir.) \\
\hline
du & konntest & Du konntest nichts hören. (Tu ne pouvais rien entendre.) \\
\hline
er/sie/es & konnte & Er konnte nicht lernen. (Il ne pouvait pas travailler.) \\
\hline
wir & konnten & Wir konnten nicht spielen. (Nous ne pouvions pas jouer.) \\
\hline
ihr & konntet & Ihr konntet nicht kommen. (Vous ne pouviez pas venir.) \\
\hline
sie/Sie & konnten & Sie konnten nichts machen. (Ils ne pouvaient rien faire.) \\
\hline
\end{tabular}
\end{center}

Sur le même modèle : \all{ich musste, du musstest, er musste, wir mussten, ihr musstet, sie mussten} ; \all{ich durfte, du durftest, er durfte, wir durften, ihr durftet, sie durften} ; \all{ich wollte, du wolltest, er wollte, wir wollten, ihr wolltet, sie wollten}.

\begin{attention}[Perte de l'umlaut au prétérit]
Les modaux à umlaut le \textbf{perdent} au prétérit : \all{können} $\rightarrow$ \all{konnte}, \all{müssen} $\rightarrow$ \all{musste}, \all{dürfen} $\rightarrow$ \all{durfte}, \all{mögen} $\rightarrow$ \all{mochte}. Attention : \all{könnte} avec umlaut n'est \emph{pas} un prétérit, c'est le \cle{Konjunktiv II} et il signifie « pourrait » (nuance d'hypothèse ou de politesse : \all{Ich könnte kommen} = « je pourrais venir »). Cette forme sera étudiée en classe de Terminale ; pour l'instant, retenez : passé = \textbf{sans} umlaut.
\end{attention}

\courssub{Emploi : pourquoi le prétérit et pas le parfait ?}

\begin{propriete}[Règle d'emploi]
\all{sein}, \all{haben} et les verbes modaux s'emploient \textbf{presque toujours au prétérit} (et non au parfait) pour exprimer le passé, à l'écrit \emph{comme à l'oral}. On dit :
\begin{itemize}
\item \all{Ich war müde.} (et non \all{ich bin müde gewesen})
\item \all{Ich hatte Kopfschmerzen.} (et non \all{ich habe Kopfschmerzen gehabt})
\item \all{Ich konnte nicht lernen.} (et non \all{ich habe nicht lernen können})
\end{itemize}
C'est une exception majeure à la règle générale selon laquelle l'oral allemand préfère le parfait. Pour ces trois familles de verbes, le prétérit est le temps \emph{normal} du passé, même dans la conversation quotidienne. Pour tous les autres verbes, rappelez-vous : le parfait domine à l'oral (\all{ich habe geschlafen}), le prétérit domine dans les récits écrits (\all{ich schlief}).
\end{propriete}

Comparaison avec le français :

\begin{center}
\begin{tabularx}{\linewidth}{|X|X|X|}
\hline
\rowcolor{popblueL} Français & Deutsch & Remarque \\
\hline
J'étais fatigué. & Ich war müde. & prétérit simple, d'usage courant (contrairement à « je fus ») \\
\hline
Nous avions mal à la tête. & Wir hatten Kopfschmerzen. & \all{haben} + nom, sans article \\
\hline
Elle devait se lever tôt. & Sie musste früh aufstehen. & modal au prétérit + infinitif final \\
\hline
Il voulait dormir, mais il ne pouvait pas. & Er wollte schlafen, aber er konnte nicht. & l'infinitif peut être sous-entendu, comme en français \\
\hline
\end{tabularx}
\end{center}

\begin{exemplebox}[Exemples traduits]
\begin{itemize}
\item \all{Gestern war ich sehr müde.} « Hier j'étais très fatigué. »
\item \all{Wir hatten keine Ruhe.} « Nous n'avions pas de calme. »
\item \all{Sie musste früh aufstehen.} « Elle devait se lever tôt. »
\item \all{Er wollte schlafen, aber er konnte nicht.} « Il voulait dormir, mais il ne pouvait pas. »
\item \all{Die Kinder durften am Abend nicht fernsehen.} « Les enfants n'avaient pas le droit de regarder la télé le soir. »
\item \all{Als ich klein war, mochte ich laute Musik nicht.} « Quand j'étais petit, je n'aimais pas la musique forte. »
\end{itemize}
\end{exemplebox}

\courssub{La structure de la phrase avec un modal au prétérit}

\begin{propriete}[Place des verbes]
Avec un modal au prétérit, la structure est la même qu'au présent : le modal conjugué occupe la \textbf{2\textsuperscript{e} position}, et l'infinitif se place à la \textbf{toute fin} de la phrase :
\begin{center}
\all{Ich} \quad \textbf{konnte} \quad \all{wegen des Lärms nicht} \quad \textbf{schlafen}.\\
(sujet) \quad (modal V2) \quad (milieu de phrase) \quad (infinitif final)
\end{center}
« Je ne pouvais pas dormir à cause du bruit. »
\end{propriete}

\begin{attention}[Pièges des francophones]
\begin{itemize}
\item Ne placez jamais l'infinitif juste après le modal comme en français (« je pouvais dormir »). En allemand, tout le milieu de la phrase passe \emph{avant} l'infinitif : \all{Ich konnte gestern wegen des Lärms bis zwei Uhr nicht schlafen.}
\item N'ajoutez pas \all{zu} devant l'infinitif après un modal : \all{ich konnte schlafen} (et non \all{ich konnte zu schlafen}).
\item Ne confondez pas \all{war} (était) et \all{wahr} (vrai) : même prononciation « var », mais sens et orthographe différents.
\end{itemize}
\end{attention}

\begin{popfigure}
\begin{tikzpicture}[scale=1]
\draw[-{Stealth[length=3mm]}, thick, popdark] (-0.5,0) -- (11.5,0);
\foreach \x/\lab in {1.5/{Vergangenheit}, 5.5/{Gegenwart}, 9.5/{Zukunft}} {
  \draw[thick, popdark] (\x,0.15) -- (\x,-0.15);
  \node[below, popdark, align=center] at (\x,-0.2) {\lab};
}
\node[above, popblue, align=center] at (1.5,0.25) {\textbf{Präteritum}\\ \all{war, hatte, konnte}};
\node[above, popgreen, align=center] at (5.5,0.25) {\textbf{Präsens}\\ \all{bin, habe, kann}};
\node[above, poporange, align=center] at (9.5,0.25) {\textbf{Futur}\\ \all{werde sein/haben/können}};
\node[popblue, align=center, fill=popblueL, rounded corners] at (1.5,-1.6) {\all{Gestern war es laut.}\\ \all{Ich konnte nicht schlafen.}};
\node[popgreen, align=center, fill=popgreenL, rounded corners] at (5.5,-1.6) {\all{Heute bin ich müde.}\\ \all{Ich habe Kopfschmerzen.}};
\node[poporange, align=center, fill=poporangeL, rounded corners] at (9.5,-1.6) {\all{Morgen werde ich}\\ \all{früh schlafen.}};
\end{tikzpicture}
\legende{Frise des temps : le prétérit de \all{sein}, \all{haben} et des modaux raconte le passé.}
\end{popfigure}

\courssub{Entraînement : du présent au prétérit}

\begin{exoresolu}[Transformer au prétérit]
Énoncé : mettez les phrases suivantes au prétérit. 1. \all{Ich bin sehr müde.} 2. \all{Wir haben keine Ruhe.} 3. \all{Sie kann nicht schlafen.} 4. \all{Die Schüler müssen viel lernen.} 5. \all{Er will sich beschweren, aber er darf nicht.}
\tcblower
Solution détaillée : on identifie le verbe conjugué, on le remplace par sa forme de prétérit à la même personne, et l'infinitif éventuel reste en fin de phrase, inchangé.
\begin{enumerate}
\item \all{Ich \textbf{war} sehr müde.} (\all{bin} $\rightarrow$ \all{war})
\item \all{Wir \textbf{hatten} keine Ruhe.} (\all{haben} $\rightarrow$ \all{hatten})
\item \all{Sie \textbf{konnte} nicht schlafen.} (\all{kann} $\rightarrow$ \all{konnte}, infinitif \all{schlafen} en fin de phrase)
\item \all{Die Schüler \textbf{mussten} viel lernen.} (\all{müssen} $\rightarrow$ \all{mussten}, 3\textsuperscript{e} pers. pl.)
\item \all{Er \textbf{wollte} sich beschweren, aber er \textbf{durfte} nicht.} (\all{will} $\rightarrow$ \all{wollte} ; \all{darf} $\rightarrow$ \all{durfte})
\end{enumerate}
\end{exoresolu}

\begin{exoresolu}[Thème guidé : traduction]
Énoncé : traduisez en allemand. 1. « Hier soir, c'était très bruyant dans notre rue. » 2. « Je ne pouvais pas dormir à cause du bruit. » 3. « Mes parents n'étaient pas à la maison. » 4. « Nous voulions nous plaindre, mais nous n'avions pas le droit. »
\tcblower
Solution détaillée :
\begin{enumerate}
\item \all{Gestern Abend \textbf{war} es in unserer Straße sehr laut.} (\all{sein} au prétérit, 3\textsuperscript{e} pers. sg. ; « dans notre rue » = \all{in unserer Straße}, datif après \all{in} pour le lieu)
\item \all{Ich \textbf{konnte} wegen des Lärms nicht schlafen.} (modal en V2, infinitif \all{schlafen} en toute fin ; \all{wegen} + génitif : \all{des Lärms})
\item \all{Meine Eltern \textbf{waren} nicht zu Hause.} (\all{sein} au prétérit, 3\textsuperscript{e} pers. pl.)
\item \all{Wir \textbf{wollten} uns beschweren, aber wir \textbf{durften} nicht.} (verbe réfléchi \all{sich beschweren} : \all{uns} à la 1\textsuperscript{re} pers. pl. ; l'infinitif est sous-entendu après \all{durften nicht}, comme en français)
\end{enumerate}
\end{exoresolu}

\begin{exercice}[Mini-récit au prétérit]
Intégrez les cinq phrases transformées dans un petit récit cohérent sur une nuit bruyante. Commencez par : \all{Gestern Abend war es in unserer Straße sehr laut ...} Utilisez au moins une subordonnée avec \all{als} et une proposition relative.
\end{exercice}

\begin{corrige}[Mini-récit (proposition)]
\all{Gestern Abend war es in unserer Straße sehr laut. Die Nachbarn, die nebenan wohnen, hatten eine Party. Ich war sehr müde, aber ich konnte nicht schlafen. Als die Musik um Mitternacht noch lauter wurde, hatten wir keine Ruhe mehr. Die Schüler mussten am nächsten Tag viel lernen, aber sie waren zu müde. Mein Vater wollte sich beschweren, aber er durfte nicht, denn die Nachbarn sind unsere Freunde.} « Hier soir, c'était très bruyant dans notre rue. Les voisins, qui habitent à côté, avaient une fête. J'étais très fatigué, mais je ne pouvais pas dormir. Quand la musique est devenue encore plus forte à minuit, nous n'avions plus de calme. Les élèves devaient beaucoup travailler le lendemain, mais ils étaient trop fatigués. Mon père voulait se plaindre, mais il ne pouvait pas, car les voisins sont nos amis. »
\end{corrige}

\begin{aretenir}[L'essentiel]
\begin{itemize}
\item \all{sein} $\rightarrow$ \all{war} (ich/er war, du warst, wir/sie waren, ihr wart) ; \all{haben} $\rightarrow$ \all{hatte}.
\item Modaux au prétérit \textbf{sans umlaut} : \all{konnte, musste, durfte, wollte, sollte, mochte}.
\item 1\textsuperscript{re} et 3\textsuperscript{e} personnes du singulier identiques, sans terminaison : \all{ich konnte} = \all{er konnte}.
\item Ces verbes expriment le passé au prétérit même à l'oral : jamais \all{ich bin gewesen} dans la langue courante.
\item Structure : sujet + modal au prétérit (V2) + ... + infinitif en toute fin de phrase, sans \all{zu}.
\end{itemize}
\end{aretenir}

\coursec{Production guidée et méthode du commentaire}

Dans la section précédente, nous avons appris à conjuguer au prétérit les verbes \all{sein}, \all{haben} et les verbes modaux (\all{ich war}, \all{ich hatte}, \all{ich konnte}, \all{ich musste}...). Ces formes vont maintenant nous servir concrètement : pour raconter une nuisance subie, rédiger une plainte polie et commenter le texte \all{Lärm in der Stadt}. C'est l'objectif de cette dernière section : passer de la \cle{compréhension} à la \cle{production}, écrite et orale, en réutilisant tout ce que la leçon nous a appris — le vocabulaire du bruit, les subordonnées avec \all{wenn} et \all{als}, les propositions relatives et le prétérit.

\courssub{La méthode du commentaire de texte en 4 étapes}

Au lycée, on vous demandera souvent de « commenter » un texte allemand. Un commentaire n'est ni une traduction ni un simple résumé : c'est un texte organisé qui présente, résume, analyse et juge le document. Retenez la méthode en quatre étapes.

\begin{methode}[Commenter un texte en 4 étapes]
\begin{enumerate}
\item \textbf{Présenter} : indiquer le titre, le type de texte (article, lettre, dialogue), le thème et éventuellement l'auteur ou la source. Formule utile : \all{Der Text „...“ handelt von...} (Le texte « ... » traite de...).
\item \textbf{Résumer} : dégager les idées principales en 2 ou 3 phrases, avec ses propres mots, sans recopier le texte.
\item \textbf{Analyser} : observer le vocabulaire (champ lexical du bruit : \all{der Lärm, laut, stören}), les structures grammaticales (subordonnées avec \all{wenn/als}, relatives) et le point de vue de l'auteur (est-il critique, neutre, optimiste ?).
\item \textbf{Donner son avis} : exprimer une opinion argumentée avec \all{meiner Meinung nach} (à mon avis), \all{ich finde, dass...} (je trouve que...), \all{ich bin der Meinung, dass...} (je suis d'avis que...).
\end{enumerate}
\end{methode}

\begin{popfigure}
\begin{tikzpicture}[node distance=0.55cm,
etape/.style={rectangle, rounded corners=3pt, draw=popblue, fill=popblueL, text width=5.2cm, align=center, minimum height=0.9cm, font=\small},
fleche/.style={-{Stealth[length=2.5mm]}, thick, poporange}]
\node[etape, align=center] (e1){\textbf{1. Présentation}\\ Titre, type de texte, thème};
\node[etape, below=of e1, align=center] (e2){\textbf{2. Résumé}\\ Idées principales en 2--3 phrases};
\node[etape, below=of e2, align=center] (e3){\textbf{3. Analyse}\\ Vocabulaire, structures, point de vue};
\node[etape, below=of e3, align=center] (e4){\textbf{4. Avis personnel}\\ \all{meiner Meinung nach...}};
\draw[fleche] (e1) -- (e2);
\draw[fleche] (e2) -- (e3);
\draw[fleche] (e3) -- (e4);
\end{tikzpicture}
\legende{Organigramme des 4 étapes du commentaire de texte}
\end{popfigure}

\courssub{Commentaire modèle sur le texte étudié}

Voici un commentaire rédigé sur le texte \all{Lärm in der Stadt}. Observez les connecteurs qui donnent sa logique au texte : \all{zuerst} (d'abord), \all{dann} (ensuite), \all{außerdem} (de plus), \all{deshalb} (c'est pourquoi), \all{meiner Meinung nach} (à mon avis).

\begin{textebox}[Kommentar — Modell]
Der Text „Lärm in der Stadt“ handelt von den Geräuschen und dem Lärm in einer großen Stadt. Zuerst beschreibt der Autor die vielen Lärmquellen: Autos, Motorräder, Baustellen und laute Nachbarn. Dann erklärt er die Folgen des Lärms: Man kann nicht schlafen, man ist müde und nervös. Außerdem zeigt der Text, dass viele Menschen sich mehr Ruhe wünschen. Der Autor benutzt oft das Wortfeld „Lärm“: laut, leise, stören, die Ruhe. Er ist kritisch, aber nicht pessimistisch. Meiner Meinung nach ist das Thema sehr wichtig, denn auch in Yaoundé und Douala gibt es viel Lärm. Ich finde, dass die Nachbarn mehr Rücksicht nehmen müssen. Deshalb ist dieser Text aktuell und nützlich für uns alle.
\end{textebox}

\begin{definition}[Glossaire du commentaire]
\begin{itemize}
\item \all{das Geräusch, -e} : le bruit (neutre, plus faible que \all{der Lärm})
\item \all{die Lärmquelle, -n} : la source de bruit
\item \all{die Baustelle, -n} : le chantier
\item \all{die Folge, -n} : la conséquence
\item \all{nervös} : nerveux, irritable
\item \all{sich wünschen} : souhaiter
\item \all{Rücksicht nehmen} : faire preuve de considération
\item \all{aktuell} : d'actualité (faux ami : ne signifie pas « actuel » au sens de « qui existe maintenant » ; pour cela, l'allemand dit \all{gegenwärtig} ou \all{im Moment})
\end{itemize}
\end{definition}

\begin{exemplebox}[Les connecteurs du commentaire]
\begin{itemize}
\item \all{Zuerst beschreibt der Autor die Lärmquellen.} « D'abord, l'auteur décrit les sources de bruit. »
\item \all{Dann erklärt er die Folgen.} « Ensuite, il explique les conséquences. »
\item \all{Außerdem benutzt er oft das Wortfeld „Lärm“.} « De plus, il utilise souvent le champ lexical du bruit. »
\item \all{Deshalb ist der Text aktuell.} « C'est pourquoi le texte est d'actualité. »
\item \all{Meiner Meinung nach ist das Thema wichtig.} « À mon avis, le sujet est important. »
\end{itemize}
\end{exemplebox}

\begin{attention}[Place des connecteurs]
Après \all{zuerst}, \all{dann}, \all{außerdem}, \all{deshalb} et \all{meiner Meinung nach}, le verbe conjugué vient en \cle{deuxième position} : le connecteur occupe la première place. On dit \all{Deshalb ist der Text aktuell} et jamais \all{Deshalb der Text ist aktuell}. C'est l'inversion sujet-verbe, typique de l'allemand : quelle que soit la première place occupée (sujet, complément de temps, connecteur), le verbe reste fidèle à la deuxième place.
\end{attention}

\courssub{Production écrite guidée n°1 : questions et résumé}

\begin{exercice}[Comprendre et résumer]
Répondez en allemand par des phrases complètes, puis rédigez un résumé de 3 phrases.
\begin{enumerate}
\item \all{Was ist das Thema des Textes „Lärm in der Stadt“?}
\item \all{Welche Lärmquellen nennt der Autor?}
\item \all{Was sind die Folgen des Lärms für die Menschen?}
\item \all{Gibt es in Ihrer Stadt auch viel Lärm? Wo?}
\end{enumerate}
\end{exercice}

\begin{exoresolu}[Résumé guidé — corrigé type]
Rédigez un résumé de 3 phrases avec \all{zuerst}, \all{dann} et \all{deshalb}.
\tcblower
\all{Zuerst beschreibt der Text die vielen Lärmquellen in der Stadt, zum Beispiel Autos und Baustellen. Dann erklärt er, dass die Menschen wegen des Lärms nicht schlafen können und nervös sind. Deshalb wünschen sich alle mehr Ruhe.}

« D'abord, le texte décrit les nombreuses sources de bruit en ville, par exemple les voitures et les chantiers. Ensuite, il explique que les gens ne peuvent pas dormir à cause du bruit et qu'ils sont nerveux. C'est pourquoi tous souhaitent plus de calme. »

Remarquez : trois phrases, trois connecteurs, aucun mot recopié tel quel en bloc — c'est un vrai résumé. Notez aussi la subordonnée avec \all{dass} : le verbe conjugué \all{können} va à la fin.
\end{exoresolu}

\courssub{Production écrite guidée n°2 : la lettre de plainte (die Beschwerde)}

Se plaindre en allemand est un exercice codifié : il faut être \cle{ferme} mais \cle{poli}. Une plainte bien construite suit toujours la même architecture.

\begin{methode}[Structure d'une plainte polie]
\begin{enumerate}
\item \textbf{Salutation} : \all{Lieber Herr Ngono,} / \all{Liebe Frau Mbarga,} (Cher Monsieur / Chère Madame).
\item \textbf{Constat} : \all{Seit einer Woche gibt es jeden Abend laute Musik.} (Depuis une semaine, il y a de la musique forte chaque soir.)
\item \textbf{Conséquence} : \all{Ich kann nicht schlafen, weil die Musik so laut ist.} (Je ne peux pas dormir parce que la musique est si forte.)
\item \textbf{Demande} : \all{Ich bitte Sie, die Musik leiser zu machen.} (Je vous prie de baisser la musique.)
\item \textbf{Formule de politesse} : \all{Mit freundlichen Grüßen} (Cordialement) + signature.
\end{enumerate}
\end{methode}

\begin{exemplebox}[Phrases utiles pour la plainte — traduites]
\begin{itemize}
\item \all{Seit Montag gibt es jeden Abend laute Musik.} « Depuis lundi, il y a de la musique forte chaque soir. »
\item \all{Wenn ich schlafen will, höre ich immer Lärm.} « Quand je veux dormir, j'entends toujours du bruit. »
\item \all{Ich bitte Sie, die Musik leiser zu machen.} « Je vous prie de baisser la musique. »
\item \all{Als ich gestern nach Hause kam, war die Musik schon laut.} « Quand je suis rentré hier à la maison, la musique était déjà forte. »
\item \all{Der Lärm, der aus Ihrer Wohnung kommt, stört mich sehr.} « Le bruit qui vient de votre appartement me dérange beaucoup. »
\end{itemize}
\end{exemplebox}

\begin{attention}[Rester poli : le piège de l'impératif]
En allemand comme en français, on ne donne pas d'ordre sec à un inconnu ou à un voisin. Évitez \all{Mach leiser!} (« Baisse ! »), trop brutal. Utilisez toujours :
\begin{itemize}
\item le vouvoiement \all{Sie} (avec majuscule, même au milieu de la phrase !) ;
\item la formule \all{Ich bitte Sie, ... zu + Infinitiv} (je vous prie de...) ;
\item éventuellement le conditionnel de politesse : \all{Könnten Sie bitte die Musik leiser machen?} (Pourriez-vous baisser la musique, s'il vous plaît ?).
\end{itemize}
Une plainte agressive est une plainte ratée — et à l'examen, elle coûte des points de cohérence sociale.
\end{attention}

\begin{textebox}[Modèle de plainte — Eine Beschwerde an den Nachbarn]
Lieber Herr Etoundi,

seit einer Woche gibt es in Ihrer Wohnung jeden Abend laute Musik. Wenn ich abends schlafen will, höre ich immer diesen Lärm, der aus Ihrem Zimmer kommt. Ich kann nicht schlafen, weil die Musik bis Mitternacht so laut ist. Am Montag war ich sehr müde, und ich konnte in der Schule nicht gut arbeiten. Gestern hatte ich auch Kopfschmerzen. Ich habe einen kleinen Bruder, der auch nicht schlafen kann.

Ich bitte Sie, die Musik nach 22 Uhr leiser zu machen. Dann können alle in Ruhe schlafen, und wir bleiben gute Nachbarn.

Vielen Dank und mit freundlichen Grüßen

Amina Bello
\end{textebox}

\begin{definition}[Glossaire de la plainte]
\begin{itemize}
\item \all{die Wohnung, -en} : l'appartement
\item \all{bis Mitternacht} : jusqu'à minuit
\item \all{die Kopfschmerzen (pl.)} : les maux de tête
\item \all{in Ruhe schlafen} : dormir en paix
\item \all{Vielen Dank} : merci beaucoup
\end{itemize}
\end{definition}

Observez comment le modèle respecte les consignes de la leçon : une subordonnée avec \all{wenn} (\all{Wenn ich abends schlafen will...}, verbe \all{will} à la fin), une subordonnée avec \all{weil} (\all{weil die Musik ... so laut ist}), une relative (\all{diesen Lärm, der aus Ihrem Zimmer kommt} — \all{der} est nominatif masculin singulier car il remplace \all{der Lärm} et est sujet de \all{kommt}), et des prétérits (\all{ich war}, \all{ich konnte}, \all{ich hatte}).

\begin{center}
\begin{tabularx}{\linewidth}{|X|X|X|}
\hline
\rowcolor{popblueL} Étape & Français & Deutsch \\
\hline
Salutation & Cher Monsieur Etoundi & Lieber Herr Etoundi, \\
\hline
Constat & Depuis une semaine, il y a... & Seit einer Woche gibt es... \\
\hline
Conséquence & Je ne peux pas dormir parce que... & Ich kann nicht schlafen, weil... \\
\hline
Demande & Je vous prie de baisser... & Ich bitte Sie, ... leiser zu machen. \\
\hline
Politesse & Cordialement & Mit freundlichen Grüßen \\
\hline
\end{tabularx}
\end{center}

\begin{exercice}[À vous d'écrire]
Rédigez une plainte de 6 à 8 lignes à un voisin bruyant. Consignes obligatoires : au moins une subordonnée avec \all{wenn} ou \all{als}, une proposition relative, deux verbes au prétérit, et une demande polie avec \all{Ich bitte Sie, ...}.
\end{exercice}

\begin{corrige}[Exercice : la plainte]
Éléments de corrigé possibles :
\all{Liebe Frau Mbarga, seit drei Tagen höre ich jeden Abend laute Musik aus Ihrer Wohnung. Als ich gestern schlafen wollte, war der Lärm sehr groß. Ich konnte nicht schlafen und ich war heute müde. Der Lärm, der aus Ihrem Zimmer kommt, stört auch meine Familie. Ich bitte Sie, die Musik leiser zu machen. Mit freundlichen Grüßen, Paul Essomba.}

Vérifiez : subordonnée avec \all{als} (événement unique au passé, verbe \all{wollte} à la fin), relative avec \all{der} (nominatif masculin), prétérits \all{war} et \all{konnte}, demande polie. Tout y est.
\end{corrige}

\courssub{Entraînement à la traduction (thème)}

La traduction du français vers l'allemand (le \cle{thème}) fait partie des épreuves de classe. Entraînez-vous sur des phrases qui réutilisent exactement les acquis de la leçon.

\begin{exoresolu}[Thème guidé]
Traduisez en allemand : « Quand je suis rentré hier soir, les voisins faisaient du bruit. Je n'ai pas pu dormir. »
\tcblower
\all{Als ich gestern Abend nach Hause kam, machten die Nachbarn Lärm. Ich konnte nicht schlafen.}

Points de méthode : « quand » au passé pour un événement unique = \all{als} (et non \all{wenn}) ; verbe conjugué \all{kam} à la fin de la subordonnée ; « rentrer à la maison » = \all{nach Hause kommen} ; « je n'ai pas pu » se dit au prétérit \all{ich konnte nicht} (plus naturel que le parfait avec un modal).
\end{exoresolu}

\courssub{Production orale : le dialogue des voisins}

\begin{experience}[Jeu de rôle en binôme : Nachbar gegen Nachbar]
Par groupes de deux, jouez la scène devant la classe.
\begin{itemize}
\item \textbf{Élève A} (le voisin qui se plaint) : saluez poliment, décrivez le problème (musique, fête, chien), expliquez la conséquence (\all{Ich kann nicht schlafen...}), formulez une demande polie.
\item \textbf{Élève B} (le voisin bruyant) : excusez-vous (\all{Es tut mir leid!} — Je suis désolé !), expliquez-vous brièvement, proposez une solution (\all{Ich mache die Musik nach 22 Uhr leiser.}).
\end{itemize}
Exemple d'ouverture : \all{— Guten Abend, Herr Ngono. Ich muss mit Ihnen sprechen. Seit einer Woche...} / \all{— Oh, das tut mir leid! Ich wusste das nicht. Ich mache die Musik leiser, versprochen!}
\end{experience}

\courssub{Grille d'auto-évaluation}

Après votre production, évaluez-vous honnêtement avec cette grille. Chaque critère validé est un point gagné.

\begin{center}
\begin{tabularx}{\linewidth}{|X|c|c|}
\hline
\rowcolor{popblueL} Critère & Oui & Non \\
\hline
J'ai utilisé au moins 2 mots du vocabulaire du thème (\all{der Lärm, laut, stören, die Ruhe}) & & \\
\hline
Ma subordonnée avec \all{wenn} ou \all{als} est correcte (verbe à la fin, bon choix de conjonction) & & \\
\hline
Ma proposition relative est correcte (pronom relatif au bon cas, verbe à la fin) & & \\
\hline
J'ai employé au moins deux prétérits corrects (\all{war, hatte, konnte, musste}) & & \\
\hline
Ma demande est polie (\all{Sie}, \all{Ich bitte Sie...}) & & \\
\hline
\end{tabularx}
\end{center}

\begin{savaistu}[Et au Bac ?]
Au baccalauréat, la production écrite est évaluée sur trois critères : la \cle{richesse du vocabulaire}, la \cle{correction grammaticale} et la \cle{cohérence} du texte. Bonne nouvelle : réutiliser les structures apprises dans la leçon (subordonnées \all{wenn/als}, relatives, prétérits, connecteurs \all{zuerst, dann, deshalb}) garantit des points sur les trois critères à la fois. Un texte simple mais correct vaut toujours mieux qu'un texte ambitieux mais fautif !
\end{savaistu}

\begin{aretenir}[L'essentiel de la section]
\begin{itemize}
\item Un commentaire se construit en 4 étapes : présentation, résumé, analyse, avis personnel.
\item Les connecteurs \all{zuerst, dann, außerdem, deshalb, meiner Meinung nach} organisent le propos et imposent l'inversion sujet-verbe.
\item Une plainte polie suit le schéma : salutation, constat, conséquence, demande (\all{Ich bitte Sie, ... zu...}), formule de politesse.
\item Jamais d'impératif sec avec un inconnu : toujours \all{Sie} (avec majuscule) et des formules atténuées.
\item Au passé, « quand » se traduit par \all{als} pour un événement unique et par \all{wenn} pour une habitude.
\item L'auto-évaluation systématique transforme chaque production en progrès mesurable.
\end{itemize}
\end{aretenir}

\newpage

\coursec{Activité d'intégration}

\courssub{Objectifs de l'activité}

Cette activité d'intégration te permet de réinvestir, dans une situation de communication réelle, tout ce que tu as appris dans la leçon « Les bruits et les nuisances sonores ». À la fin de l'activité, tu seras capable de :

\begin{itemize}
\item rédiger une \cle{lettre de plainte polie} (\all{eine Beschwerde}) de 8 à 10 lignes en allemand correct ;
\item employer le \cle{vocabulaire du bruit et de la tranquillité} : \all{der Lärm}, \all{laut}, \all{leise}, \all{die Nachbarn}, \all{stören}, \all{die Ruhe}, \all{die Musik}, \all{schlafen}, \all{sich konzentrieren} ;
\item utiliser correctement les \cle{subordonnées temporelles} avec \all{wenn} et \all{als} (verbe à la fin) ;
\item construire une \cle{proposition relative simple} (\all{die Musik, die ...}) ;
\item conjuguer au \cle{prétérit} les verbes \all{sein}, \all{haben} et les verbes modaux (\all{war}, \all{hatte}, \all{konnte}, \all{musste}, \all{wollte}) ;
\item formuler une \cle{demande polie} avec le conditionnel (\all{Ich würde Sie bitten, ...} / \all{Könnten Sie bitte ... ?}) ;
\item rédiger une courte \cle{réponse d'excuse} et présenter l'échange à l'oral en binôme.
\end{itemize}

\courssub{Situation}

Tu t'appelles Amina (ou Arnaud) et tu es élève en classe de Première. Depuis septembre, tu vis dans un petit immeuble à Bonn, en Allemagne, où tu suis un programme d'échange scolaire. Depuis une semaine, ton voisin du deuxième étage, Herr Keller, écoute de la musique très forte chaque nuit, souvent jusqu'à deux heures du matin. Tu ne dors plus, tu es fatigué(e) en classe et tu as un contrôle d'allemand dans trois jours. Tu décides d'écrire une lettre de plainte polie à ton voisin.

\begin{textebox}[Affiche dans l'entrée de l'immeuble]
„Liebe Mieterinnen und Mieter,\\
bitte denken Sie an die Ruhe im Haus!\\
Nach 22 Uhr: keine laute Musik, keine lauten Partys.\\
Bei Problemen: sprechen Sie zuerst mit Ihren Nachbarn oder schreiben Sie an die Hausverwaltung.\\
Vielen Dank!\\
Ihre Hausverwaltung Bonn-Süd“

\emph{Traduction : « Chers locataires, pensez à la tranquillité de l'immeuble ! Après 22 heures : pas de musique forte, pas de fêtes bruyantes. En cas de problème : parlez d'abord à vos voisins ou écrivez à l'administration de l'immeuble. Merci ! »}
\end{textebox}

\begin{definition}[Vocabulaire utile pour la situation]
\all{der Mieter, -} : le locataire ; \all{die Hausverwaltung, -en} : l'administration de l'immeuble ; \all{die Ruhe} : le calme, la tranquillité ; \all{die Beschwerde, -n} : la plainte ; \all{sich beschweren (über + A)} : se plaindre (de) ; \all{der Schlaf} : le sommeil ; \all{müde} : fatigué.
\end{definition}

\courssub{Consignes}

\begin{enumerate}
\item \textbf{Lecture et repérage.} Lis l'affiche de la Hausverwaltung. Souligne les mots qui concernent le bruit et le calme. Réponds en français : à partir de quelle heure doit-on être silencieux ? Que propose l'administration en cas de problème ?
\item \textbf{Préparation du lexique.} Complète mentalement la liste : cinq mots ou expressions de la leçon que tu veux absolument utiliser dans ta lettre (par exemple \all{der Lärm}, \all{stören}, \all{die Ruhe}, \all{laut}, \all{schlafen}).
\item \textbf{Plan de la lettre.} Ta lettre doit contenir quatre parties : (a) la présentation du problème avec \textbf{une proposition relative} (\all{die Musik, die ...}) ; (b) le récit de ce qui s'est passé au \textbf{prétérit} (\all{war}, \all{hatte}, \all{konnte}, \all{musste}) ; (c) les conséquences sur ta santé et tes études avec \textbf{une subordonnée en} \all{wenn} ; (d) une \textbf{demande polie} et une formule de fin.
\item \textbf{Rédaction.} Écris ta lettre de 8 à 10 lignes. Respecte la structure d'une lettre allemande : lieu et date, formule d'appel (\all{Sehr geehrter Herr Keller,}), virgule puis \textbf{minuscule} à la ligne suivante, formule de politesse finale (\all{Mit freundlichen Grüßen}).
\item \textbf{Relecture.} Vérifie : verbe conjugué en deuxième position dans les principales, verbe à la fin dans les subordonnées, majuscule à tous les noms, prétérit correct.
\item \textbf{Échange.} Donne ta lettre à un camarade. Celui-ci joue le rôle de Herr Keller et rédige une courte réponse d'excuse (4 à 5 lignes) : il s'excuse (\all{Es tut mir leid}), explique brièvement et promet de faire attention.
\item \textbf{Oral.} En binôme, présentez le dialogue à la classe : l'un lit sa plainte, l'autre sa réponse.
\end{enumerate}

\courssub{Ressources à mobiliser}

\begin{aretenir}[Rappels grammaticaux]
\begin{itemize}
\item \textbf{Subordonnée temporelle} : le verbe conjugué va \textbf{à la fin}. \all{Wenn ich schlafen will, höre ich die Musik.} — \all{Als ich gestern nach Hause kam, war es sehr laut.} Rappel : \all{wenn} pour le présent, le futur et les actions répétées (aussi au passé) ; \all{als} pour une action unique accomplie au passé.
\item \textbf{Proposition relative} : le pronom relatif reprend le genre et le nombre du nom, son cas dépend de sa fonction ; verbe à la fin. \all{Die Musik, die mein Nachbar hört, ist zu laut.}
\item \textbf{Prétérit} : \all{sein} → \all{ich war} ; \all{haben} → \all{ich hatte} ; \all{können} → \all{ich konnte} ; \all{müssen} → \all{ich musste} ; \all{wollen} → \all{ich wollte}.
\item \textbf{Demande polie} : \all{Ich würde Sie bitten, die Musik leiser zu machen.} / \all{Könnten Sie bitte nach 22 Uhr leise sein ?}
\end{itemize}
\end{aretenir}

\begin{center}
\begin{tabularx}{\linewidth}{|X|X|X|}\hline
\rowcolor{popblueL} Français & Deutsch & Remarque \\ \hline
le bruit & der Lärm (Sg.) & sans pluriel en général \\ \hline
fort, bruyant & laut & contraire : \all{leise} \\ \hline
déranger & stören & verbe faible régulier \\ \hline
le voisin / la voisine & der Nachbar, -n / die Nachbarin, -nen & attention au genre (n-déclinaison) \\ \hline
dormir & schlafen (schläft, schlief) & verbe fort \\ \hline
se concentrer & sich konzentrieren & verbe réfléchi \\ \hline
se plaindre de & sich beschweren über + A & \all{Ich beschwere mich über den Lärm.} \\ \hline
je suis désolé & Es tut mir leid & formule d'excuse standard \\ \hline
\end{tabularx}
\end{center}

\begin{attention}[Pièges fréquents des francophones]
\begin{itemize}
\item Ne dis pas \all{*Ich bin einverstanden} pour « je suis d'accord » dans ce contexte ; pour t'excuser, dis \all{Es tut mir leid} (et non \all{*Ich bin traurig}, qui signifie « je suis triste »).
\item Après la formule d'appel suivie d'une virgule, la ligne suivante commence par une \textbf{minuscule} : \all{Sehr geehrter Herr Keller,} / \all{ich schreibe Ihnen, weil ...}
\item \all{bekommen} ne signifie pas « devenir » mais « recevoir » : \all{Ich habe Ihren Brief bekommen.}
\item Dans la subordonnée, n'oublie pas le verbe à la fin : \all{Wenn ich nachts nicht schlafen kann, ...} et non \all{*Wenn ich kann nicht schlafen}.
\end{itemize}
\end{attention}

\courssub{Proposition de production et corrigé commenté}

\begin{exoresolu}[Lettre de plainte — modèle et commentaire]
\textbf{Énoncé.} Rédige la lettre d'Amina à Herr Keller (8 à 10 lignes) en respectant le plan en quatre parties.

\tcblower

\textbf{Production modèle :}

Bonn, den 12. März\\
Sehr geehrter Herr Keller,\\
ich heiße Amina und ich wohne seit September in der Wohnung unter Ihnen. Ich schreibe Ihnen, weil es seit einer Woche ein Problem gibt: die Musik, die Sie jede Nacht hören, ist sehr laut.\\
Als ich am Montag um zwei Uhr nachts noch wach war, konnte ich nicht schlafen. Am Dienstag hatte ich große Kopfschmerzen, und ich musste in der Schule eine Arbeit schreiben. Wenn ich nachts nicht schlafen kann, bin ich am nächsten Tag müde und ich kann mich nicht konzentrieren. In drei Tagen habe ich eine wichtige Prüfung.\\
Ich würde Sie bitten, nach 22 Uhr keine laute Musik mehr zu hören. Das ist auch die Regel der Hausverwaltung.\\
Ich danke Ihnen und hoffe auf Ihr Verständnis.\\
Mit freundlichen Grüßen\\
Amina

\emph{Traduction : « Monsieur, je m'appelle Amina et j'habite depuis septembre dans l'appartement en dessous du vôtre. Je vous écris parce qu'il y a un problème depuis une semaine : la musique que vous écoutez chaque nuit est très forte. Lundi, quand j'étais encore réveillée à deux heures du matin, je n'ai pas pu dormir. Mardi, j'avais de gros maux de tête et j'ai dû rédiger un devoir à l'école. Quand je ne peux pas dormir la nuit, je suis fatiguée le lendemain et je ne peux pas me concentrer. Dans trois jours, j'ai un examen important. Je vous demanderais de ne plus écouter de musique forte après 22 heures. C'est aussi la règle de l'administration de l'immeuble. Je vous remercie et espère votre compréhension. Cordialement, Amina. »}

\textbf{Commentaire des choix de langue :}
\begin{itemize}
\item \textbf{Proposition relative} : \all{die Musik, die Sie jede Nacht hören, ist sehr laut} — le relatif \all{die} reprend \all{die Musik} (féminin singulier, accusatif car complément de \all{hören}) et le verbe \all{hören} est placé à la fin.
\item \textbf{Subordonnée avec} \all{als} : \all{Als ich am Montag ... noch wach war, ...} — fait unique au passé, donc \all{als} ; verbe \all{war} en fin de subordonnée, et la principale commence par le verbe (\all{konnte}) car la subordonnée occupe la première position.
\item \textbf{Subordonnée avec} \all{wenn} : \all{Wenn ich nachts nicht schlafen kann, bin ich ... müde} — fait répété, vérité générale, donc \all{wenn} ; inversion sujet-verbe dans la principale (\all{bin ich}).
\item \textbf{Prétérit} : \all{war} (sein), \all{hatte} (haben), \all{konnte} (können), \all{musste} (müssen) — exactement les formules exigées par la leçon.
\item \textbf{Demande polie} : \all{Ich würde Sie bitten, ... zu hören} — conditionnel de politesse suivi d'un infinitif avec \all{zu}.
\item \textbf{Registre} : vouvoiement systématique (\all{Sie}, \all{Ihnen}, \all{Ihr}) car on s'adresse à un voisin adulte ; ton ferme mais courtois, sans insulte.
\end{itemize}

\textbf{Réponse d'excuse modèle (rôle de Herr Keller) :}

Liebe Amina,\\
vielen Dank für Ihren Brief. Es tut mir sehr leid! Ich wusste nicht, dass die Musik so laut war. Letzte Woche hatte ich Besuch von Freunden, und wir haben oft Musik gehört. Ich verspreche Ihnen: nach 22 Uhr mache ich die Musik jetzt leise. Viel Glück für Ihre Prüfung!\\
Mit freundlichen Grüßen\\
Herr Keller

\emph{Traduction : « Chère Amina, merci beaucoup pour votre lettre. Je suis vraiment désolé ! Je ne savais pas que la musique était si forte. La semaine dernière, j'avais la visite d'amis et nous avons souvent écouté de la musique. Je vous promets : après 22 heures, je baisse maintenant la musique. Bonne chance pour votre examen ! Cordialement, Herr Keller. »}
\end{exoresolu}

\courssub{Bilan de l'activité}

À travers cette lettre de plainte, tu as mobilisé l'ensemble des compétences de la leçon dans une tâche authentique de la vie quotidienne dans un pays germanophone. Vérifie que tu as bien atteint chaque objectif :

\begin{itemize}
\item \textbf{Lexique} : as-tu employé au moins cinq mots du thème du bruit (\all{der Lärm}, \all{laut}, \all{leise}, \all{stören}, \all{die Ruhe}, \all{die Musik}, \all{schlafen}) ?
\item \textbf{Grammaire} : ta lettre contient-elle au moins une subordonnée avec \all{wenn}, une avec \all{als} et une proposition relative, toutes avec le verbe à la fin ?
\item \textbf{Conjugaison} : as-tu utilisé correctement au moins trois prétérits parmi \all{war}, \all{hatte}, \all{konnte}, \all{musste}, \all{wollte} ?
\item \textbf{Communication} : ta demande est-elle polie (\all{Ich würde Sie bitten ...}) et ta lettre respecte-t-elle la présentation allemande (formule d'appel, minuscule après la virgule, \all{Mit freundlichen Grüßen}) ?
\item \textbf{Interaction} : la réponse d'excuse de ton camarade est-elle cohérente avec ta plainte ?
\end{itemize}

\begin{experience}[Pour aller plus loin : le débat de voisinage]
En groupes de quatre, jouez une réunion de la Hausverwaltung : deux voisins plaignants, Herr Keller et un médiateur. Chacun doit utiliser au moins deux subordonnées (\all{wenn} ou \all{als}) et une relative. Le médiateur conclut avec une phrase au conditionnel : \all{Ich würde vorschlagen, dass ...}
\end{experience}

Cette activité t'a aussi appris une compétence de vie : dans les pays germanophones, la tranquillité après 22 heures (\all{die Nachtruhe}) est une règle sociale très respectée, et savoir se plaindre poliment par écrit est une démarche courante et appréciée.

\newpage

\coursec{Méthodes et exercices résolus}

\coursec{Leçon ALA13 : Les bruits et les nuisances sonores}

\courssub{Texte support : Lärm in der Stadt}

\begin{textebox}[Lärm in der Stadt — Texte original pour la compréhension]
In Yaoundé, wie in allen Großstädten, gehört der Lärm zum Alltag. Herr Abena wohnt im Viertel Mvog-Ada, in der Nähe einer stark befahrenen Kreuzung. Schon um 5 Uhr morgens hört man den Lärm der Motoren von Bendskins und Taxis. Gegen 19 Uhr macht die Bar „Chez Papa“ die Musik sehr laut. Die Gäste schreien, lachen und singen. Herr Abena kann nicht schlafen. Er hat Kopfschmerzen. Er hat schon mit dem Barbesitzer gesprochen, aber der hat gesagt: „Das ist mein Geschäft!“. Herr Abena denkt daran, umzuziehen. Er sucht eine ruhige Wohnung, weit weg von der Straße und den Bars. Ruhe ist wichtig für die Gesundheit.
\end{textebox}

\begin{definition}[Traduction et glossaire des mots difficiles]
\textbf{Traduction :} « À Yaoundé, comme dans toutes les grandes villes, le bruit fait partie du quotidien. M. Abena habite dans le quartier Mvog-Ada, près d'un carrefour très fréquenté. Dès 5 heures du matin, on entend le bruit des moteurs des bendskins et des taxis. Vers 19 heures, le bar „Chez Papa“ met la musique très forte. Les clients crient, rient et chantent. M. Abena ne peut pas dormir. Il a mal à la tête. Il a déjà parlé au patron du bar, mais celui-ci a dit : „C'est mon commerce !“. M. Abena pense à déménager. Il cherche un logement calme, loin de la route et des bars. Le calme est important pour la santé. »

\textbf{Glossaire (nom + article + pluriel) :}
\begin{itemize}
\item \all{der Lärm} (kein Plural / die Lärme selten) — le bruit, le vacarme
\item \all{die Ruhe} (kein Plural) — le calme, le repos
\item \all{der Nachbar, -n} / \all{die Nachbarin, -nen} — le voisin / la voisine
\item \all{die Störung, -en} — la nuisance, le dérangement
\item \all{das Geräusch, -e} — le bruit (son), le son
\item \all{die Bar, -s} — le bar
\item \all{der Besitzer, -} — le propriétaire
\item \all{die Kreuzung, -en} — le carrefour
\item \all{der Bendskin, -s} — le bendskin (moto-taxi)
\item \all{der Kopfschmerz, -en} (meist Pl. \all{die Kopfschmerzen}) — le mal de tête
\item \all{umziehen} (trennbar) — déménager
\item \all{laut / leise} — fort, bruyant / doux, bas, silencieux
\item \all{stören} — déranger
\item \all{sich beschweren über + Akk.} — se plaindre de
\item \all{Lärm machen} — faire du bruit
\end{itemize}
\end{definition}

\courssub{Compréhension du texte}

\begin{methode}[Lire et exploiter un texte de compréhension écrite]
\begin{enumerate}
\item \textbf{Première lecture globale} : lis le texte une fois sans dictionnaire. Repère le thème (ici : \all{der Lärm}), le lieu (\all{in Yaoundé, im Viertel Mvog-Ada}), les personnes (\all{Herr Abena, der Barbesitzer}) et le problème exposé.
\item \textbf{Repérage des mots-clés} : souligne le vocabulaire du bruit : \all{der Lärm, laut, leise, stören, die Ruhe, die Nachbarn, die Bar, der Besitzer}. Devine le sens des mots inconnus grâce au contexte et aux familles de mots (\all{laut} $\rightarrow$ \all{der Lärm} ; \all{stören} $\rightarrow$ \all{die Störung} ; \all{ruhig} $\rightarrow$ \all{die Ruhe}).
\item \textbf{Lecture détaillée} : réponds aux questions dans l'ordre du texte. Les réponses se trouvent presque toujours dans l'ordre des paragraphes.
\item \textbf{Rédaction de la réponse} : réponds par une phrase complète en allemand, en reprenant les mots de la question, sans recopier tout le texte. Exemple : \all{Warum kann Herr Abena nicht schlafen?} $\rightarrow$ \all{Herr Abena kann nicht schlafen, weil die Bar sehr laute Musik macht.}
\item \textbf{Vérification} : verbe en deuxième position dans la principale, verbe à la fin dans la subordonnée (\all{weil, dass, als, wenn}), majuscules à tous les noms, cas corrects (accusatif après \all{wegen}, datif après \all{mit}, etc.).
\end{enumerate}
\end{methode}

\begin{exoresolu}[Compréhension écrite — type Bac]
\textbf{Énoncé.} Lis le texte \emph{Lärm in der Stadt} et réponds aux questions par des phrases complètes en allemand.\par
\textbf{Questions :}
1) \all{Wo wohnt Herr Abena?}
2) \all{Was hört man morgens um 5 Uhr?}
3) \all{Warum hat Herr Abena Kopfschmerzen?}
4) \all{Was hat der Barbesitzer geantwortet?}
5) \all{Was will Herr Abena jetzt machen?}
\tcblower
\textbf{Solution détaillée.}\par
\textbf{1)} Question sur le lieu (\all{wo}). Texte : \all{wohnt im Viertel Mvog-Ada}. Réponse : \all{Herr Abena wohnt im Viertel Mvog-Ada, in der Nähe einer stark befahrenen Kreuzung.} « M. Abena habite dans le quartier Mvog-Ada, près d'un carrefour très fréquenté. »\par
\textbf{2)} Question sur l'objet (\all{was}). Texte : \all{hört man den Lärm der Motoren}. Réponse : \all{Morgens um 5 Uhr hört man den Lärm der Motoren von Bendskins und Taxis.} « Le matin à 5h, on entend le bruit des moteurs de bendskins et de taxis. »\par
\textbf{3)} Question sur la cause (\all{warum}). Texte : \all{die Musik sehr laut ... kann nicht schlafen ... hat Kopfschmerzen}. Réponse avec \all{weil} (verbe à la fin) : \all{Herr Abena hat Kopfschmerzen, weil die Musik in der Bar so laut ist und er nicht schlafen kann.} « M. Abena a mal à la tête parce que la musique au bar est si forte et qu'il ne peut pas dormir. »\par
\textbf{4)} Question sur le contenu d'une parole (\all{was}). Texte : \all{„Das ist mein Geschäft!“}. Discours indirect ou citation : \all{Der Barbesitzer hat gesagt, dass das sein Geschäft sei.} (Konjunktiv I \all{sei} pour le discours indirect formel) ou citation directe : \all{Der Barbesitzer hat geantwortet: „Das ist mein Geschäft!“.} « Le patron du bar a répondu : „C'est mon commerce !“ »\par
\textbf{5)} Question sur l'intention future (\all{was ... machen}). Texte : \all{denkt daran, umzuziehen ... sucht eine ruhige Wohnung}. Réponse : \all{Herr Abena will umziehen und sucht eine ruhige Wohnung, weit weg von der Straße und den Bars.} « M. Abena veut déménager et cherche un logement calme, loin de la route et des bars. »\par
\textbf{Conclusion :} chaque réponse reprend le sujet et le verbe de la question, utilise la conjonction adaptée (\all{weil, dass}) avec le verbe à la fin, et respecte la majuscule des noms.
\end{exoresolu}

\begin{exercice}[Compréhension écrite — Entraînement]
\textbf{Énoncé.} Réponds aux questions suivantes sur le texte \emph{Lärm in der Stadt} par des phrases complètes en allemand.\par
1) \all{Wie heißt der Barbesitzer im Text?} (Piège : le texte ne donne pas son nom propre)\par
2) \all{Was machen die Gäste in der Bar?}\par
3) \all{Wann macht die Bar die Musik laut?}\par
4) \all{Was ist Herrn Abena wichtig für die Gesundheit?}\par
5) \all{Warum ist der Lärm in Yaoundé ein Problem für die Anwohner?} (Réponse libre basée sur le texte)
\end{exercice}

\begin{corrige}[Exercice 1]
1) \all{Der Text nennt keinen Namen für den Barbesitzer. Er wird nur „der Barbesitzer“ oder „der Besitzer“ genannt.}\par
2) \all{Die Gäste schreien, lachen und singen.}\par
3) \all{Die Bar macht die Musik gegen 19 Uhr laut.}\par
4) \all{Ruhe ist Herrn Abena wichtig für die Gesundheit.}\par
5) \all{Der Lärm ist ein Problem, weil die Anwohner (wie Herr Abena) nicht schlafen können, Kopfschmerzen bekommen und ihre Ruhe brauchen.}
\end{corrige}

\courssub{Lexique thématique : der Lärm und die Ruhe}

\begin{propriete}[Vocabulaire de base : noms, verbes, adjectifs et expressions]
\begin{center}
\begin{tabularx}{\linewidth}{|X|X|X|}
\hline
\rowcolor{popblueL} \textbf{Deutsch (Article, Pluriel)} & \textbf{Français} & \textbf{Remarque / Construction} \\
\hline
\all{der Lärm} (kein Pl.) & le bruit, le vacarme & Souvent sans pluriel. \all{der Lärmpegel} (niveau sonore). \\
\all{die Ruhe} (kein Pl.) & le calme, le repos & \all{Ruhe haben / brauchen / suchen}. \\
\all{der Nachbar, -n} / \all{die Nachbarin, -nen} & le voisin / la voisine & \textbf{n-Déclinaison} : \all{dem Nachbarn, den Nachbarn} (tous cas sauf Nom. sg). \\
\all{die Störung, -en} & la nuisance, le dérangement & De \all{stören}. \all{eine Störung melden}. \\
\all{das Geräusch, -e} & le bruit (son), le son & \all{ein lautes / leises Geräusch}. \\
\all{die Lärmbelästigung, -en} & la pollution sonore & Terme administratif / juridique. \\
\all{die Nachtruhe} (kein Pl.) & le calme nocturne & Protégée par la loi (\all{Nachtruhezeit} 22h--6h). \\
\hline
\all{stören} (regelmäßig) & déranger & \textbf{+ Akkusativ} : \all{Der Lärm stört mich.} \\
\all{sich beschweren über + Akk.} & se plaindre de & Réfléchi + préposition : \all{Ich beschwere mich über den Lärm.} \\
\all{Lärm machen} & faire du bruit & \all{Die Kinder machen Lärm.} (Pas \all{*Lärm tun}). \\
\all{ruhig sein / werden} & être / devenir calme & Adjectif \all{ruhig} $\neq$ nom \all{die Ruhe}. \\
\all{umziehen} (trennbar) & déménager & \all{Er zieht um.} / \all{Er ist umgezogen} (Perfekt avec \all{sein}). \\
\hline
\all{laut} & fort, bruyant & Antonyme : \all{leise}. Steigerung : \all{laut -- lauter -- am lautesten}. \\
\all{leise} & doux, bas, silencieux & \all{leise sprechen / sein}. \\
\all{ohrenbetäubend} & assourdissant & Composé : \all{Ohren} + \all{betäuben}. \\
\hline
\end{tabularx}
\end{center}
\end{propriete}

\begin{methode}[Maîtriser le lexique thématique]
\begin{enumerate}
\item \textbf{Apprends chaque nom avec son article et son pluriel} : note la n-déclinaison de \all{der Nachbar} (\all{den Nachbarn, dem Nachbarn, des Nachbarn}) et l'absence de pluriel pour \all{der Lärm} et \all{die Ruhe}.
\item \textbf{Associe les contraires} : \all{laut} $\leftrightarrow$ \all{leise} ; \all{der Lärm} $\leftrightarrow$ \all{die Ruhe} ; \all{stören} $\leftrightarrow$ \all{nicht stören / ruhig sein}.
\item \textbf{Mémorise les verbes et leurs constructions (Regierung)} :
      \begin{itemize}
      \item \all{stören + Akkusativ} (\all{Die Musik stört die Nachbarn.})
      \item \all{sich beschweren über + Akkusativ} (\all{Wir beschweren uns über den Lärm.})
      \item \all{Lärm machen} (collocation fixe)
      \item \all{umziehen} (verbe à particule séparable, auxiliaire \all{sein} au Perfekt : \all{ist umgezogen})
      \end{itemize}
\item \textbf{Utilise les familles de mots} pour enrichir l'expression :
      \all{laut} (adj.) $\rightarrow$ \all{der Lärm} (nom) $\rightarrow$ \all{lärmen} (verbe, fam. faire du vacarme) $\rightarrow$ \all{lärmend} (part. prés.).
\item \textbf{Réemploie le lexique dans tes propres phrases} : c'est la condition pour qu'il soit disponible à l'écrit comme à l'oral.
\end{enumerate}
\end{methode}

\begin{exoresolu}[Lexique — compléter et réemployer]
\textbf{Énoncé.} Complète les phrases avec le mot qui convient (forme correcte) : \all{der Lärm, leise, stört, die Ruhe, die Nachbarn, die Störung, laut, ruhige}.\par
1) In der Nacht brauche ich \dots\ zum Schlafen.\par
2) \dots\ in der Stadt kommt oft von den Autos und den Bars.\par
3) Bitte sprich \dots, das Baby schläft!\par
4) Die \dots\ Musik \dots\ die ganze Familie.\par
5) \dots\ haben sich über die Party bei der Polizei beschwert.\par
6) Er sucht eine \dots\ Wohnung ohne \dots\.
\tcblower
\textbf{Solution détaillée.}\par
1) \all{In der Nacht brauche ich \textbf{Ruhe} zum Schlafen.} (Akkusatif de \all{die Ruhe} $\rightarrow$ \all{Ruhe}).\par
2) \all{\textbf{Der Lärm} in der Stadt kommt oft von den Autos und den Bars.} (Sujet Nominatif masculin \all{der Lärm}).\par
3) \all{Bitte sprich \textbf{leise}, das Baby schläft!} (Adverbe / adjectif prédicatif).\par
4) \all{Die \textbf{laute} Musik \textbf{stört} die ganze Familie.} (Adj. \all{laute} accordé fém. nom/akk après art. déf. ; verbe \all{stören} 3ème pers. sg).\par
5) \all{\textbf{Die Nachbarn} haben sich über die Party bei der Polizei \textbf{beschwert}.} (Sujet Nominatif pluriel ; Participe passé \all{beschwert} avec \all{haben} car pas de mouvement).\par
6) \all{Er sucht eine \textbf{ruhige} Wohnung ohne \textbf{Lärm}.} (Adj. \all{ruhige} fém. akk après art. indéf. ; \all{ohne + Akk} $\rightarrow$ \all{Lärm}).\par
\textbf{Conclusion :} toujours vérifier le cas demandé par la fonction du mot (sujet $\rightarrow$ Nominatif, COD $\rightarrow$ Akkusatif, après préposition $\rightarrow$ cas imposé) et l'accord de l'adjectif.
\end{exoresolu}

\begin{exercice}[Lexique — Entraînement]
\textbf{Énoncé.} Traduis en allemand en utilisant le vocabulaire de la leçon.\par
1) Le bruit des voitures me dérange. (stören)\par
2) Les voisins se plaignent du vacarme. (sich beschweren über)\par
3) Nous avons besoin de calme pour apprendre. (brauchen, Ruhe)\par
4) Fais moins de bruit, s'il te plaît ! (Lärm machen, bitte)\par
5) C'est une zone résidentielle calme. (ruhig, Wohngegend)\par
6) La pollution sonore est un grand problème dans les grandes villes. (Lärmbelästigung, Großstadt)
\end{exercice}

\begin{corrige}[Exercice 2]
1) \all{Der Lärm der Autos stört mich.}\par
2) \all{Die Nachbarn beschweren sich über den Lärm.}\par
3) \all{Wir brauchen Ruhe zum Lernen.}\par
4) \all{Mach bitte weniger Lärm!} (ou \all{Mach nicht so viel Lärm!})\par
5) \all{Das ist eine ruhige Wohngegend.}\par
6) \all{Die Lärmbelästigung ist ein großes Problem in Großstädten.}
\end{corrige}

\courssub{Grammaire 1 : Les subordonnées temporelles (wenn, als)}

\begin{propriete}[Règle de choix : \all{wenn} ou \all{als} ?]
\begin{center}
\begin{tabularx}{\linewidth}{|X|X|X|}
\hline
\rowcolor{popblueL} \textbf{Conjonction} & \textbf{Emploi} & \textbf{Exemple} \\
\hline
\all{als} & \textbf{Passé} (\all{Präteritum} ou \all{Perfekt}) + \textbf{Événement unique / Ponctuel} (une seule fois, daté). & \all{Als ich gestern ankam, war es still.} \\
\all{wenn} & \textbf{Présent / Futur} (régularité, condition) & \all{Wenn es regnet, bleibe ich zu Hause.} \\
\all{wenn} & \textbf{Passé} + \textbf{Événement répété / Habitude} (souvent avec \all{früher, immer, oft, jedes Mal}). & \all{Früher, wenn ich Zeit hatte, las ich viel.} \\
\hline
\end{tabularx}
\end{center}
\textbf{Structure commune :} \all{, [Konjunktion] Subjekt ... Verb\_conjugué.} Virgule obligatoire. Verbe conjugué \textbf{à la fin} de la subordonnée. Si la subordonnée est en tête : \textbf{Inversion} dans la principale (Verbe en position 1).
\end{propriete}

% TikZ Decision Tree für wenn/als
\begin{popfigure}
\begin{tikzpicture}[
    node distance=1.2cm and 1.5cm,
    decision/.style={diamond, draw=popblue, thick, fill=popblueL, text width=2.8cm, align=center, inner sep=2pt},
    answer/.style={rectangle, draw=popgreen, thick, fill=popgreenL, text width=2.5cm, align=center, inner sep=2pt},
    start/.style={ellipse, draw=poporange, thick, fill=poporangeL, text width=2cm, align=center},
    arrow/.style={-Stealth, thick, draw=popdark}
]
\node[start, align=center] (start){Subordonnée\\temporelle};
\node[decision, below of=start, align=center] (past){Le verbe de la\\principale est-il au\\\textbf{Passé} ?};
\node[answer, below left=1.5cm and -0.5cm of past, align=center] (present){\all{wenn}\\(Présent / Futur)};
\node[decision, below right=1.5cm and -0.5cm of past, align=center] (unique){Action \textbf{unique}\\ponctuelle ?\\(\all{gestern, letzte Woche, 2010})};
\node[answer, below left=1.5cm and -0.5cm of unique, align=center] (wenn_past){\all{wenn}\\(Habitude passée:\\\all{früher, immer, oft})};
\node[answer, below right=1.5cm and -0.5cm of unique, align=center] (als_past){\all{als}\\(Événement unique\\passé)};

\draw[arrow] (start) -- (past);
\draw[arrow] (past) -- node[left, pos=0.5, align=center] {\textbf{NON}} (present);
\draw[arrow] (past) -- node[right, pos=0.5, align=center] {\textbf{OUI}} (unique);
\draw[arrow] (unique) -- node[left, pos=0.5, align=center] {\textbf{NON}} (wenn_past);
\draw[arrow] (unique) -- node[right, pos=0.5, align=center] {\textbf{OUI}} (als_past);
\end{tikzpicture}
\legende{Arbre de décision : choisir entre \all{wenn} et \all{als}.}
\end{popfigure}

\begin{methode}[Choisir entre \all{wenn} et \all{als} — Démarche en 3 étapes]
\begin{enumerate}
\item \textbf{Identifie le temps de la phrase principale} : est-il au Présent/Futur ou au Passé (Präteritum/Perfekt) ?
\item \textbf{Si Présent/Futur} $\rightarrow$ toujours \all{wenn}.
\item \textbf{Si Passé} : demande-toi si l'action de la subordonnée est \textbf{unique} (datée : \all{gestern, vor einem Jahr, als Kind} = période close unique) ou \textbf{répétée} (\all{früher, immer, oft, jedes Mal, wann immer}).
      \begin{itemize}
      \item Unique $\rightarrow$ \all{als}
      \item Répétée $\rightarrow$ \all{wenn}
      \end{itemize}
\item \textbf{Applique la structure} : Virgule + Sujet + ... + Verbe conjugué \textbf{à la fin}. Si la subordonnée commence la phrase : Verbe de la principale en position 1 (inversion).
\end{enumerate}
\end{methode}

\begin{exoresolu}[Grammaire — \all{wenn} ou \all{als} ?]
\textbf{Énoncé.} Complète avec \all{wenn} ou \all{als}, conjugue le verbe entre parenthèses au temps qui convient, puis traduis.\par
1) \all{\dots\ ich klein \dots\ (sein), gab es bei uns keinen Lärm.}\par
2) \all{\dots\ die Nachbarn heute Musik \dots\ (machen), kann ich nicht lernen.}\par
3) \all{Früher, \dots\ mein Vater nach Hause \dots\ (kommen), war es immer still.}\par
4) \all{\dots\ es gestern Abend laut \dots\ (werden), habe ich die Polizei gerufen.}\par
5) \all{Ich mache das Fenster zu, \dots\ es \dots\ (sein) zu laut.}
\tcblower
\textbf{Solution détaillée.}\par
1) \all{\textbf{Als} ich klein \textbf{war}, gab es bei uns keinen Lärm.} (Période unique passée « quand j'étais petit » $\rightarrow$ \all{als} + Präteritum). « Quand j'étais petit, il n'y avait pas de bruit chez nous. »\par
2) \all{\textbf{Wenn} die Nachbarn heute Musik \textbf{machen}, kann ich nicht lernen.} (Présent, situation répétée/conditionnelle $\rightarrow$ \all{wenn}). « Si les voisins font de la musique aujourd'hui, je ne peux pas travailler. »\par
3) \all{Früher, \textbf{wenn} mein Vater nach Hause \textbf{kam}, war es immer still.} (Passé + \all{früher} + habitude répétée $\rightarrow$ \all{wenn} + Präteritum). « Autrefois, quand mon père rentrait, c'était toujours calme. »\par
4) \all{\textbf{Als} es gestern Abend laut \textbf{wurde}, habe ich die Polizei gerufen.} (Passé + \all{gestern Abend} = unique $\rightarrow$ \all{als} + Präteritum \all{wurde}). « Quand ça est devenu bruyant hier soir, j'ai appelé la police. »\par
5) \all{Ich mache das Fenster zu, \textbf{wenn} es \textbf{ist} zu laut.} (Présent/Futur conditionnel $\rightarrow$ \all{wenn}). « Je ferme la fenêtre s'il fait trop de bruit. »\par
\textbf{Conclusion :} le critère décisif est « Passé + Unique = \all{als} » ; tout le reste (Présent, Futur, Passé répété) = \all{wenn}.
\end{exoresolu}

\begin{exercice}[Grammaire — \all{Wenn} / \all{Als}]
\textbf{Énoncé.} Complète par \all{wenn} ou \all{als} et mets le verbe au temps indiqué.\par
1) \all{\dots\ ich in Yaoundé \dots\ (wohnen, Präteritum), war es oft laut.}\par
2) \all{Wir gehen spazieren, \dots\ das Wetter \dots\ (sein, Präsens) gut.}\par
3) \all{\dots\ der Lehrer \dots\ (kommen, Präteritum), wurden alle still.}\par
4) \all{Meine Mutter kocht immer, \dots\ sie \dots\ (haben, Präsens) Zeit.}\par
5) \all{\dots\ wir \dots\ (ankommen, Perfekt) gestern, war die Party schon zu Ende.}
\end{exercice}

\begin{corrige}[Exercice 3]
1) \all{Als} ich in Yaoundé \all{wohnte}, war es oft laut. (Période unique passée vécue comme un bloc).\par
2) \all{Wir gehen spazieren, \textbf{wenn} das Wetter \textbf{ist} gut.}\par
3) \all{Als} der Lehrer \all{kam}, wurden alle still. (Événement unique ponctuel).\par
4) \all{Meine Mutter kocht immer, \textbf{wenn} sie \textbf{hat} Zeit.}\par
5) \all{Als} wir \all{sind angekommen} gestern, war die Party schon zu Ende. (Arrivée unique hier).
\end{corrige}

\courssub{Grammaire 2 : Les propositions relatives simples}

\begin{propriete}[Tableau des pronoms relatifs (Déclinaison)]
Le pronom relatif s'accorde en \textbf{genre} et \textbf{nombre} avec son antécédent. Son \textbf{cas} dépend de sa fonction \textbf{dans la proposition relative}.
\begin{center}
\begin{tabular}{|c|c|c|c|c|}
\hline
\rowcolor{popblueL} \textbf{Cas} & \textbf{Maskulin} & \textbf{Féminin} & \textbf{Neutre} & \textbf{Pluriel} \\
\hline
Nominatif & \all{der} & \all{die} & \all{das} & \all{die} \\
\hline
Akkusativ & \all{den} & \all{die} & \all{das} & \all{die} \\
\hline
Dativ & \all{dem} & \all{der} & \all{dem} & \all{denen} \\
\hline
Genitif & \all{dessen} & \all{deren} & \all{dessen} & \all{deren} \\
\hline
\end{tabular}
\end{center}
\textbf{Règles d'or :}
\begin{enumerate}
\item Place la relative \textbf{juste après l'antécédent}.
\item Mets une \textbf{virgule} avant le pronom relatif.
\item Le verbe conjugué va \textbf{à la fin} de la relative.
\item \textbf{Jamais} de répétition de l'antécédent dans la relative.
\item Pour les prépositions : \all{Präposition + Relativpronomen} (ex. \all{mit dem}, \all{für den}, \all{wegen dessen}).
\end{enumerate}
\end{propriete}

\begin{methode}[Relier deux phrases par un pronom relatif — Démarche en 4 étapes]
\begin{enumerate}
\item \textbf{Trouve l'antécédent} : le nom commun aux deux phrases.
\item \textbf{Détermine son genre et son nombre} (ex. \all{die Bar} = fém. sg).
\item \textbf{Détermine la fonction du pronom DANS la relative} (Sujet $\rightarrow$ Nominatif ; COD $\rightarrow$ Akkusatif ; après préposition/verbe datif $\rightarrow$ Datif).
\item \textbf{Choisis la forme} dans le tableau, place la relative après l'antécédent avec virgule, verbe à la fin.
\end{enumerate}
\end{methode}

\begin{exoresolu}[Grammaire — les pronoms relatifs]
\textbf{Énoncé.} Relie les deux phrases par une proposition relative. Indique le cas du pronom.\par
1) \all{Der Nachbar macht viel Lärm. Der Nachbar wohnt neben uns.}\par
2) \all{Das ist ein Geräusch. Ich kenne das Geräusch nicht.}\par
3) \all{Die Musik ist zu laut. Die Musik kommt aus der Bar.}\par
4) \all{Ich spreche mit dem Besitzer. Der Besitzer ist unfreundlich.}\par
5) \all{Die Nachbarn sind wütend. Ich habe mich bei den Nachbarn beschwert.}
\tcblower
\textbf{Solution détaillée.}\par
1) Antécédent : \all{der Nachbar} (Mask. Sg). Dans relative : Sujet $\rightarrow$ \textbf{Nominatif} $\rightarrow$ \all{der}. \all{Der Nachbar, \textbf{der} neben uns wohnt, macht viel Lärm.}\par
2) Antécédent : \all{das Geräusch} (Neutre Sg). Dans relative : COD de \all{kennen} $\rightarrow$ \textbf{Akkusativ} $\rightarrow$ \all{das} (identique au Nom. pour neutre). \all{Das ist ein Geräusch, \textbf{das} ich nicht kenne.}\par
3) Antécédent : \all{die Musik} (Fém. Sg). Dans relative : Sujet $\rightarrow$ \textbf{Nominatif} $\rightarrow$ \all{die}. \all{Die Musik, \textbf{die} aus der Bar kommt, ist zu laut.}\par
4) Antécédent : \all{der Besitzer} (Mask. Sg). Dans relative : Sujet $\rightarrow$ \textbf{Nominatif} $\rightarrow$ \all{der}. \all{Ich spreche mit dem Besitzer, \textbf{der} unfreundlich ist.}\par
5) Antécédent : \all{die Nachbarn} (Pluriel). Dans relative : Objet de la préposition \all{bei} (Datif) $\rightarrow$ \textbf{Datif Pluriel} $\rightarrow$ \all{denen}. \all{Die Nachbarn, \textbf{bei denen} ich mich beschwert habe, sind wütend.} (Préposition devant le pronom).\par
\textbf{Conclusion :} Genre/Nombre = Antécédent ; Cas = Fonction dans la relative ; Verbe à la fin.
\end{exoresolu}

\begin{exercice}[Grammaire — Relatives]
\textbf{Énoncé.} Relie par une relative. Attention au cas (Prépositions : \all{mit, bei, wegen, für}).\par
1) \all{Das ist die Bar. Die Bar ist sehr laut.}\par
2) \all{Ich kenne den Mann. Der Mann repariert das Auto.}\par
3) \all{Wir warten auf den Bus. Der Bus hat Verspätung.}\par
4) \all{Das sind die Kinder. Die Kinder machen Lärm.}\par
5) \all{Herr Abena sucht eine Wohnung. Die Wohnung ist ruhig.}
\end{exercice}

\begin{corrige}[Exercice 4]
1) \all{Das ist die Bar, \textbf{die} sehr laut ist.} (Fém. Sg, Sujet $\rightarrow$ Nom. \all{die}).\par
2) \all{Ich kenne den Mann, \textbf{der} das Auto repariert.} (Mask. Sg, Sujet $\rightarrow$ Nom. \all{der}).\par
3) \all{Wir warten auf den Bus, \textbf{der} Verspätung hat.} (Mask. Sg, Sujet $\rightarrow$ Nom. \all{der}). Note : \all{warten auf + Akk} dans la principale, mais Sujet dans la relative.\par
4) \all{Das sind die Kinder, \textbf{die} Lärm machen.} (Pluriel, Sujet $\rightarrow$ Nom. \all{die}).\par
5) \all{Herr Abena sucht eine Wohnung, \textbf{die} ruhig ist.} (Fém. Sg, Sujet $\rightarrow$ Nom. \all{die}).
\end{corrige}

\courssub{Conjugaison : Prétérit de \all{sein}, \all{haben} et des verbes modaux}

\begin{propriete}[Conjugaison au Präteritum (Passé simple narratif)]
Au passé (récit, écrit), on utilise le \textbf{Präteritum} pour \all{sein}, \all{haben} et les \textbf{verbes modaux} (même à l'oral). Pas d'Umlaut aux modaux au Präteritum !
\begin{center}
\begin{tabular}{|l|c|c|c|c|c|c|}
\hline
\rowcolor{popblueL} \textbf{Personne} & \textbf{sein} & \textbf{haben} & \textbf{können} & \textbf{müssen} & \textbf{wollen} & \textbf{sollen} \\
\hline
ich & war & hatte & konnte & musste & wollte & sollte \\
du & warst & hattest & konntest & musstest & wolltest & solltest \\
er/sie/es & war & hatte & konnte & musste & wollte & sollte \\
wir & waren & hatten & konnten & mussten & wollten & sollten \\
ihr & wart & hattet & konntet & musstet & wolltet & solltet \\
sie/Sie & waren & hatten & konnten & mussten & wollten & sollten \\
\hline
\end{tabular}
\end{center}
\textbf{Structure avec modal au Präteritum :} \all{Modal\_conjugué (Pos. 2) ... Infinitif (à la fin)}. Ex. \all{Ich konnte nicht schlafen.}
\textbf{Attention :} \all{könnte, müsste, wollte} (avec Umlaut) = \textbf{Konjunktiv II} (irréel / politesse : « je pourrais / devrais / voudrais »), \textbf{PAS} Präteritum.
\end{propriete}

\begin{methode}[Raconter au passé avec les auxiliaires et les modaux]
\begin{enumerate}
\item \textbf{Apprends les formes par cœur} (tableau ci-dessus). Note l'alternance vocalique pour \all{sein} (\all{ich war, du warst}) et la régularité apparente des modaux (dentale \all{-t-} + terminaisons faibles).
\item \textbf{Identifie le verbe principal} : si c'est \all{sein} ou \all{haben} $\rightarrow$ conjugue directement au Präteritum. Si c'est un modal $\rightarrow$ conjugue le modal au Präteritum + place l'infinitif du verbe lexical \textbf{à la fin} de la phrase.
\item \textbf{Vérifie l'ordre des mots} : \all{Als der Lärm anfing, konnte ich nicht schlafen.} (Subordonnée \all{als} $\rightarrow$ verbe \all{anfing} à la fin ; Principale $\rightarrow$ modal \all{konnte} pos 2, infinitif \all{schlafen} à la fin).
\item \textbf{Génitif après \all{wegen}} : \all{wegen des Lärms} (Mask./Neutre : \all{des + -s}), \all{wegen der Ruhe} (Fém.), \all{wegen der Nachbarn} (Pluriel).
\end{enumerate}
\end{methode}

\begin{exoresolu}[Conjugaison et thème — Mettre au Präteritum]
\textbf{Énoncé.} Mets les phrases au Präteritum. Traduis la phrase 3.\par
1) \all{Die Musik ist zu laut.}\par
2) \all{Wir haben keine Ruhe.}\par
3) \all{Die Kinder können wegen des Lärms nicht schlafen.}\par
4) \all{Ich muss mit dem Besitzer sprechen.}\par
5) \all{Die Bar darf bis 22 Uhr offen bleiben.}
\tcblower
\textbf{Solution détaillée.}\par
1) \all{Die Musik \textbf{war} zu laut.} (\all{sein}, 3ème sg). « La musique était trop forte. »\par
2) \all{Wir \textbf{hatten} keine Ruhe.} (\all{haben}, 1ère pl). « Nous n'avions pas de calme. »\par
3) \all{Die Kinder \textbf{konnten} wegen des Lärms nicht \textbf{schlafen}.} (\all{können} $\rightarrow$ \all{konnten} sans Umlaut ! Infinitif \all{schlafen} à la fin ; \all{wegen} + Génitif \all{des Lärms}). Traduction : « Les enfants ne pouvaient pas dormir à cause du bruit. »\par
4) \all{Ich \textbf{musste} mit dem Besitzer \textbf{sprechen}.} (\all{müssen} $\rightarrow$ \all{musste}, infinitif final). « Je devais parler avec le propriétaire. »\par
5) \all{Die Bar \textbf{durfte} bis 22 Uhr offen \textbf{bleiben}.} (\all{dürfen} $\rightarrow$ \all{durfte}, infinitif \all{bleiben} final). « Le bar avait le droit / était autorisé à rester ouvert jusqu'à 22h. »\par
\textbf{Conclusion :} au Präteritum, le modal (ou l'auxiliaire) porte la marque du passé et de la personne ; l'infinitif lexical reste à la fin, inchangé.
\end{exoresolu}

\begin{exercice}[Conjugaison — Präteritum]
\textbf{Énoncé.} Conjugue le verbe entre parenthèses au Präteritum.\par
1) Gestern \all{(sein)} \dots\ der Lärm unerträglich.\par
2) Wir \all{(haben)} \dots\ große Kopfschmerzen.\par
3) Der Besitzer \all{(können)} \dots\ die Musik nicht leiser \all{(machen)}.\par
4) Ich \all{(müssen)} \dots\ die Polizei \all{(rufen)}.\par
5) Die Nachbarn \all{(wollen)} \dots\ schlafen, aber die Bar \all{(sollen)} \dots\ lauter \all{(sein)}.
\end{exercice}

\begin{corrige}[Exercice 5]
1) Gestern \all{war} der Lärm unerträglich.\par
2) Wir \all{hatten} große Kopfschmerzen.\par
3) Der Besitzer \all{konnte} die Musik nicht leiser \all{machen}.\par
4) Ich \all{musste} die Polizei \all{rufen}.\par
5) Die Nachbarn \all{wollten} schlafen, aber die Bar \all{sollte} lauter \all{sein}. (Note: \all{sollte} = Prétérit de \all{sollen} « devait / était censé » ; \all{sollte} avec Umlaut n'existe pas).
\end{corrige}

\courssub{Expression écrite et orale guidée}

\begin{methode}[Rédiger un court texte guidé — « Le bruit dans mon quartier »]
\begin{enumerate}
\item \textbf{Plan en trois temps} :
      \begin{enumerate}
      \item \textbf{Situation} : \all{Ich wohne in ... / In meiner Straße gibt es ...} (Présent).
      \item \textbf{Problème \& Conséquences} : \all{Der Lärm stört mich, deshalb kann ich nicht ... / Als / Wenn ..., war / ist es laut.} (Intègre \all{weil, denn, deshalb}, vocabulaire thème, \all{wenn/als}, relative).
      \item \textbf{Solution / Avis} : \all{Meiner Meinung nach sollte / müsste man ... / Ich werde ...} (Modal au Prétérit ou Présent + infinitif).
      \end{enumerate}
\item \textbf{Contraintes obligatoires (niveau Bac)} : Minimum 6--8 phrases. Utilise \textbf{au moins} : 4 mots du lexique (\all{Lärm, Ruhe, stören, Nachbar, laut/leise, umziehen}), 1 subordonnée \all{wenn} ou \all{als}, 1 proposition relative, 2 verbes au Prétérit parmi \all{sein, haben, können, müssen, wollen, sollen, dürfen}.
\item \textbf{Connecteurs logiques} : \all{deshalb, aber, denn, weil, außerdem, zuerst, dann, schließlich}.
\item \textbf{Relis-toi (Check-list)} : Verbe position 2 (principale) / fin (subordonnée/relative) ? Majuscules aux noms ? Cas corrects (Akk/Dat/Gén) ? Accords adjectifs ? Particules séparables (\all{umziehen, anrufen, aufmachen}) à la fin ?
\end{enumerate}
\end{methode}

\begin{exoresolu}[Production écrite guidée — Modèle type Bac]
\textbf{Énoncé.} Rédige un texte de 6 à 8 phrases : « Le bruit dans ton quartier ». Contraintes : 1 subordonnée \all{wenn/als}, 1 relative, Prétérit de 2 verbes parmi \all{sein, haben, können, müssen}.\par
\tcblower
\textbf{Modèle de rédaction :}\par
\all{Ich wohne in Douala, in einem Viertel, das sehr lebendig ist. Direkt neben unserem Haus gibt es eine Bar, die jeden Abend laute Musik spielt. Als ich letzte Woche für das Abitur lernen wollte, war der Lärm unerträglich. Ich konnte mich nicht konzentrieren und hatte starke Kopfschmerzen. Meine kleine Schwester konnte nicht schlafen, weil die Musik so laut war. Mein Vater musste mit dem Besitzer sprechen, aber der wollte nichts ändern. Meiner Meinung nach sollten die Bars die Musik nach 22 Uhr leiser machen. Dann hätten alle mehr Ruhe.}\par
\textbf{Traduction :} « J'habite à Douala, dans un quartier très animé. Juste à côté de notre maison, il y a un bar qui joue de la musique forte chaque soir. Quand j'ai voulu réviser pour le Bac la semaine dernière, le bruit était insupportable. Je ne pouvais pas me concentrer et j'avais de forts maux de tête. Ma petite sœur ne pouvait pas dormir parce que la musique était si forte. Mon père a dû parler au propriétaire, mais celui-ci n'a rien voulu changer. À mon avis, les bars devraient baisser la musique après 22 heures. Alors tout le monde aurait plus de calme. »\par
\textbf{Points de grammaire vérifiés :}
\begin{itemize}
\item Relatives : \all{das sehr lebendig ist} (Neutre, Nom., antécédent \all{Viertel}) ; \all{die jeden Abend ... spielt} (Fém., Nom., antécédent \all{Bar}).
\item Subordonnée temporelle : \all{Als ich ... lernen wollte} (Unique passé $\rightarrow$ \all{als}, verbe \all{wollte} à la fin).
\item Subordonnée causale : \all{weil die Musik so laut war} (verbe \all{war} à la fin).
\item Prétérit : \all{war} (\all{sein}), \all{konnte} (\all{können}), \all{hatte} (\all{haben}), \all{musste} (\all{müssen}), \all{wollte} (\all{wollen}), \all{sollten} (\all{sollen} Pluriel Prétérit).
\item Ordre des mots : Verbe 2ème position principales (\all{wohne, gibt, war, konnte, hatte, konnte, musste, wollte, sollten, hätten}) ; Verbe fin subordonnées/relatives.
\item Majuscules : \all{Lärm, Musik, Bar, Besitzer, Vater, Schwester, Kopfschmerzen, Meinung, Bars, Uhr, Ruhe}.
\end{itemize}
\end{exoresolu}

\begin{experience}[Production orale : Jeu de rôle « Chez le voisin / Au bar »]
\textbf{Consigne :} En binôme. L'un joue \all{der Anwohner} (le riverain), l'autre \all{der Barbesitzer} ou \all{der Nachbar}.
\textbf{Scénario A :} Le riverain sonne chez le voisin qui fait une fête bruyante (\all{zu laut, nach 22 Uhr, Ruhe brauchen, sich beschweren}). Le voisin réagit (\all{Entschuldigung, wir machen leiser / Das ist mein Recht}).
\textbf{Scénario B :} Le riverain va voir le patron du bar (\all{Barbesitzer, Lärmbelästigung, Gesetz, Nachtruhe, Ordnungsamt}). Le patron négocie (\all{Wirtschaft, Gäste, Kompromiss, 23 Uhr}).
\textbf{Objectif :} Réemployer le lexique (\all{stören, Lärm, Ruhe, laut, leise, sich beschweren, umziehen, müssen, können, wollen}) et la grammaire (modaux, \all{weil/als/wenn}, relatives) à l'oral.
\end{experience}

\begin{exercice}[Production écrite — Entraînement Bac]
\textbf{Énoncé.} « \all{Lärm in meiner Umgebung} ». Écris un paragraphe d'environ 80 mots. Décris une situation de bruit réel ou imaginé (quartier, école, marché, transport). Intègre obligatoirement :
\begin{itemize}
\item 5 mots du vocabulaire thème (noms, verbes, adjectifs).
\item 1 subordonnée avec \all{als} (passé unique) OU \all{wenn} (habitude/présent).
\item 1 proposition relative (sujet ou objet).
\item 2 verbes au Prétérit (\all{sein, haben, können, müssen, wollen, sollen, dürfen}).
\item Connecteurs : \all{deshalb, weil, aber, außerdem}.
\end{itemize}
\end{exercice}

\begin{corrige}[Exercice 6]
\textbf{Proposition de correction (exemple) :}\par
\all{In meiner Straße in Yaoundé gibt es viele Bars und Restaurants. Der Lärm am Wochenende ist ohrenbetäubend. Wenn die Musik anfängt, kann niemand schlafen. Letzte Woche war es besonders schlimm, als ein DJ bis 3 Uhr morgens auflegte. Ich hatte Kopfschmerzen und konnte mich nicht erholen. Meine Nachbarn haben sich bei der Polizei beschwert, aber die Beamten konnten nichts machen. Meiner Meinung nach muss die Stadt strengere Regeln für die Nachtruhe einführen. Dann hätten wir endlich wieder Ruhe.}\par
\textbf{Vérification :} Lexique (Lärm, Musik, Kopfschmerzen, Nachbarn, Ruhe, Nachtruhe) ; \all{Wenn} (habitude) + \all{als} (unique passé) ; Relative \all{die Beamten konnten nichts machen} (sujet) ou \all{als ein DJ ... auflegte} (subordonnée) ; Prétérit : \all{war, hatte, konnte, konnten, mussten, hatten} ; Connecteurs : \all{weil, aber, dann}.
\end{corrige}

\courssub{Point culture et pièges à éviter}

\begin{savaistu}[La \all{Ruhezeit} en Allemagne et au Cameroun]
En Allemagne, la protection du calme (\all{Lärmschutz}) est strictement réglementée par la \all{Bundesimmissionsschutzgesetz} (Loi fédérale sur la protection contre les immissions) et les règlements locaux (\all{Landesimmissionsschutzgesetze}).
\begin{itemize}
\item \textbf{\all{Nachtruhe} (Calme nocturne) :} Généralement de \textbf{22h00 à 6h00} (parfois 7h00). Tout bruit évitable est interdit (\all{verbotene Geräusche} : tondeuse, perceuse, musique forte, fêtes).
\item \textbf{\all{Mittagsruhe} (Calme de midi) :} Dans de nombreux \all{Bundesländer} (Länder), de \textbf{13h00 à 15h00}.
\item \textbf{\all{Sonntags- / Feiertagsruhe} :} Toute la journée le dimanche et jours fériés. Interdiction de travaux bruyants, tonte de pelouse, lessive bruyante (dans certains vieux immeubles), recyclage du verre.
\item \textbf{Sanctions :} Amendes (\all{Bußgeld}) pouvant aller jusqu'à 5000 €. L'\all{Ordnungsamt} (office de l'ordre public) intervient sur plainte (\all{Anzeige}).
\end{itemize}
\textbf{Au Cameroun :} La loi n°96/12 du 5 août 1996 relative à la gestion de l'environnement (art. 17, 18) et les arrêtés municipaux prévoient la lutte contre les nuisances sonores. Cependant, l'application est souvent laxiste. Les bars, églises, mosquées, et \all{bendskins} (moto-taxis) sont les principales sources. La sensibilisation (\all{Sensibilisierung}) et le dialogue de bon voisinage restent les premiers recours.
\end{savaistu}

\begin{attention}[Les 5 erreurs qui coûtent des points au Bac]
\begin{enumerate}
\item \textbf{Confondre \all{wenn} et \all{als}} : Pour un événement \textbf{unique et daté au passé}, seul \all{als} est correct. \all{Als ich gestern nach Hause kam, ...} (pas \all{*wenn}). \all{Wenn} = répétition (passé) ou présent/futur.
\item \textbf{Oublier le verbe à la fin} dans les subordonnées (\all{weil, dass, als, wenn}) et les relatives. \all{Die Musik, die aus der Bar \textbf{kommt}, ist laut.} — Jamais \all{*die kommt aus der Bar}. \all{Ich schlafe nicht, weil es \textbf{laut ist}.} (pas \all{*weil es ist laut}).
\item \textbf{Mettre un Umlaut aux modaux au Präteritum} : On écrit \all{konnte, musste, durfte, wollte, sollte, mochte} — \textbf{jamais} \all{könnte, müsste, dürfte} (ce sont des Konjunktiv II : « pourrais, devrais »).
\item \textbf{Oublier la majuscule des noms} : \all{der Lärm, die Ruhe, die Nachbarn, das Geräusch, die Bar, der Besitzer}. Une minuscule sur un nom commun est une faute d'orthographe sanctionnée.
\item \textbf{Calquer le français (Faux amis / Collocations)} :
      \begin{itemize}
      \item « Faire du bruit » = \all{Lärm machen} (pas \all{*Lärm tun} ou \all{*Bruit machen}).
      \item « Se plaindre de » = \all{sich beschweren über + Akkusativ} (pas \all{*sich beklagen über} — \all{beklagen} prend souvent un objet direct ou \all{über} mais plus formel, \all{beschweren} est standard pour « se plaindre auprès de qqn de qqch »).
      \item \all{stören} = déranger (pas « stocker » $\rightarrow$ \all{lagern}).
      \item \all{die Mappe} = la pochette / le classeur (pas « la carte » $\rightarrow$ \all{die Karte}).
      \item \all{wegen + GÉNITIF} : \all{wegen des Lärms} (pas \all{*wegen den Lärm} — bien que l'allemand courant accepte de plus en plus l'accusatif, l'examen exige le génitif).
      \end{itemize}
\end{enumerate}
\end{attention}

\newpage

\coursec{Exercices}

\courssub{Je vérifie mes connaissances}

\begin{exercice}[QCM : le vocabulaire du bruit]
Entoure la bonne réponse.
\begin{enumerate}
\item \all{Der Lärm} signifie :
\begin{itemize}
\item[a)] la musique \item[b)] le bruit \item[c)] le silence
\end{itemize}
\item Le contraire de \all{laut} est :
\begin{itemize}
\item[a)] \all{leise} \item[b)] \all{schnell} \item[c)] \all{schwer}
\end{itemize}
\item \all{Die Nachbarn stören mich} signifie :
\begin{itemize}
\item[a)] Les voisins m'aident. \item[b)] Les voisins me dérangent. \item[c)] Les voisins m'écoutent.
\end{itemize}
\item Quel est le pluriel de \all{die Ruhezeit} ?
\begin{itemize}
\item[a)] \all{die Ruhezeiten} \item[b)] \all{die Ruhezeiter} \item[c)] \all{die Ruhezeit}
\end{itemize}
\item \all{Wenn ich schlafen will, brauche ich ...}
\begin{itemize}
\item[a)] \all{die Ruhe} \item[b)] \all{das Radio} \item[c)] \all{die Baustelle}
\end{itemize}
\end{enumerate}
\end{exercice}

\begin{exercice}[Vrai ou faux ? Justifie ta réponse]
Réponds par \all{richtig} ou \all{falsch} et justifie en une phrase en français.
\begin{enumerate}
\item \all{Als} s'emploie pour une action répétée au passé.
\item Dans une subordonnée introduite par \all{wenn}, le verbe conjugué se place à la fin.
\item Dans la phrase \all{Der Mann, der nebenan wohnt, ist laut}, \all{der} est un pronom relatif.
\item Le prétérit de \all{ich kann} est \all{ich konnte}.
\item \all{Wenn} s'emploie pour une action unique au passé.
\end{enumerate}
\end{exercice}

\begin{exercice}[Vocabulaire : trouve le mot]
Complète chaque phrase avec le mot qui convient : \all{der Lärm, die Ruhe, die Nachbarn, stören, leise, die Baustelle}.
\begin{enumerate}
\item Vor unserem Haus gibt es eine große \all{...........} : die Arbeiter bauen ein neues Haus.
\item Meine \all{...........} hören jeden Abend laute Musik.
\item Bitte sprich \all{...........} ! Das Baby schläft.
\item Der \all{...........} von den Motorrädern ist sehr stark.
\item Nachts brauche ich \all{...........}, um gut zu schlafen.
\item Die Autos \all{...........} die Schüler beim Lernen.
\end{enumerate}
\end{exercice}

\begin{exercice}[Conjugaison : le prétérit]
Mets les verbes entre parenthèses au prétérit.
\begin{enumerate}
\item Gestern \all{(sein)} die Straße sehr laut.
\item Wir \all{(haben)} keine Ruhe vor den Motorrädern.
\item Ich \all{(können)} nicht schlafen, weil die Musik so laut war.
\item Die Schüler \all{(müssen)} die Fenster schließen.
\item Mein Bruder \all{(wollen)} die Polizei rufen.
\item Ihr \all{(sein)} gestern Abend sehr müde.
\end{enumerate}
\end{exercice}

\courssub{Je m'entraîne}

\begin{exercice}[Grammaire : wenn ou als ?]
Complète les phrases avec \all{wenn} ou \all{als} et conjugue le verbe entre parenthèses à la bonne place.
\begin{enumerate}
\item \all{...........} ich ein Kind \all{(sein)}, wohnten wir neben einer Baustelle.
\item \all{...........} die Nachbarn Musik \all{(hören)}, kann ich nicht lernen.
\item \all{...........} ich gestern nach Hause \all{(kommen)}, war die Straße sehr laut.
\item \all{...........} es nachts laut \all{(sein)}, rufen wir die Polizei.
\item \all{...........} wir in Douala \all{(leben)}, gab es viel Lärm von den Motorrädern.
\item \all{...........} mein Vater früher in Yaoundé \all{(arbeiten)}, war das Büro sehr ruhig.
\item \all{...........} die Baustelle um sechs Uhr \all{(anfangen)}, wache ich immer auf.
\item \all{...........} das Fest letztes Jahr zu Ende \all{(sein)}, waren alle froh.
\end{enumerate}
\end{exercice}

\begin{exercice}[Grammaire : les propositions relatives]
Relie les deux phrases avec un pronom relatif (\all{der, die, das, den, dem}).
\begin{enumerate}
\item Der Lärm kommt von der Baustelle. Der Lärm stört mich.
\item Die Nachbarn hören laute Musik. Die Nachbarn wohnen nebenan.
\item Das Radio ist zu laut. Mein Bruder hat das Radio gekauft.
\item Die Straße ist sehr laut. Wir wohnen an der Straße.
\item Der Motorradfahrer macht viel Lärm. Ich sehe den Motorradfahrer jeden Morgen.
\item Das Zimmer ist ruhig. Ich lerne in dem Zimmer.
\end{enumerate}
\end{exercice}

\begin{exercice}[Texte à trous]
Complète le texte avec les mots suivants : \all{störte, als, konnte, Ruhe, der, musste, laut, wenn}.

\begin{textebox}[Lärm in der Nacht]
Gestern Nacht war es in unserer Straße sehr \all{...........}. \all{...........} ich ins Bett ging, begann eine Party bei den Nachbarn. Die Musik, \all{...........} aus dem Garten kam, \all{...........} die ganze Familie. Ich \all{...........} nicht schlafen und meine kleine Schwester weinte. Mein Vater \all{...........} zu den Nachbarn gehen und um \all{...........} bitten. \all{...........} er zurückkam, war die Musik endlich leiser.
\end{textebox}
\end{exercice}

\begin{exercice}[Appariement : mots et définitions]
Relie chaque mot allemand à sa définition française.
\begin{center}
\begin{tabularx}{\linewidth}{|X|X|}
\hline
\rowcolor{popblueL} \textbf{Deutsch} & \textbf{Français} \\
\hline
1. \all{die Baustelle} & a) le calme, le silence \\
\hline
2. \all{die Ruhe} & b) un chantier de construction \\
\hline
3. \all{der Nachbar} & c) gêner, déranger \\
\hline
4. \all{stören} & d) la personne qui habite à côté \\
\hline
5. \all{die Ruhezeit} & e) la période où l'on doit être silencieux \\
\hline
6. \all{leise} & f) qui fait peu de bruit \\
\hline
\end{tabularx}
\end{center}
\end{exercice}

\courssub{Je me prépare au Bac}

\begin{exercice}[Compréhension écrite type Bac]
Lis le texte suivant, puis réponds aux questions.

\begin{textebox}[Lärm in München — ein Problem für alle]
München ist eine schöne Stadt, aber viele Einwohner klagen über den Lärm. Der Verkehr ist das größte Problem: Autos, Busse und Motorräder fahren Tag und Nacht durch die Straßen. Familie Huber wohnt seit zehn Jahren im Zentrum. „Als wir hier einzogen, war die Straße noch ruhig“, erzählt Frau Huber. „Aber jetzt können wir nachts oft nicht schlafen.“ Auch die Baustellen stören die Menschen. Vor dem Haus der Familie bauen Arbeiter seit sechs Monaten eine neue U-Bahn. Die Arbeit beginnt schon um sechs Uhr morgens. Wenn die Maschinen arbeiten, muss die Familie die Fenster schließen. Die Stadt München will das Problem lösen: Nachts gilt eine Ruhezeit von 22 Uhr bis 6 Uhr. Wer laute Musik spielt, muss eine Strafe zahlen. Frau Huber sagt: „Wenn alle Menschen die Regeln respektieren, wird unsere Stadt wieder ruhiger.“ (environ 150 mots)
\end{textebox}

\textbf{A. Questions de compréhension} — Réponds en allemand par des phrases complètes.
\begin{enumerate}
\item Was ist das größte Lärmproblem in München?
\item Seit wann wohnt Familie Huber im Zentrum?
\item Was bauen die Arbeiter vor dem Haus der Familie?
\item Wann beginnt die Ruhezeit in München?
\item Was muss man zahlen, wenn man nachts laute Musik spielt?
\end{enumerate}

\textbf{B. Vrai ou faux ?} — Réponds par \all{richtig} ou \all{falsch} et justifie avec une phrase du texte (en allemand).
\begin{enumerate}
\item Die Straße von Familie Huber war früher laut.
\item Die Baustelle gibt es schon seit einem Jahr.
\item Frau Huber glaubt, dass die Stadt ruhiger werden kann.
\end{enumerate}

\textbf{C. Questions de langue}
\begin{enumerate}
\item Trouve dans le texte une subordonnée avec \all{als} et une avec \all{wenn}, et explique la différence d'emploi.
\item Relève le prétérit de \all{können} et de \all{müssen} dans le texte.
\end{enumerate}
\end{exercice}

\begin{exercice}[Thème et version]
\textbf{A. Version} — Traduis en français le paragraphe suivant.

\begin{textebox}[Die Ruhezeiten in Deutschland]
In Deutschland gibt es klare Regeln für die Ruhe. Die Nachtruhe beginnt um 22 Uhr und endet um 6 Uhr. Wenn die Nachbarn in dieser Zeit laute Musik hören, kann man die Polizei rufen. Auch am Sonntag muss man leise sein: Man darf nicht mit lauten Maschinen im Garten arbeiten. Viele Deutsche respektieren diese Regeln, denn die Ruhe ist ihnen wichtig.
\end{textebox}

\textbf{B. Thème} — Traduis en allemand les phrases suivantes.
\begin{enumerate}
\item Quand j'étais petit, nous habitions près d'un chantier. (utiliser \all{als})
\item Quand les voisins écoutent de la musique forte, je ne peux pas dormir. (utiliser \all{wenn})
\item Le bruit qui vient de la rue dérange les élèves. (proposition relative)
\item Hier soir, je devais fermer la fenêtre parce que la fête était trop bruyante. (modal au prétérit)
\item La nuit, tout le monde a besoin de calme.
\end{enumerate}
\end{exercice}

\begin{exercice}[Production écrite type Bac]
\textbf{Sujet :} \all{Beschreibe eine Lärmsituation, die du erlebt hast, und schlage eine Lösung vor.}

Décris une situation de nuisance sonore que tu as vécue (chez toi, à l'école, dans ton quartier) et propose une solution. Rédige un texte d'environ \textbf{10 lignes} (80 à 100 mots) en allemand.

\textbf{Consignes :}
\begin{itemize}
\item Situe la situation (où ? quand ? qui ?).
\item Utilise au moins une subordonnée avec \all{als} (passé) et une avec \all{wenn}.
\item Utilise au moins une proposition relative (\all{der, die, das}...).
\item Emploie au moins deux verbes au prétérit parmi : \all{sein, haben, können, müssen, wollen}.
\item Termine par une proposition de solution.
\end{itemize}

\textbf{Grille critériée :}
\begin{center}
\begin{tabularx}{\linewidth}{|X|c|}
\hline
\rowcolor{popblueL} \textbf{Critère} & \textbf{Points} \\
\hline
Respect du sujet et de la longueur & 4 \\
\hline
Vocabulaire du thème (Lärm, Ruhe, stören...) & 4 \\
\hline
Subordonnées \all{wenn}/\all{als} correctes & 4 \\
\hline
Proposition relative correcte & 4 \\
\hline
Prétérit et correction grammaticale générale & 4 \\
\hline
\end{tabularx}
\end{center}
\end{exercice}

\newpage

\coursec{Corrigés des exercices}

\begin{corrige}[Exercice 1 — QCM : le vocabulaire du bruit]
\begin{enumerate}
\item \textbf{b)} \all{Der Lärm} signifie « le bruit ». (\all{die Musik} = la musique ; \all{die Stille} = le silence.)
\item \textbf{a)} Le contraire de \all{laut} (bruyant, fort) est \all{leise} (silencieux, doux). \all{schnell} = rapide ; \all{schwer} = lourd, difficile.
\item \textbf{b)} \all{Die Nachbarn stören mich} = « Les voisins me dérangent. » \all{stören} = déranger, gêner. Attention au faux ami : \all{helfen} = aider, \all{zuhören} = écouter.
\item \textbf{a)} Le pluriel de \all{die Ruhezeit} est \all{die Ruhezeiten}. Les noms féminins en \all{-zeit} forment leur pluriel en \all{-en}.
\item \textbf{a)} \all{Wenn ich schlafen will, brauche ich die Ruhe.} = « Quand je veux dormir, j'ai besoin de calme. »
\end{enumerate}
\end{corrige}

\begin{corrige}[Exercice 2 — Vrai ou faux ?]
\begin{enumerate}
\item \textbf{falsch} : \all{als} s'emploie pour une action \emph{unique} au passé (un seul événement, souvent au prétérit). Pour une action répétée au passé, on utilise \all{wenn}.
\item \textbf{richtig} : \all{wenn} est une conjonction de subordination ; le verbe conjugué se rejette à la fin de la subordonnée : \all{Wenn ich Zeit habe, ...}.
\item \textbf{richtig} : dans \all{Der Mann, der nebenan wohnt, ist laut}, \all{der} reprend \all{der Mann} (masculin, nominatif, sujet de \all{wohnt}) : c'est bien un pronom relatif.
\item \textbf{richtig} : le prétérit de \all{können} est \all{konnte} : \all{ich konnte, du konntest, er konnte...} (sans Umlaut au prétérit).
\item \textbf{falsch} : pour une action unique au passé, on utilise \all{als}. \all{wenn} s'emploie au présent/futur ou pour les habitudes (même au passé : \all{Wenn ich früher kam, ...} = « chaque fois que je venais »).
\end{enumerate}
\end{corrige}

\begin{corrige}[Exercice 3 — Vocabulaire : trouve le mot]
\begin{enumerate}
\item Vor unserem Haus gibt es eine große \textbf{Baustelle} : die Arbeiter bauen ein neues Haus. « Devant notre maison, il y a un grand chantier. »
\item Meine \textbf{Nachbarn} hören jeden Abend laute Musik. « Mes voisins écoutent de la musique forte chaque soir. »
\item Bitte sprich \textbf{leise} ! Das Baby schläft. « Parle doucement, s'il te plaît ! Le bébé dort. »
\item Der \textbf{Lärm} von den Motorrädern ist sehr stark. « Le bruit des motos est très fort. »
\item Nachts brauche ich \textbf{Ruhe}, um gut zu schlafen. « La nuit, j'ai besoin de calme pour bien dormir. »
\item Die Autos \textbf{stören} die Schüler beim Lernen. « Les voitures dérangent les élèves pendant qu'ils travaillent. »
\end{enumerate}
\end{corrige}

\begin{corrige}[Exercice 4 — Conjugaison : le prétérit]
\begin{enumerate}
\item Gestern \textbf{war} die Straße sehr laut. (\all{sein} : \all{ich/er war})
\item Wir \textbf{hatten} keine Ruhe vor den Motorrädern. (\all{haben} : \all{wir hatten})
\item Ich \textbf{konnte} nicht schlafen, weil die Musik so laut war. (\all{können} : \all{ich konnte}, sans Umlaut au prétérit)
\item Die Schüler \textbf{mussten} die Fenster schließen. (\all{müssen} : \all{sie mussten})
\item Mein Bruder \textbf{wollte} die Polizei rufen. (\all{wollen} : \all{er wollte})
\item Ihr \textbf{wart} gestern Abend sehr müde. (\all{sein} : \all{ihr wart})
\end{enumerate}
Rappel des paradigmes : \all{sein} : \all{war, warst, war, waren, wart, waren} ; \all{haben} : \all{hatte, hattest, hatte, hatten, hattet, hatten} ; modaux : \all{konnte, musste, wollte, durfte, sollte} (toujours sans Umlaut au prétérit).
\end{corrige}

\begin{corrige}[Exercice 5 — Grammaire : wenn ou als ?]
\begin{enumerate}
\item \textbf{Als} ich ein Kind \textbf{war}, wohnten wir neben einer Baustelle. (période unique du passé $\rightarrow$ \all{als})
\item \textbf{Wenn} die Nachbarn Musik \textbf{hören}, kann ich nicht lernen. (présent, habitude $\rightarrow$ \all{wenn})
\item \textbf{Als} ich gestern nach Hause \textbf{kam}, war die Straße sehr laut. (action unique au passé $\rightarrow$ \all{als}, verbe fort \all{kommen} : \all{kam})
\item \textbf{Wenn} es nachts laut \textbf{ist}, rufen wir die Polizei. (présent, cas général $\rightarrow$ \all{wenn})
\item \textbf{Als} wir in Douala \textbf{lebten}, gab es viel Lärm von den Motorrädern. (période passée unique $\rightarrow$ \all{als} ; \all{leben} est faible : \all{lebten})
\item \textbf{Als} mein Vater früher in Yaoundé \textbf{arbeitete}, war das Büro sehr ruhig. (période passée $\rightarrow$ \all{als} ; \all{arbeiten} : \all{arbeitete}, avec \all{-e-} de liaison)
\item \textbf{Wenn} die Baustelle um sechs Uhr \textbf{anfängt}, wache ich immer auf. (habitude au présent $\rightarrow$ \all{wenn} ; verbe à particule : \all{anfängt} en fin de subordonnée, particule recollée)
\item \textbf{Als} das Fest letztes Jahr zu Ende \textbf{war}, waren alle froh. (événement unique au passé $\rightarrow$ \all{als})
\end{enumerate}
\end{corrige}

\begin{corrige}[Exercice 6 — Grammaire : les propositions relatives]
\begin{enumerate}
\item Der Lärm, \textbf{der} mich stört, kommt von der Baustelle. (\all{der Lärm} : masculin, nominatif, sujet de \all{stört}.)
\item Die Nachbarn, \textbf{die} nebenan wohnen, hören laute Musik. (\all{die Nachbarn} : pluriel, nominatif $\rightarrow$ \all{die}.)
\item Das Radio, \textbf{das} mein Bruder gekauft hat, ist zu laut. (\all{das Radio} : neutre, accusatif car complément de \all{gekauft hat} ; pour le neutre, accusatif = nominatif : \all{das}.)
\item Die Straße, \textbf{an der} wir wohnen, ist sehr laut. (féminin avec préposition \all{an} + datif : \all{an der}.)
\item Der Motorradfahrer, \textbf{den} ich jeden Morgen sehe, macht viel Lärm. (masculin, accusatif car complément de \all{sehe} : \all{den}.)
\item Das Zimmer, \textbf{in dem} ich lerne, ist ruhig. (neutre avec préposition \all{in} + datif : \all{in dem}.)
\end{enumerate}
Point de méthode : le pronom relatif prend le \emph{genre et le nombre} du nom qu'il reprend, et le \emph{cas} exigé par sa fonction dans la relative ; le verbe conjugué se place à la fin.
\end{corrige}

\begin{corrige}[Exercice 7 — Texte à trous]
\begin{textebox}[Lärm in der Nacht — corrigé]
Gestern Nacht war es in unserer Straße sehr \textbf{laut}. \textbf{Als} ich ins Bett ging, begann eine Party bei den Nachbarn. Die Musik, \textbf{die} aus dem Garten kam, \textbf{störte} die ganze Familie. Ich \textbf{konnte} nicht schlafen und meine kleine Schwester weinte. Mein Vater \textbf{musste} zu den Nachbarn gehen und um \textbf{Ruhe} bitten. \textbf{Als} er zurückkam, war die Musik endlich leiser.
\end{textebox}
Justifications : \all{Als} (deux fois) car il s'agit d'événements uniques au passé ; \all{die} reprend \all{die Musik} (féminin, nominatif) ; \all{störte} est le prétérit faible de \all{stören} ; \all{konnte} et \all{musste} sont les prétérits des modaux \all{können} et \all{müssen} ; \all{um Ruhe bitten} = « demander le calme ».
\end{corrige}

\begin{corrige}[Exercice 8 — Appariement : mots et définitions]
\begin{center}
\begin{tabularx}{\linewidth}{|X|X|}
\hline
\rowcolor{popblueL} \textbf{Deutsch} & \textbf{Français} \\
\hline
1. \all{die Baustelle} & \textbf{b)} un chantier de construction \\
\hline
2. \all{die Ruhe} & \textbf{a)} le calme, le silence \\
\hline
3. \all{der Nachbar} & \textbf{d)} la personne qui habite à côté \\
\hline
4. \all{stören} & \textbf{c)} gêner, déranger \\
\hline
5. \all{die Ruhezeit} & \textbf{e)} la période où l'on doit être silencieux \\
\hline
6. \all{leise} & \textbf{f)} qui fait peu de bruit \\
\hline
\end{tabularx}
\end{center}
Réponses : 1-b, 2-a, 3-d, 4-c, 5-e, 6-f.
\end{corrige}

\begin{corrige}[Exercice 9 — Compréhension écrite type Bac]
\textbf{A. Questions de compréhension} (réponses attendues, formulations voisines acceptées)
\begin{enumerate}
\item \all{Das größte Lärmproblem in München ist der Verkehr: Autos, Busse und Motorräder.} « Le plus grand problème de bruit à Munich, c'est la circulation. »
\item \all{Familie Huber wohnt seit zehn Jahren im Zentrum.} « La famille Huber habite au centre depuis dix ans. »
\item \all{Die Arbeiter bauen eine neue U-Bahn vor dem Haus der Familie.} « Les ouvriers construisent une nouvelle ligne de métro. »
\item \all{Die Ruhezeit beginnt um 22 Uhr (und endet um 6 Uhr).} « La période de silence commence à 22 heures. »
\item \all{Man muss eine Strafe zahlen.} « On doit payer une amende. »
\end{enumerate}
\textbf{B. Vrai ou faux ?}
\begin{enumerate}
\item \textbf{falsch} : \all{„Als wir hier einzogen, war die Straße noch ruhig.“} — la rue était calme autrefois.
\item \textbf{falsch} : \all{„...bauen Arbeiter seit sechs Monaten eine neue U-Bahn.“} — le chantier dure depuis six mois, pas un an.
\item \textbf{richtig} : \all{„Wenn alle Menschen die Regeln respektieren, wird unsere Stadt wieder ruhiger.“}
\end{enumerate}
\textbf{C. Questions de langue}
\begin{enumerate}
\item \all{als} : \all{„Als wir hier einzogen, war die Straße noch ruhig“} — événement unique au passé (prétérit). \all{wenn} : \all{„Wenn die Maschinen arbeiten, muss die Familie die Fenster schließen“} — situation répétée/générale au présent. Règle : \all{als} = une seule fois au passé ; \all{wenn} = présent, futur ou répétition.
\item \all{können} : \all{„können wir nachts oft nicht schlafen“} est au \emph{présent} ; le prétérit dans le texte : il n'y a pas de forme de \all{können} au prétérit ; en revanche \all{müssen} apparaît au présent (\all{muss}). Formes de prétérit à connaître : \all{können} $\rightarrow$ \all{konnte}, \all{müssen} $\rightarrow$ \all{musste}. (Accepter aussi la citation de \all{„Als wir hier einzogen“} pour \all{einziehen} : \all{zogen ein}.)
\end{enumerate}
\textbf{Barème indicatif (sur 20)} : A. 5 $\times$ 2 = 10 pts (1 pt sens, 1 pt correction de la phrase) ; B. 3 $\times$ 2 = 6 pts (1 pt \all{richtig}/\all{falsch}, 1 pt justification citée) ; C. 2 $\times$ 2 = 4 pts.
\textbf{Points de vigilance} : répondre par des phrases complètes avec le verbe en deuxième position ; citer la justification en allemand, sans la traduire ; ne pas confondre \all{seit zehn Jahren} (depuis dix ans) et \all{vor zehn Jahren} (il y a dix ans).
\end{corrige}

\begin{corrige}[Exercice 10 — Thème et version]
\textbf{A. Version} — Traduction proposée :
« En Allemagne, il existe des règles claires pour le calme. Le silence nocturne commence à 22 heures et se termine à 6 heures. Quand les voisins écoutent de la musique forte pendant cette période, on peut appeler la police. Le dimanche aussi, on doit être silencieux : on n'a pas le droit de travailler dans le jardin avec des machines bruyantes. Beaucoup d'Allemands respectent ces règles, car le calme est important pour eux. »
Points de vigilance : \all{die Nachtruhe} = le silence nocturne ; \all{man darf nicht} = « on n'a pas le droit de » (et non « on ne doit pas obligatoirement ») ; \all{denn} = car.

\textbf{B. Thème} — Traductions proposées :
\begin{enumerate}
\item \all{Als ich klein war, wohnten wir in der Nähe einer Baustelle.} (\all{als} + prétérit ; verbe en fin de subordonnée.)
\item \all{Wenn die Nachbarn laute Musik hören, kann ich nicht schlafen.} (\all{wenn} pour une habitude ; \all{schlafen} = dormir.)
\item \all{Der Lärm, der von der Straße kommt, stört die Schüler.} (relatif \all{der} : masculin nominatif.)
\item \all{Gestern Abend musste ich das Fenster schließen, weil die Party zu laut war.} (\all{musste} : prétérit de \all{müssen} ; \all{schließen} avec ß.)
\item \all{Nachts braucht jeder (Mensch) Ruhe.} ou \all{In der Nacht brauchen alle Menschen Ruhe.}
\end{enumerate}
Barème indicatif : version 10 pts (fidélité 6, qualité du français 4) ; thème 5 $\times$ 2 = 10 pts (1 pt grammaire visée, 1 pt exactitude globale).
\end{corrige}

\begin{corrige}[Exercice 11 — Production écrite type Bac]
\textbf{Proposition de rédaction modèle} (environ 95 mots) :
\begin{textebox}[Eine Lärmsituation in meinem Viertel]
Letztes Jahr wohnte ich mit meiner Familie in der Nähe einer großen Straße in Douala. Als ich in die neue Wohnung einzog, war ich froh, aber der Lärm war ein großes Problem. Die Motorräder, die Tag und Nacht vorbeifuhren, störten uns sehr. Wenn ich meine Hausaufgaben machen wollte, konnte ich mich nicht konzentrieren. Nachts hatte ich keine Ruhe, und ich musste oft die Fenster schließen. Mein Vater wollte sogar umziehen. Ich schlage eine Lösung vor: Die Stadt soll eine Ruhezeit nach 22 Uhr einführen, und alle Motorradfahrer sollen leiser fahren. Dann wird unser Viertel wieder ruhig.
\end{textebox}
Traduction : « L'année dernière, j'habitais avec ma famille près d'une grande rue à Douala. Quand j'ai emménagé dans le nouvel appartement, j'étais content, mais le bruit était un grand problème. Les motos, qui passaient jour et nuit, nous dérangeaient beaucoup. Quand je voulais faire mes devoirs, je ne pouvais pas me concentrer. La nuit, je n'avais pas de calme et je devais souvent fermer les fenêtres. Mon père voulait même déménager. Je propose une solution : la ville devrait instaurer une période de silence après 22 heures, et tous les motards devraient rouler moins bruyamment. Alors notre quartier redeviendra calme. »
\textbf{Barème indicatif (sur 20)} : respect du sujet et de la longueur 4 pts ; vocabulaire du thème 4 pts ; subordonnées \all{wenn}/\all{als} 4 pts ; proposition relative 4 pts ; prétérit et correction générale 4 pts.
\textbf{Points de vigilance} : verbe conjugué en fin de subordonnée (\all{Als ich ... einzog}) ; verbe en fin de relative (\all{die ... vorbeifuhren}) ; prétérits corrects (\all{war, hatte, konnte, musste, wollte}) ; majuscules à tous les noms ; ne pas oublier la solution finale.
\end{corrige}

\newpage

\coursec{Fiche bilan}

\begin{popfigure}
\begin{tikzpicture}[mindmap, grow cyclic, every node/.style={concept}, concept color=popblue!70,
root concept/.append style={concept color=popblue, font=\bfseries, text=white},
level 1/.append style={level distance=3.2cm, sibling angle=51.4, font=\small},
level 2/.append style={level distance=2.2cm, font=\footnotesize}]
\node[root concept, align=center]{Lärm und Ruhe\\ (ALA13)}
  child[concept color=poporange!80] { node {Lexique} 
    child { node {der Lärm, laut, leise} }
    child { node {stören, die Ruhe} }
  }
  child[concept color=popgreen!80] { node[align=center]{Subordonnées\\ temporelles}
    child { node[align=center]{wenn : présent,\\ futur, répétition} }
    child { node[align=center]{als : passé,\\ une seule fois} }
  }
  child[concept color=popgold!90] { node[align=center]{Propositions\\ relatives}
    child { node[align=center]{der, die, das,\\ die (pl.)} }
    child { node[align=center]{verbe\\ à la fin} }
  }
  child[concept color=poppurple!80] { node {Präteritum}
    child { node {war, hatte} }
    child { node[align=center]{musste, konnte,\\ wollte, durfte} }
  }
  child[concept color=poppink!80] { node[align=center]{Compréhension\\ du texte}
    child { node[align=center]{global puis\\ détaillé} }
  }
  child[concept color=popgreen!60] { node {Production}
    child { node[align=center]{résumé,\\ paragraphe guidé} }
  }
  child[concept color=poporange!60] { node[align=center]{Méthode\\ commentaire}
    child { node[align=center]{lire, repérer,\\ reformuler} }
  };
\end{tikzpicture}
\legende{Carte mentale de la leçon ALA13 : les bruits et les nuisances sonores}
\end{popfigure}

\begin{bilan}

\textbf{1. Lexique clé à connaître par cœur}

\begin{center}
\begin{tabularx}{\linewidth}{|l|X|}
\hline
\rowcolor{popblueL} \textbf{Deutsch} & \textbf{Français} \\
\hline
der Lärm (Sg.) & le bruit, le vacarme \\
\hline
laut / leise & bruyant, fort / silencieux, doux \\
\hline
der Nachbar, -n / die Nachbarin, -nen & le voisin / la voisine \\
\hline
stören & déranger, troubler \\
\hline
die Ruhe (Sg.) & le calme, le silence \\
\hline
die Lärmbelästigung, -en & la nuisance sonore \\
\hline
sich beschweren (über + Akk.) & se plaindre (de) \\
\hline
der Verkehr (Sg.) & la circulation \\
\hline
die Baustelle, -n & le chantier \\
\hline
ohne / gegen den Lärm & sans / contre le bruit \\
\hline
\end{tabularx}
\end{center}

\textbf{2. Grammaire : les règles essentielles}

\begin{center}
\begin{tabularx}{\linewidth}{|l|X|X|}
\hline
\rowcolor{popblueL} \textbf{Point} & \textbf{Règle} & \textbf{Exemple} \\
\hline
\all{wenn} & présent, futur ou répétition au passé ; verbe en fin de subordonnée & \all{Wenn der Lärm zu laut ist, kann ich nicht schlafen.} \\
\hline
\all{als} & événement unique au passé (Präteritum/Perfekt) & \all{Als ich nach Hause kam, war es endlich ruhig.} \\
\hline
Relatives & pronom relatif = genre et nombre du nom, cas selon la fonction ; verbe en fin & \all{Der Nachbar, der nebenan wohnt, ist sehr laut.} \\
\hline
Präteritum & \all{sein : war ; haben : hatte ; müssen : musste ; können : konnte ; wollen : wollte ; dürfen : durfte ; sollen : sollte} & \all{Ich konnte wegen des Lärms nicht lernen.} \\
\hline
\end{tabularx}
\end{center}

\textbf{3. Méthodes express}
\begin{itemize}
\item \textbf{Comprendre le texte} : 1) lire le titre et anticiper ; 2) lecture globale (qui ? où ? quel problème ?) ; 3) lecture détaillée avec le lexique ; 4) répondre par des phrases complètes.
\item \textbf{Choisir entre \all{wenn} et \all{als}} : phrase au passé + action unique $\rightarrow$ \all{als} ; sinon $\rightarrow$ \all{wenn}.
\item \textbf{Construire une relative} : repérer le nom de référence $\rightarrow$ choisir \all{der/die/das} selon le genre $\rightarrow$ adapter le cas $\rightarrow$ placer le verbe conjugué à la fin.
\item \textbf{Produire un paragraphe} : 4 à 6 phrases, une idée par phrase, connecteurs \all{weil, wenn, aber, deshalb}.
\end{itemize}

\textbf{4. Pièges de francophones}
\begin{itemize}
\item Dans toute subordonnée (\all{wenn, als, weil, dass}), le verbe conjugué va \textbf{à la fin}.
\item \all{wenn} ne signifie pas « si » hypothétique ici ; attention à ne pas traduire « quand » au passé par \all{wenn} pour un événement unique.
\item Le prétérit des modaux ne prend \textbf{jamais} de \all{ge-} : \all{musste}, pas « gemusst ».
\item Majuscule aux noms : \all{der Lärm, die Ruhe, der Nachbar}.
\end{itemize}

\end{bilan}

\begin{aretenir}[Checklist avant le Bac]
\begin{itemize}
\item[$\square$] Je connais le lexique du bruit : \all{der Lärm, laut, leise, stören, die Ruhe, die Nachbarn}.
\item[$\square$] Je sais donner l'article et le pluriel de chaque nom du thème.
\item[$\square$] Je distingue \all{wenn} (répétition, présent, futur) et \all{als} (passé, une seule fois).
\item[$\square$] Je place correctement le verbe conjugué à la fin des subordonnées temporelles.
\item[$\square$] Je choisis le bon pronom relatif (\all{der, die, das, die} au pluriel) selon le genre et le cas.
\item[$\square$] Je connais par cœur le prétérit de \all{sein} (war) et \all{haben} (hatte).
\item[$\square$] Je connais le prétérit des modaux : \all{musste, konnte, wollte, durfte, sollte, mochte}.
\item[$\square$] Je sais résumer le texte \all{Lärm in der Stadt} en trois ou quatre phrases.
\item[$\square$] Je sais répondre aux questions de compréhension par des phrases complètes.
\item[$\square$] Je sais rédiger un court paragraphe sur les nuisances sonores dans mon quartier.
\item[$\square$] Je n'oublie pas la majuscule aux noms allemands et les guillemets „…“ dans une citation.
\item[$\square$] Je relis ma production : place du verbe, accords, orthographe.
\end{itemize}
\end{aretenir}

\findecours

\end{document}
